\documentclass[11pt]{article}
\usepackage[T1]{fontenc}
\usepackage[utf8]{inputenc}
\usepackage{lmodern,amsmath,amssymb,amsthm,mathtools,mathrsfs}
\usepackage[margin=26mm]{geometry}
\usepackage{booktabs,longtable,array}
\usepackage[hidelinks]{hyperref}
\usepackage{microtype}
\setlength{\emergencystretch}{2em}
\raggedbottom
\numberwithin{equation}{section}
\newtheorem{theorem}{Theorem}[section]
\newtheorem{lemma}[theorem]{Lemma}
\newtheorem{proposition}[theorem]{Proposition}
\newtheorem{corollary}[theorem]{Corollary}
\theoremstyle{definition}
\newtheorem{definition}[theorem]{Definition}
\newtheorem{remark}[theorem]{Remark}
\newcommand{\C}{\mathbb C}
\newcommand{\R}{\mathbb R}
\newcommand{\Q}{\mathbb Q}
\newcommand{\Z}{\mathbb Z}
\newcommand{\T}{\mathbb T}
\newcommand{\ii}{\mathrm i}
\newcommand{\dd}{\,\mathrm d}
\newcommand{\eps}{\varepsilon}
\renewcommand{\DH}{\mathcal D_H}
\newcommand{\DL}{\mathcal D_L}
\newcommand{\XH}{\mathcal X_H}
\newcommand{\XL}{\mathcal X_L}
\newcommand{\cO}{\mathcal O}
\newcommand{\norm}[1]{\left\lVert#1\right\rVert}
\DeclareMathOperator{\Pf}{Pf}
\DeclareMathOperator{\Res}{Res}
\DeclareMathOperator{\dist}{dist}
\hypersetup{pdftitle={Natural boundaries of the full square-lattice Ising susceptibility in two complex coupling variables: Endpoint geometry and constrained transport},pdfauthor={Dehao Lin}}
\title{Natural boundaries of the full square-lattice Ising\newline
susceptibility in two complex coupling variables\newline
\large Endpoint geometry and constrained transport}
\author{Dehao Lin\\School of Physics, Sun Yat-sen University\\\small ORCID: \href{https://orcid.org/0009-0001-4551-8490}{0009-0001-4551-8490}\\\small\texttt{lindh9@mail2.sysu.edu.cn}}
\date{October 2026}
\begin{document}
\maketitle
\begin{abstract}
We study the complete connected zero-field response of the infinite
nearest-neighbor square-lattice Ising model with two independent complex
couplings. We give explicit high- and low-temperature domains in
\(\C^2\), identify them as components of an independently specified
spectral complement, construct the physical full-sum continuations there,
and prove local nonextendability at every boundary point. At the selected
regular boundary points the obstruction is stable along arbitrary
nonramified holomorphic transverse curves,
with genuine local uniformity in complex dispersion parameters. Its
negative odd-particle edge mechanism uses exact pole periods and a single
dominant index in selected high Taylor coefficients of the entire sum.
The collinear isotropic case is treated by a separate, explicitly proved
even--odd tail estimate. Exact pullbacks recover both prescribed
fixed-positive-ratio temperature domains for every ratio, including all
exceptional boundary points. The map \(s=\sinh z\) then gives the full
isotropic interior and exterior descendants with their distinct physical
germs. The onsite term, all lattice sites, all particle indices, the
plus-phase subtraction and the physical relation \(\chi=\beta\mathcal X\)
are retained. No conclusion about singularity throughout the interior
of the spectral set is asserted.
\end{abstract}

\section{Main results and the physical observables}
\label{main:statement}

The zero-field susceptibility sums spin correlations over the entire
square lattice. To derive a natural boundary from its form-factor
expansion, singularities of finite-particle integrals must be combined
with estimates for the remaining infinite sum
\cite[Section 3.2]{ONGP}\cite[Section 6]{BHM}.
With independent complex couplings, nonextension on physical temperature
lines also requires a separate argument.

The variables \(K_x,K_y\) are independent coupling coordinates.
Put
\[
 \begin{gathered}
 z=2K_x=x+\ii y,\qquad w=2K_y=u+\ii v,\\
 S=\sinh z,\quad T=\sinh w,\quad A=\cosh z\cosh w,\quad p=ST.
 \end{gathered}
\]
For real positive couplings, \(\mathcal X_H\) is the sum of all
two-spin correlations in the unique high-temperature Gibbs state;
\(\mathcal X_L\) is the sum of all connected two-spin correlations
in the infinite-volume plus phase. Their precise definitions and the
physical interface are in Section~\ref{ana:physical}.

Define independently of any susceptibility
\[
 \Sigma=\{(K_x,K_y):A=Sa+Tb\text{ for some }a,b\in[-1,1]\}.
\]
The two domains used in this paper are explicitly
\begin{align}
 \DH=\bigl\{&|y|<\pi/2,\ |v|<\pi/2,\nonumber\\[-2pt]
       &\cosh x\tanh|u|<\cos y,\quad
        \cosh u\tanh|x|<\cos v\bigr\},\label{main:DH}\\
 \DL=\bigl\{&x>0,\ u>0,\quad
       \cosh x\tanh u>|\cos y|,\quad
       \cosh u\tanh x>|\cos v|\bigr\}.\label{main:DL}
\end{align}
These are open sets in the original coupling space, using the displayed
invertible linear change of coordinates. The HT domain is unbounded
along decoupled directions; boundedness is asserted only for its
positive-ratio temperature slices.

\begin{theorem}[Joint physical natural boundaries]\label{main:joint}
The following assertions hold for both phases of the full bulk reduced
susceptibility.
\begin{enumerate}
\item The sets \(\DH\) and \(\DL\) are precisely the components of
\(\C^2\setminus\Sigma\) containing, respectively, the real positive
regions \(0<p<1\) and \(p>1\). They are contractible and
pseudoconvex. On them \(|p|<1\) and \(|p|>1\), respectively.
\item The physical germs extend to single-valued jointly holomorphic
functions \(\XH\in\mathcal O(\DH)\) and
\(\XL\in\mathcal O(\DL)\). Their complete odd and connected-even
form-factor sums converge normally on compact subsets, with every
fixed mixed complex derivative.
\item At every point of the respective boundary, there is no
holomorphic function on a neighborhood agreeing with that physical
continuation on an accumulating connected component of its domain
side. These are thus natural boundaries also for continuation in
unramified Riemann domains over \(\C^2\), reached from the given
physical germ.
\item The proof gives more than an ambient obstruction. At a dense
set on every regular boundary face, the complete physical sum is
nonextendable along every holomorphic transverse curve on which the
local edge parameter is nonramified. The estimates are locally
uniform for a fixed selected phase and compact admissible subfamilies
of bounded holomorphic complex-parameter families on common polydisks
in a regular-edge chart. No uniformity over the primes or through
\(\Im B=0\) is required. The degenerate equal-coupling restriction has the additional
dense local obstructions of Proposition~\ref{iso:local-obstruction}.
\end{enumerate}
Here regular faces and their edge parameters are specified by
Definition~\ref{geo:regular-definition} and
Proposition~\ref{geo:local-coordinates}. The final assertion is part
of the parent construction needed for the exact restrictions below.
\end{theorem}

For every finite \(\lambda>0\), put
\(\iota_\lambda(\zeta)=(\zeta/2,\lambda\zeta/2)\). Let
\[
 P_\lambda(a)=\cosh a\tanh(\lambda a),\qquad
 Q_\lambda(a)=\cosh(\lambda a)\tanh a,\qquad
 Y_\lambda=\frac{\pi}{2\max(1,\lambda)}.
\]
Both functions on \([0,\infty)\) have increasing inverses. Define
\begin{align*}
 H_\lambda(b)&=\min\{P_\lambda^{-1}(\cos b),
                      Q_\lambda^{-1}(\cos(\lambda b))\},
                 &&|b|<Y_\lambda,\\
 L_\lambda(b)&=\max\{P_\lambda^{-1}(|\cos b|),
                      Q_\lambda^{-1}(|\cos(\lambda b)|)\},
                 &&b\in\R,\\
 \Omega_{H,\lambda}&=\{a+\ii b:|b|<Y_\lambda,\ |a|<H_\lambda(b)\},\\
 \Omega_{L,\lambda}&=\{a+\ii b:a>L_\lambda(b)\}.
\end{align*}

\begin{theorem}[Exact recovery of every positive-ratio slice]
\label{main:slices}
For every finite positive \(\lambda\),
\[
 \iota_\lambda^{-1}(\DH)=\Omega_{H,\lambda},\qquad
 \iota_\lambda^{-1}(\DL)=\Omega_{L,\lambda}.
\]
The restrictions of the two joint physical functions are the respective
physical reduced susceptibilities, holomorphic on these exact domains,
and nonextendable through every boundary point in the componentwise
sense of Theorem~\ref{main:joint}. They retain the complete lattice
sum and every odd or nonzero even particle order. On a physical line
with fixed \(J_x>0,\ J_y=\lambda J_x\), the same conclusions hold
for the energy-field response
\[
       \chi(\zeta)=\frac{\zeta}{2J_x}\,
                         \mathcal X(\iota_\lambda(\zeta)).
\]
The proof uses the
curve-level part of Theorem~\ref{main:joint} and its analytic
construction, rather than inferring a restriction obstruction solely
from ambient nonextension.
\end{theorem}

\begin{corollary}[Exact isotropic descendants]\label{main:isotropic}
Set \(\lambda=1\) and \(s=\sinh\zeta\).
The reduced HT physical germ is single valued and holomorphic on
\(|s|<1\), with value one at zero; the reduced connected-plus-phase
LT physical germ is single valued and holomorphic on \(|s|>1\).
Every point of \(|s|=1\) is a natural-boundary point for each germ.
Their complete normalized representations are
\[
 \XH(s)=\frac{(1-s^4)^{1/4}}{s}
                      \sum_{N\ge1,\ N\ {\rm odd}}C_N(1/s),\qquad
 \XL(s)=(1-s^{-4})^{1/4}
                      \sum_{N\ge2,\ N\ {\rm even}}C_N(s),
\]
where the coefficients are the full-site integrals in
Section~\ref{iso:section}. The LT onsite term is
\(1-(1-s^{-4})^{1/4}\).
The map \(\sinh\) is a biholomorphism from \(\Omega_{H,1}\)
to the disk and is the universal covering from \(\Omega_{L,1}\)
to the exterior, with deck translations \(\zeta\mapsto\zeta+2\pi\ii k\).
The energy-field results follow with \(\beta=\zeta/(2J)\) on that
physical continuation. No single-valued inverse temperature on the
whole exterior is required.
\end{corollary}


\subsection{Related work and proof outline}

The unsquared anisotropic Nickel equation is studied in
\cite[Section 6]{BHM}. Its spectral set has interior; the present
claims concern the two physical complementary components and their
frontiers. No claim of a physical singularity at every point inside
that spectral set is needed.

The only exact-solution imports are real noncritical spin form
factors, together with the published Gibbs-state and ferromagnetic
inequalities identified in Section~\ref{ana:physical}. All complex
continuations and infinite-particle boundary estimates used below
are proved here. 
The geometry follows from endpoint incidence: the endpoint locus
contains the entire spectral boundary, which is enough to identify
the two physical components. Double contours then construct the
jointly holomorphic physical sums. A fixed collision block controls
every higher regular-edge coefficient. At the negative odd edge,
exact residues on the original real angular plane, an angular
concentration estimate and a complex beta period isolate a single
particle index in high Taylor coefficients. On the collinear slice,
one constrained transport based on an analytic variance controls
both parities and all higher orders. The four first-resonance charts
are evaluated on the literal constrained real cycle.

The distinction between joint and restricted continuation is essential.
The exact identity $\XH(K_x,0)=e^{2K_x}$ already prevents an inference
from ambient nonextension to nonextension on arbitrary lines.
For each positive physical ratio, the proof instead supplies a dense
set of actual nonramified slice transversals and retains their
holomorphic anisotropy $B$. The equal-coupling case uses the separate
complete-sum argument of Section~\ref{iso:section}.

\par\medskip\noindent\textbf{AI Declaration.} ChatGPT was used extensively in the mathematical research,
proof development, and preparation of this manuscript.\par

\clearpage
\tableofcontents
\clearpage


\section{The physical components from endpoint incidence}
\label{geo:section}

Put
\[
 z=2K_x=x+iy,\quad w=2K_y=u+iv,\quad
 S=\sinh z,\quad T=\sinh w,\quad A=\cosh z\cosh w,\quad p=ST,
\]
and retain the unsquared spectral set
\begin{equation}\label{geo:sigma}
 \Sigma=\{(z,w): A=Sa+Tb\text{ for some }a,b\in[-1,1]\}.
\end{equation}
The change between $(z,w)$ and the coupling coordinates is invertible.
Compactness of the representing square makes $\Sigma$ closed.
Define
\begin{equation}\label{geo:F}
 F_1=\cosh^2u\cos^2y-\cosh^2x\sinh^2u,\qquad
 F_2=\cosh^2x\cos^2v-\cosh^2u\sinh^2x,
 \qquad E=\{F_1F_2=0\}.
\end{equation}
Only the boundary of the spectral set is needed to determine its two
physical complementary components.  The following endpoint criterion
avoids a classification of the spectral interior.

\begin{lemma}[Endpoint boundary criterion]\label{geo:endpoint-boundary}
The set $E$ is exactly the set of parameters admitting a representation
in \eqref{geo:sigma} with $a$ or $b$ equal to $1$ or $-1$.
Consequently
\begin{equation}\label{geo:endpoint-sandwich}
 E\subset\Sigma,\qquad \partial\Sigma\subset E.
\end{equation}
These statements hold for independent complex $z,w$, including zero
real parts, zeros of $S$ or $T$, and real-linear dependence of $S,T$.
\end{lemma}
\begin{proof}
If $S\ne0$, $\tanh(z/2)$ is finite and nonzero.  For
$\epsilon\in\{1,-1\}$ put $r_\epsilon=e^{\epsilon w}\tanh(z/2)$.
Direct multiplication gives
\begin{equation}\label{geo:endpoint-identities}
 \frac{r_\epsilon+r_\epsilon^{-1}}2=\frac{A-\epsilon T}{S},
 \qquad
 |\tanh(z/2)|^2=\frac{\cosh x-\cos y}{\cosh x+\cos y}.
\end{equation}
Both terms in the last ratio are positive when $S\ne0$.  Therefore
\[
 |r_\epsilon|=1
 \quad\Longleftrightarrow\quad
 \cos y=\epsilon\cosh x\tanh u.
\]
The preimage of $[-1,1]$ under $(r+r^{-1})/2$ is the unit circle.
Thus $A=Sa+\epsilon T$ with $a\in[-1,1]$ is equivalent to the
last equation.  Taking the union over $\epsilon$ gives exactly $F_1=0$.
If $S=0$, then $x=0$, $y\in\pi\mathbb Z$, and $F_1=1$.
Also $A^2-T^2=1$, so no such endpoint representation exists.
This proves the equivalence globally.  Exchange $z,w$ for $F_2$.

For $P\in\partial\Sigma$, consider the real-linear map
$L_P(a,b)=S(P)a+T(P)b$.  If it has rank two and the representing
pair is in the open square, continuity of $L_K^{-1}A(K)$ would make
$P$ an interior point of $\Sigma$.  Its pair must therefore have an
endpoint coordinate.  If the rank is at most one, the affine fiber
of any representing pair has a nonzero kernel direction; following
that direction inside the closed square reaches its boundary.
There is again an endpoint representation.  The equivalence just
proved gives $P\in E$.
\end{proof}

\subsection{The two entire physical components}
For $a,b\ge0$ write
\[
 r_1(a,b)=\cosh a\tanh b,\qquad r_2(a,b)=\cosh b\tanh a.
\]
Both functions are nondecreasing in each argument and strictly
increasing in their hyperbolic-tangent argument.  Also
\begin{equation}\label{geo:ri-one}
 r_1(a,b)^2-1=\frac{\sinh^2a\sinh^2b-1}{\cosh^2b},\qquad
 r_2(a,b)^2-1=\frac{\sinh^2a\sinh^2b-1}{\cosh^2a}.
\end{equation}

\begin{theorem}[Physical components and branches]\label{geo:components}
The component of $\mathbb C^2\setminus\Sigma$ containing the positive
real region $0<p<1$ is
\begin{equation}\label{geo:DH}
 \mathcal D_H=\{ |y|,|v|<\pi/2,\quad
 r_1(|x|,|u|)<\cos y,\quad r_2(|x|,|u|)<\cos v\}.
\end{equation}
The component containing the positive real region $p>1$ is
\begin{equation}\label{geo:DL}
 \mathcal D_L=\{x,u>0,\quad r_1(x,u)>|\cos y|,\quad
                              r_2(x,u)>|\cos v|\}.
\end{equation}
Both domains are contractible, and
\begin{equation}\label{geo:p-domains}
 |p|<1\text{ on }\mathcal D_H,\qquad |p|>1\text{ on }\mathcal D_L.
\end{equation}
\end{theorem}
\begin{proof}
Simultaneous dilation of all four real coordinates toward zero
preserves the strict inequalities of $\mathcal D_H$: their left
sides decrease and their right sides increase.  Hence this domain
is star shaped about zero.  Its closure is exactly
\begin{equation}\label{geo:HT-closure}
 |y|,|v|\le\pi/2,\quad
 r_1(|x|,|u|)\le\cos y,\quad r_2(|x|,|u|)\le\cos v.
\end{equation}
Indeed, every point satisfying these weak inequalities is moved
into the strict domain by every dilation factor $t<1$.
For an active first inequality, either $u\ne0$, in which case
$\tanh(t|u|)$ strictly decreases, or $u=0$ and $|y|=\pi/2$,
in which case $\cos(ty)>0$; the second inequality is the exchanged
calculation.  Thus every finite boundary point has $F_1F_2=0$.

For $\mathcal D_L$, increasing both real parts strictly opens every
weak inequality, so its closure is exactly
\begin{equation}\label{geo:LT-closure}
 x,u\ge0,\quad r_1(x,u)\ge|\cos y|,\quad
                         r_2(x,u)\ge|\cos v|.
\end{equation}
If both $F_i$ are strictly negative at such a point, then $x,u>0$
and the point is interior.  Its finite boundary is therefore also
in $E$.  To contract the domain, fix $R>0$ with $\sinh R>1$.
First increase $(x,u)$ to $(\max(x,R),\max(u,R))$, then move to
$(R,R)$ within $[R,\infty)^2$, and finally contract $(y,v)$ to
$(0,0)$.  The first move preserves the inequalities by monotonicity;
during the other two, both $r_i$ are at least $\sinh R>1$.
These three continuous homotopies give a contraction.

On the respective domains both $F_i$ have a strict common sign,
so neither domain meets $E$.  By \eqref{geo:endpoint-sandwich},
its intersection with $\Sigma$ is relatively open and closed.
Both domains are connected.  Zero is outside $\Sigma$, and the
positive real point $(R,R)$ above is outside $\Sigma$, since
\[
 (A-S-T)(A+S+T)=(p-1)^2>0
\]
there and $A+S+T>0$.  Hence both domains are disjoint from $\Sigma$.
Their boundaries lie in $\Sigma$, so each is an entire complementary
component.  Equation \eqref{geo:ri-one} identifies the indicated
positive real regions.

Finally,
\begin{equation}\label{geo:p-modulus}
 |p|^2=(\cosh^2x-\cos^2y)(\cosh^2u-\cos^2v).
\end{equation}
If $F_1,F_2>0$, the two factors are respectively less than
$\cosh^2x/\cosh^2u$ and $\cosh^2u/\cosh^2x$; if both are negative,
both inequalities reverse.  This proves \eqref{geo:p-domains}.
\end{proof}

The binomial series define $(1-p^2)^{1/4}$ on $\mathcal D_H$ and
$(1-p^{-2})^{1/4}$ on $\mathcal D_L$, with their positive real
physical values.  The HT transfer zeros are exactly the complex
axes $z=0$ and $w=0$; neither transfer variable vanishes on the LT
domain.  No assertion that $\mathcal D_H$ is bounded is made.
The corner identity
\begin{equation}\label{geo:corner}
 A^2-(\epsilon S+\delta T)^2=(p-\epsilon\delta)^2,
 \qquad \epsilon,\delta\in\{1,-1\},
\end{equation}
will also be useful.

\begin{lemma}[Canonical decaying root]\label{geo:canonical-root}
For $(z,w)\notin\Sigma$ and real $q$, the equation
\begin{equation}\label{geo:root-equation}
 S\rho^2-2(A-T\cos q)\rho+S=0
\end{equation}
has a unique root with $|\rho|<1$, interpreted as zero at $S=0$.
It is holomorphic in both couplings, locally jointly holomorphic
with the complex angle, and has a uniform modulus gap below one
on each compact parameter set and the real angular circle.
On either positive real physical region it is the positive decaying root.
\end{lemma}
\begin{proof}
For $S\ne0$, $W=(A-T\cos q)/S$ is outside $[-1,1]$ by the unsquared
definition of $\Sigma$.  The reciprocal roots of
$\rho+\rho^{-1}=2W$ are distinct and neither is on the unit circle;
exactly one is inside.  Its local implicit-function branches agree
by uniqueness.  At $S=0$, $A-T\cos q\ne0$, so the implicit-function
theorem extends that root through zero.  Compactness gives the gap
and a common thin complex angular neighborhood.  On positive real
noncritical parameters, the corner identity gives $A-S-T>0$, hence
$W>1$ and $0<\rho<1$.
\end{proof}

\begin{proposition}[Pseudoconvexity]\label{geo:pseudoconvex}
Every component of $\mathbb C^2\setminus\Sigma$, in particular
$\mathcal D_H$ and $\mathcal D_L$, has a continuous
plurisubharmonic exhaustion.
\end{proposition}
\begin{proof}
The continuous function
\[
 m(K)=\min_{a,b\in[-1,1]}|A-Sa-Tb|
\]
is positive on a complementary component and zero exactly on $\Sigma$.
There $-\log m$ is the continuous locally bounded supremum of the
pluriharmonic functions $-\log|A-Sa-Tb|$, so it is plurisubharmonic
by the submean inequality on complex lines.  Thus
\[
 |K_x|^2+|K_y|^2+\max(0,-\log m(K))
\]
is a continuous plurisubharmonic exhaustion: it diverges at infinity
and at every finite boundary point.  A bounded sublevel has $m$ bounded
away from zero, so all its limit points remain in the same component.
\end{proof}

\subsection{Regular joint faces and their density}
\begin{definition}\label{geo:regular-definition}
A boundary point is a regular joint edge when $S,T\ne0$,
$D=\Im(\overline S T)\ne0$, exactly one of its representing
coefficients is an endpoint, the other lies in $(-1,1)$, and $p^2\ne1$.
\end{definition}

For face regularity we use the following endpoint identity.  Suppose $F_1=0$ and use its
endpoint representation $A=Sa+\epsilon T$, $a\in[-1,1]$.
Set $K=\cosh z\,T-\epsilon\cosh w$.  Then
\[
 K^2=S^2(a^2-1),\qquad
 \cos y=\epsilon\cosh x\tanh u.
\]
Writing $c=\cosh x$, $C=\cosh u$, $H=\sinh u$, we obtain
\[
 \Re(\overline K T)
 =c\cos y\,|T|^2-\epsilon HC
 =\epsilon\frac HC\{c^2(C^2-\cos^2v)-C^2\}
 =-\epsilon\tanh u\,F_2.
\]
Since $K/S=\pm i\sqrt{1-a^2}$, including the zero case, this proves
\begin{equation}\label{geo:strict-face}
 D^2(1-a^2)=\tanh^2u\,F_2^2.
\end{equation}
Also direct substitution in \eqref{geo:p-modulus} gives
\begin{equation}\label{geo:p-edge}
 |p|^2=1-\frac{F_2}{\cosh^2u}\quad(F_1=0),\qquad
 |p|^2=1-\frac{F_1}{\cosh^2x}\quad(F_2=0).
\end{equation}
Thus $x,u\ne0$ and exactly one vanishing $F_i$ imply all the
regular-edge conditions on either physical boundary.

\begin{proposition}[Density on the complete joint boundary]
\label{geo:joint-density}
Points with $x,u\ne0$ and exactly one vanishing $F_i$ are dense in
both $\partial\mathcal D_H$ and $\partial\mathcal D_L$.
In particular regular joint edges are dense, including near all
zero-transfer, zero-real-part, collinear and intersection strata.
\end{proposition}
\begin{proof}
For the HT closure \eqref{geo:HT-closure},
$\sinh|x|\sinh|u|\le1$ by \eqref{geo:ri-one}.  For fixed real-part
magnitudes $a,b$ with product strictly below one, its angular fiber is
\[
 |y|\le\arccos r_1(a,b),\qquad |v|\le\arccos r_2(a,b).
\]
At product one both intervals collapse to zero.  At any boundary point
at least one inequality is active.  Approximate its magnitudes by
$a_j,b_j>0$ with $\sinh a_j\sinh b_j<1$.  If the first face is active,
choose $|y_j|=\arccos r_1(a_j,b_j)$ with $y_j\to y$, and choose
$v_j$ strictly inside its second interval with $v_j\to v$.
Preserve each nonzero real-part sign, and choose either sign at a
zero real part.  This gives regular points $F_1=0,F_2>0$ approaching
the prescribed point.  Exchange the faces when the second is active.
The construction also covers strip endpoints and collapsed fibers.

A point of the LT boundary has $\sinh x\sinh u\le1$, since a larger
product would make both inequalities strict.  It has an active cosine
inequality: $x=0$ forces the second active and $u=0$ the first.
Choose $x_j,u_j>0$ with product below one approaching the given point.
For an active first face choose
\[
 |\cos y_j|=r_1(x_j,u_j),\quad y_j\to y,
 \qquad |\cos v_j|<r_2(x_j,u_j),\quad v_j\to v.
\]
Continuity of the local cosine intervals permits these choices also
when an endpoint tends to a cosine zero or to a multiple of $\pi$.
They give $F_1=0,F_2<0$.  Exchange the faces for the other case.
\end{proof}

\begin{proposition}[Independent holomorphic edge coordinates]
\label{geo:local-coordinates}
Where $S\ne0$, put $C=A/S$ and $B=T/S$.  Then
\begin{equation}\label{geo:Jacobian}
 \det\frac{\partial(C,B)}{\partial(z,w)}=\frac{p^2-1}{S^3}.
\end{equation}
At a regular edge these are local holomorphic coordinates and
$\Im B\ne0$.  At an edge with $b=\epsilon$, the coordinate
$h=C-\epsilon B$ has boundary $\Im h=0$, $\Re h\in(-1,1)$;
the physical component is the entire local complementary side.
For $a=\epsilon$, exchange transfer axes.
On every regular joint boundary patch, points whose free coefficient is
$\cos(2\pi k/q)$, with $q$ an odd prime and $1\le k\le(q-1)/2$,
are dense.
\end{proposition}
\begin{proof}
The derivatives
\[
 C_z=-\frac{\cosh w}{S^2},\quad C_w=\frac{\cosh z\,T}{S},\quad
 B_z=-\frac{T\cosh z}{S^2},\quad B_w=\frac{\cosh w}{S}
\]
give \eqref{geo:Jacobian}; also $D=|S|^2\Im B$.
In these coordinates the spectral condition is the parallelogram
$C\in[-1,1]+B[-1,1]$.  Near the interior of its $b=\epsilon$ edge,
\[
 b-\epsilon=\frac{\Im(C-\epsilon B)}{\Im B},
\]
while the other coefficient remains strictly interior.  The two local
sides are therefore the spectral side and one connected complementary
side.  The latter meets the physical component and lies entirely in it.
Holding $B$ fixed makes $h$ range over a real interval.  Prime-grid
angles have arbitrarily small mesh, so their cosines are dense there.
\end{proof}

On an HT $F_1$ face the endpoint is $b=\operatorname{sign}u$;
on an LT $F_1$ face it is $b=\operatorname{sign}(\cos y)$.
On $F_2$ faces exchange $x,y$ with $u,v$.  Both endpoint signs
must consequently be retained in the analytic boundary theorem.

\subsection{Every positive-ratio slice}
For $\lambda>0$ define
\[
 P_\lambda(a)=\cosh a\tanh(\lambda a),\quad
 Q_\lambda(a)=\cosh(\lambda a)\tanh a,\quad
 Y_\lambda=\frac{\pi}{2\max(1,\lambda)}.
\]
These functions are odd analytic bijections of $\mathbb R$ with
strictly positive derivatives.  Their restrictions to $[0,\infty)$
thus have increasing inverses.  Write
\[
 H_\lambda(b)=\min\{P_\lambda^{-1}(\cos b),Q_\lambda^{-1}(\cos(\lambda b))\},
 \quad |b|<Y_\lambda,
\]
\[
 L_\lambda(b)=\max\{P_\lambda^{-1}(|\cos b|),Q_\lambda^{-1}(|\cos(\lambda b)|)\}.
\]

\begin{theorem}[Exact pullbacks]\label{geo:exact-pullbacks}
For $\iota_\lambda(\zeta)=(\zeta/2,\lambda\zeta/2)$,
\begin{align}
 \iota_\lambda^{-1}(\mathcal D_H)
 &=\Omega_{H,\lambda}=\{a+ib:|b|<Y_\lambda,\ |a|<H_\lambda(b)\},
 \label{geo:HT-pullback}\\
 \iota_\lambda^{-1}(\mathcal D_L)
 &=\Omega_{L,\lambda}=\{a+ib:a>L_\lambda(b)\}.
 \label{geo:LT-pullback}
\end{align}
These domains are simply connected.  If $z_c>0$ is the unique solution
of $\sinh z_c\sinh(\lambda z_c)=1$, then
$P_\lambda(z_c)=Q_\lambda(z_c)=1$ and $0\le L_\lambda\le z_c$.
The HT profile is positive in its interval and tends to zero at both ends.
\end{theorem}
\begin{proof}
For example
\[
 P_\lambda'(a)=\sinh a\tanh(\lambda a)+\lambda\cosh a\operatorname{sech}^2(\lambda a)>0;
\]
the exchanged derivative has the same property.  Substitute
$(x,y,u,v)=(a,b,\lambda a,\lambda b)$ into the exact component
inequalities and invert these two increasing functions.  This gives
both pullbacks.  The identities \eqref{geo:ri-one} give the claims
about $z_c,L_\lambda$.  At an HT endpoint at least one cosine
vanishes, proving the limit of $H_\lambda$.  The HT domain is star
shaped by dilation.  The LT map $(a,b)\mapsto(a-L_\lambda(b),b)$
is a homeomorphism to $(0,\infty)\times\mathbb R$.
\end{proof}

\begin{proposition}[Same-slice regular density]\label{geo:slice-density}
For every $\lambda>0$, $\lambda\ne1$, all but a locally finite subset
of either physical slice boundary are regular joint edges, and their
edge maps are nonramified on the actual slice.  The selected prime
cosines are dense on every remaining regular arc and hence on both
entire slice boundaries.
\end{proposition}
\begin{proof}
Use the signed real-analytic profiles
\[
 R(b)=P_\lambda^{-1}(\cos b),\qquad
 U(b)=Q_\lambda^{-1}(\cos(\lambda b)),\qquad
 E_\lambda(b)=R(b)^2-U(b)^2.
\]
Here the inverses are those of the odd bijections on all of $\mathbb R$.
Thus $E_\lambda$ is analytic on all of $\mathbb R$, including every
simultaneous cosine zero.  Near zero,
\begin{align*}
 E_\lambda(b)
 &=z_c\left(\frac{\lambda^2}{Q_\lambda'(z_c)}-
                  \frac1{P_\lambda'(z_c)}\right)b^2+O(b^4)\\
 &=z_c\{\lambda\coth(\lambda z_c)-\coth z_c\}b^2+O(b^4).
\end{align*}
For the second equality put $h=\sinh z_c$, so
$\sinh(\lambda z_c)=h^{-1}$,
$P_\lambda'(z_c)=h(1+\lambda h)/\cosh z_c$, and
$Q_\lambda'(z_c)=(1+\lambda h)/(h\cosh z_c)$.
The displayed coefficient is nonzero when $\lambda\ne1$, because
$t\coth t$ is strictly increasing on $t>0$; its derivative has
positive numerator $\sinh t\cosh t-t$.
Hence the zeros of $E_\lambda$ are locally finite.

Except at these profile contacts and the two HT tips, a boundary
point has $a\ne0$ and exactly one vanishing $F_i$.  The strict-face
identity and \eqref{geo:p-edge} give all joint regularity conditions.
There is no separate rational-ratio exception: all such contacts
are already zeros of the same analytic function.

The edge map on an $F_1$ face is
\begin{equation}\label{geo:slice-edge-maps}
 g_\epsilon(\zeta)=\frac{A-\epsilon T}{S}
   =\frac{r_\epsilon+r_\epsilon^{-1}}2,
 \qquad r_\epsilon=e^{\epsilon\lambda\zeta}\tanh(\zeta/2).
\end{equation}
At that face,
\[
 \frac{r_\epsilon'}{r_\epsilon}=\epsilon\lambda+\frac1S,
 \qquad
 \Re\frac{r_\epsilon'}{r_\epsilon}
 =\epsilon\left(\lambda+
       \frac{\sinh a\sinh(\lambda a)\cosh(\lambda a)}{\cosh a}\right)\ne0.
\]
The expression in parentheses is positive also when $a<0$.
Strictness of the free coefficient implies $r_\epsilon\ne\pm1$,
so $g_\epsilon'\ne0$.  On an $F_2$ face use
$r=e^{\epsilon\zeta}\tanh(\lambda\zeta/2)$; there
\[
 \Re\frac{r'}r=\epsilon\left(1+
 \lambda\frac{\sinh a\sinh(\lambda a)\cosh a}{\cosh(\lambda a)}\right)\ne0.
\]
Consequently each regular boundary arc maps biholomorphically to
an interval of free coefficients.  Prime-cosine density applies to
this actual arc.  The two boundaries are continuous graphs or a
continuous lens boundary and have no isolated points; deleting their
locally finite exceptional subsets leaves a dense set.
\end{proof}

At a regular selected slice point $\zeta_0\ne0$, the map
$(\zeta,\lambda)\mapsto(z,w)=(\zeta,\lambda\zeta)$ has determinant
$\zeta_0$.  Composing with the edge chart and using the just-proved
nonzero edge derivative gives, by the holomorphic implicit-function
theorem, the actual anisotropy $B=B(t,\lambda)$ on a product polydisk,
where $t$ is the nonramified edge coordinate.  Shrinking a slightly
larger product polydisk preserves the sign of $\Im B$ and supplies
common Cauchy bounds.  Thus the local analytic boundary estimates
can be applied to the variable-$B$ physical slices and their complex
parameter neighborhoods.  This construction makes no inference
from ambient nonextension to restriction nonextension.

\subsection{The local collar and propagation to exceptional points}
\begin{lemma}[Inward-map collar]\label{geo:collar}
Let $D$ be $\mathcal D_H$, $\mathcal D_L$, or either of their
positive-ratio pullbacks.  If $P\in\partial D$ and $U$ is an open
neighborhood of $P$, there is a neighborhood $V\subset U$ of $P$
such that $D\cap V$ lies in one connected component of $D\cap U$.
Every component of $D\cap U$ accumulating at $P$ contains $D\cap V$.
\end{lemma}
\begin{proof}
For HT use $H_s(Q)=(1-s)Q$; for LT in $(x,u,y,v)$ coordinates use
$H_s(Q)=Q+s(1,1,0,0)$.  These maps preserve $D$ for $s\ge0$
sufficiently small and move every boundary point strictly inside
for $s>0$, by the closure arguments above.  On a slice use the
same HT dilation and the LT translation $\zeta\mapsto\zeta+s$,
corresponding to the real direction $(1,\lambda,0,0)$.

Choose a small fixed $s>0$ so that $H_t(P)\in U$ for $0\le t\le s$,
and choose a ball around $H_s(P)$ contained in $D\cap U$.
By continuity and compactness of $[0,s]$, after shrinking $V$,
every path $H_t(Q)$, $Q\in D\cap V$, stays in $D\cap U$ and its
endpoint belongs to that ball.  The paths and straight segments
inside the ball join any two such $Q$.  An accumulating component
meets $D\cap V$ and hence is this same component.
\end{proof}

\begin{definition}\label{geo:componentwise-extension}
For $f\in\mathcal O(D)$, a local holomorphic extension through
$P\in\partial D$ is a function holomorphic on a neighborhood $U$
of $P$ that agrees with $f$ on a connected component of $D\cap U$
accumulating at $P$.
\end{definition}

\begin{corollary}[Dense obstructions imply every componentwise obstruction]
\label{geo:propagate-obstruction}
For any domain $D$ in the preceding lemma, suppose $f\in\mathcal O(D)$
has no holomorphic extension agreeing on the full local domain side
at a dense subset of $\partial D$.  Then no neighborhood of any
$P\in\partial D$ admits a holomorphic function agreeing with $f$
on an accumulating connected component of its intersection with $D$.
\end{corollary}
\begin{proof}
An extension agreeing on such a component agrees on $D\cap V$ by
the collar lemma.  Choose a dense obstructing point in $V$ and
restrict the extension to a neighborhood contained in $V$.
\end{proof}

\subsection{The isotropic disk and exterior covering}
\begin{lemma}[Exact isotropic geometry]\label{geo:isotropic-cover}
At $\lambda=1$, with $s=\sinh\zeta$,
\[
 \iota_1^{-1}(\Sigma)=\{|\sinh\zeta|=1\},\quad
 \Omega_{H,1}=\{|b|<\pi/2,\ |a|<\operatorname{arsinh}(\cos b)\},\quad
 \Omega_{L,1}=\{a>\operatorname{arsinh}(|\cos b|)\}.
\]
The map $\sinh$ is a biholomorphism from $\Omega_{H,1}$ onto the
unit disk and a universal covering from $\Omega_{L,1}$ onto the
unit-circle exterior, with deck translations $\zeta\mapsto\zeta+2\pi i k$.
\end{lemma}
\begin{proof}
Here $S=T=s$ and $A=1+s^2$.  For $s\ne0$ the unsquared condition
is $s+s^{-1}\in[-2,2]$, equivalent to $|s|=1$; $s=0$ is not
spectral.  Since $P_1=Q_1=\sinh$, the exact profiles follow.
The identity $|\sinh(a+ib)|^2=\sinh^2a+\sin^2b$ maps the bounded
HT lens into the disk and its boundary into the circle.  The map
is injective on $|\Im\zeta|<\pi/2$: equal hyperbolic sines require
$\zeta_1=\zeta_2+2\pi i k$ or
$\zeta_1=-\zeta_2+(2k+1)\pi i$, and only the first with $k=0$
is possible.  Its derivative has no zero there.  Its image is
nonempty, open and relatively closed in the disk by the boundary
identity and boundedness, hence is the whole disk.

For the exterior put $\omega=e^\zeta$.  The equation
$\omega-\omega^{-1}=2s$ has exactly one root with $|\omega|>1$
when $|s|>1$: the root moduli have product one and a unit root
would give $s\in i[-1,1]$.  The outside root is holomorphic,
using $\sqrt{s^2+1}=s\sqrt{1+s^{-2}}$ with its binomial branch.
Its image is
\[
 \{re^{ib}: r>\exp(\operatorname{arsinh}|\cos b|)\},
\]
by the same modulus identity.  This root map is a biholomorphism.
Its full exponential preimage is the simply connected LT epigraph,
so exponentiation supplies the universal covering and the stated
deck group.  The reduced susceptibility requires its own termwise
deck invariance to descend; inverse temperature remains on this cover.
\end{proof}

\section{The physical observable and joint holomorphy}
\label{ana:physical}

Throughout this paper \(z=2K_x\) and \(w=2K_y\) are independent
complex coordinates. Write
\[
 S=\sinh z,\qquad T=\sinh w,\qquad
 A=\cosh z\cosh w,\qquad p=ST.
\]
No inverse temperature is introduced as an additional independent
variable in \(\C^2\). For real positive \(K_x,K_y\), the
dimensionless uniform field is denoted by \(\eta\). The zero-field
responses are
\begin{align}
 \XH(K_x,K_y)
 &=\sum_{(m,l)\in\Z^2}\langle\sigma_{0,0}\sigma_{m,l}\rangle,
 &&0<p<1, \label{ana:physicalHT}\\
 \XL(K_x,K_y)
 &=\sum_{(m,l)\in\Z^2}
   \bigl(\langle\sigma_{0,0}\sigma_{m,l}\rangle_+-M^2\bigr),
 &&p>1, \label{ana:physicalLT}\\
 M^2&=(1-p^{-2})^{1/4}. \nonumber
\end{align}
The high-temperature state is the unique Gibbs state. The subscript
\(+\) means the infinite-volume plus state; the fourth root in the last
line is positive. These definitions include the origin, both axes and
all quadrants. In particular the low-temperature onsite contribution
is \(1-M^2\), not \(1\).

\subsection{Real-temperature form factors and state identification}
\label{ana:interface}

Only the real, noncritical spin form factors are imported from the
exact solution. The relevant finite-cylinder formulas at arbitrary
separation are \cite[(35)--(39)]{BL}; their finite-lattice matrix-element
factorization is proved in \cite[Theorem 6.1]{IL}, with the
paramagnetic vacuum change in its Remark 6.2. In our convention put
\[
 B=T/S>0,\quad C=A/S,\quad
 \cosh\gamma(q)=C-B\cos q,\quad \gamma(q)>0.
\]
Then the infinite-plane coefficient is
\begin{equation}\label{ana:ff}
 f_N(m,l)=\frac{B^{N^2/2}}{N!}
 \int_{\T^N}\prod_{i=1}^N
 \frac{e^{-|m|\gamma(q_i)+\ii lq_i}}{\sinh\gamma(q_i)}
 \prod_{i<j}
 \left\{\frac{\sin((q_i-q_j)/2)}
 {\sinh((\gamma(q_i)+\gamma(q_j))/2)}\right\}^{\!2}
 \prod_{i=1}^N\frac{\dd q_i}{2\pi}.
\end{equation}
With positive real square roots the two correlations are
\begin{align}
 \langle\sigma_0\sigma_{m,l}\rangle
 &=\xi_H\sum_{\substack{N\ge1\\N\ {\rm odd}}}f_N(m,l),
 &\xi_H&=(1-p^2)^{1/4}p^{-1/2},\label{ana:corHT}\\
 \langle\sigma_0\sigma_{m,l}\rangle_+
 &=\xi_L\left(1+\sum_{\substack{N\ge2\\N\ {\rm even}}}f_N(m,l)\right),
 &\xi_L&=(1-p^{-2})^{1/4}=M^2.\label{ana:corLT}
\end{align}
This is the normalization of \cite[(1.8)--(1.10)]{ONGP}.
In particular the anisotropy exponent in \eqref{ana:ff} is \(N^2/2\).

Here are the limit and state identifications needed in using that
interface. On a real noncritical compact set the dispersion has a
uniform positive gap and is analytic in an angular strip. The
finite-circumference corrections \(\eta,\Lambda^{-1},\log\xi_T\)
in \cite{BL} tend uniformly to zero, by the exponentially convergent
trapezoidal rule applied to their sum-minus-integral expressions.
Here is a quantitative version that also covers the highest particle
indices. If \(L\) is the circumference, the unnormalized Ramond and
Neveu--Schwarz sums of a strip-analytic function differ by
\(O(Le^{-hL})\): their Fourier expansions select the multiples of
\(L\), with respective factors \(1\) and \((-1)^k\).
Apply this to \(\gamma\) and to
\(\log(1-e^{-\gamma(q)-\gamma(q')})\). The mixed double-grid
difference of \(\log\sinh((\gamma(q)+\gamma(q'))/2)\) in
\cite[(25)]{BL} gives
\[
 |\Lambda_L^{-1}|+\sup_q|\eta_L(q)|\le CLe^{-hL},\qquad
 |\log\xi_{T,L}|\le CL^2e^{-2hL}.
\]
Thus at each fixed separation the entire correction is bounded by
\(C_{m,l}C^N\), uniformly in \(L,N\). Terms with \(N>L\)
vanish on the discrete grid. In particular no estimate at a fixed
particle number is being used to pass an unbounded particle sum
through the thermodynamic limit.
For the real roots \(0<\rho(q)=e^{-\gamma(q)}\le r<1\),
\[
 \left|\frac{\rho(q)-\rho(q')}{1-\rho(q)\rho(q')}\right|\le r<1.
\]
The matching argument in Lemma~\ref{ana:matching} is an algebraic
identity on discrete angular grids as well as on continuous circles.
It gives a common bound
\(\exp\{-cN^2+O(N\log(N+1))\}\) at fixed separation, including
the finite-size one-body corrections and factorial normalization.
Consequently the particle sum passes through the cylinder limit;
each fixed-index Riemann sum gives \eqref{ana:ff}.

Every local cylinder limit is a translation-invariant and spin-flip
invariant Gibbs state. The Gibbs-state theorem
\cite[Theorem 1]{Raoufi} applies with the vertex-transitive group
chosen to be translations: the two unequal positive nearest-neighbor
couplings satisfy its nonnegativity, invariance, irreducibility and
summability hypotheses. That state is the equal mixture of the plus
and minus states. Its two-spin correlations equal those of the plus
state. This proves the state identification in \eqref{ana:corLT};
subtracting \(M^2=\xi_L\) removes exactly its zero-particle term.
No response of a zero-mean mixed state is substituted for
\eqref{ana:physicalLT}.

\subsection{Double contours and the full lattice}
\label{ana:double}

Let \(\zeta,\eta\) now denote independent contour variables, and set
\[
 \mathcal D(\zeta,\eta;z,w)=
 A-\frac S2(\zeta+\zeta^{-1})
  -\frac T2(\eta+\eta^{-1}),\qquad
 a(v,v')=\frac{v-v'}{1-vv'}.
\]
If \(X=\prod_i\zeta_i\) and \(Y=\prod_i\eta_i\), the same coefficient
has the symmetric representation
\begin{equation}\label{ana:doublecorr}
 f_N(m,l)=\frac{p^{N/2}}{N!}\oint\cdots\oint
 X^{|m|-1}Y^{|l|-1}
 \prod_{i<j}a(\zeta_i,\zeta_j)a(\eta_i,\eta_j)
 \prod_{i=1}^N
 \frac{\dd\zeta_i\,\dd\eta_i}{(2\pi\ii)^2\mathcal D(\zeta_i,\eta_i)}.
\end{equation}
The circles have radii strictly less than and sufficiently close to
one. This formula can be checked directly on the real physical
region, as follows. At \(\eta=e^{\ii q}\) the inside root is
\(\zeta=\rho=e^{-\gamma(q)}\), and
\[
 \mathcal D_\zeta(\rho,\eta)
   =\frac{S(1-\rho^2)}{2\rho^2},\qquad
 \frac{\sqrt p}{\mathcal D_\zeta(\rho,\eta)}
   =\sqrt B\,\frac{\rho}{\sinh\gamma}.
\]
Subtracting two dispersion equations gives the exact canceled pair
\begin{equation}\label{ana:pair}
 a(\rho_i,\rho_j)a(\eta_i,\eta_j)
 =-B\frac{(\eta_i-\eta_j)^2\rho_i\rho_j}
          {\eta_i\eta_j(1-\rho_i\rho_j)^2}
 =B\left\{\frac{\sin((q_i-q_j)/2)}
       {\sinh((\gamma_i+\gamma_j)/2)}\right\}^{\!2}
 =:P_{ij}.
\end{equation}
The residues therefore give \(B^{N/2}\), and the pairs give
\(B^{N(N-1)/2}\). This proves the full exponent and the normalization
in \eqref{ana:ff}. Reflection of either spatial coordinate is
legitimate; equality under interchanging the separations while holding
unequal couplings fixed is not required.

The exact product of the two spatial geometric series is
\begin{equation}\label{ana:fullsites}
 \sum_{m,l\in\Z}X^{|m|-1}Y^{|l|-1}
 =\frac{(1+X)(1+Y)}{XY(1-X)(1-Y)}
 =\frac1{XY}+\frac{2(X^{-1}+Y^{-1})}{(1-X)(1-Y)}.
\end{equation}
The first term is the onsite coefficient. Denoting by \(T_N\)
the integral with the second numerator but without its factor two
gives \(C_N=f_N(0,0)+2T_N\), where \(C_N\) sums all sites.
The angular form of this complete coefficient is
\begin{equation}\label{ana:angular}
 C_N=\frac1{N!}
 \int_{\sum_iq_i=0\ ({\rm mod}\ 2\pi)}
 G_N(q)\frac{1+Z}{1-Z}
 \prod_{i=1}^{N-1}\frac{\dd q_i}{2\pi},\qquad
 G_N=B^{N/2}\prod_i\frac1{\sinh\gamma_i}\prod_{i<j}P_{ij},
 \quad Z=\prod_i\rho_i.
\end{equation}
There are \(N-1\), not \(N\), independent normalized measures in this
formula; \(q_N=-\sum_{i<N}q_i\) modulo \(2\pi\).


\begin{lemma}[Scalar matching identity]\label{ana:schur-identity}
For even $N$,
\[
 \operatorname{Pf}\left[\frac{v_i-v_j}{1-v_iv_j}\right]_{i,j=1}^N
       =\prod_{i<j}\frac{v_i-v_j}{1-v_iv_j}.
\]
For odd $N$, adjoin $v_{N+1}=-1$; its upper-column entries are
the removable constant one and the product on the right is the
product over the original indices.
\end{lemma}
\begin{proof}
Set $a_i=(1+v_i)/(1-v_i)$, so the entries become
$(a_i-a_j)/(a_i+a_j)$. As rational functions of $a_N$, the
Pfaffian expansion and the product have the same residues at every
$a_N=-a_i$, by induction on $N$. Both are bounded at infinity.
Their difference is therefore constant in $a_N$. At
$a_N=a_{N-1}$ both vanish, proving the induction; $N=2$ starts it.
The resulting rational identity includes all removable coincidences.
Substitution of the ghost variable proves the odd formula.
\end{proof}

\begin{proposition}[Joint normal convergence]\label{ana:normal}
On each of the two domains \(\DH,\DL\) specified by the independent
spectral geometry, the physical germ has a single-valued jointly
holomorphic continuation. It is given by
\begin{align}
 \XH&=(1-p^2)^{1/4}
       \sum_{\substack{N\ge1\\N\ {\rm odd}}}p^{(N-1)/2}F_N,
       \label{ana:globalHT}\\
 \XL&=(1-p^{-2})^{1/4}
       \sum_{\substack{N\ge2\\N\ {\rm even}}}p^{N/2}F_N,
       \label{ana:globalLT}
\end{align}
where \(F_N\) is the integral \eqref{ana:doublecorr} with its
\(p^{N/2}\) factor removed and with the multiplier
\eqref{ana:fullsites}. The sums converge normally, with every fixed
mixed complex derivative, on compact subsets. The fourth roots are
the binomial branches supplied by \(|p|<1\) on \(\DH\) and
\(|p|>1\) on \(\DL\).
\end{proposition}
\begin{proof}
On a compact parameter neighborhood in \(\C^2\setminus\Sigma\),
the polynomial \(\mathcal D\) has a nonzero minimum modulus on the
product of unit circles. By uniform continuity that gap persists on
a common thin product annulus and a slightly larger parameter
polydisk. Choose fixed radii \(r_x,r_y<1\) in the annulus.
All Schur entries on these circles are bounded by a constant
independent of \(N\). For even \(N\), the Schur Pfaffian identity
and Hadamard's determinant inequality give
\[
 \left|\prod_{i<j}a(\zeta_i,\zeta_j)\right|
 =|\operatorname{Pf}(a(\zeta_i,\zeta_j))|
 \le C^{N/2}N^{N/4}.
\]
For odd \(N\), adjoin a ghost variable \(-1\); its entries are
the removable constant one and the dimension becomes \(N+1\).
Apply the estimate separately to the two contour-variable families.
The one-body factors and contour lengths cost \(C^N\).
Also \(1-X\) and \(1-Y\) are bounded away from zero uniformly in
\(N\), whereas \(X^{-1},Y^{-1}\) cost only \(C^N\).
Thus, after increasing constants,
\begin{equation}\label{ana:compactbound}
 |F_N|\le C_0 C_1^N\frac{N^{(N+1)/2}}{N!}.
\end{equation}
The right side is summable after multiplication by any fixed
exponential in \(N\). Multivariate Cauchy estimates on the larger
polydisk give the same assertion for every fixed mixed derivative.
On overlaps, radii sufficiently close to one can be deformed into
one another inside the common gap annulus; all pair and product
denominators stay nonzero. The local integrals therefore glue.
All powers of \(p\) within \eqref{ana:globalHT}--\eqref{ana:globalLT}
are integers. The binomial fourth roots introduce no monodromy,
and the series prove joint holomorphy, including the HT axes
\(S=0\) or \(T=0\).

On a real physical compact set the same contours in
\eqref{ana:doublecorr} additionally give
\(r_x^{N|m|}r_y^{N|l|}\). Their double geometric sum is bounded
uniformly in \(N\ge1\). Thus the lattice and particle sums are
absolutely interchangeable, and \eqref{ana:fullsites} gives exactly
\eqref{ana:physicalHT}--\eqref{ana:physicalLT}, with the stated
zero-particle subtraction. Equality on a real open physical set
determines the joint germ: iterating the one-variable identity
theorem over a real rectangle proves the required uniqueness.
\end{proof}

The canonical root is that of Lemma~\ref{geo:canonical-root}.
We use $\gamma=-\log\rho$ only on local patches where $\rho\ne0$;
no logarithm through a root zero is required.

The angular representation used below is identified with this global
continuation coefficient by coefficient. Indeed
\(1/\sinh\gamma=2\rho/(1-\rho^2)\) and the rational expression
for \(P_{ij}\) in \eqref{ana:pair} contain no logarithm. After
including the physical normalization, their one-body prefactors are
\[
 (1-p^2)^{1/4}\frac{p^{(N-1)/2}}{S^N}
 \quad(N\text{ odd}),\qquad
 (1-p^{-2})^{1/4}\frac{p^{N/2}}{S^N}
 \quad(N\text{ even}).
\]
All powers here are integers. Hence the normalized angular coefficient
is single valued and holomorphic wherever \(S\ne0\) in the
corresponding physical domain; the real angular torus has uniform
root and denominator gaps on compact subsets. The set
\(\DH\setminus\{S=0\}\) is connected: radial contraction enters
a small polydisk without crossing \(z=0\), and that polydisk minus
\(z=0\) is connected. The domain \(\DL\) already has \(S\ne0\).
The identity theorem from the real physical rectangle therefore gives
equality with the double-contour coefficients throughout these sets.
Near a regular edge \(S,T\ne0\), this also fixes the compatible
local odd square root \(\sqrt B=\sqrt p/S\). Normal convergence
of the identified series follows from Proposition~\ref{ana:normal}.
Thus every boundary calculation below is a calculation of the global
physical sum on its actual interior side.

\begin{remark}[A restriction that really is removable]\label{ana:axis}
On \(K_y=0\), the lattice splits into independent horizontal chains.
For \(t=\tanh K_x\) near zero its response is
\[
 1+2\sum_{m\ge1}t^m=\frac{1+t}{1-t}=e^{2K_x}.
\]
This agrees with \eqref{ana:globalHT}: \(p=0\), only \(N=1\)
survives, and \(F_1=(A-S)^{-1}=e^{2K_x}\).
The restriction is entire, although the joint domain has boundary
points on this line. This exact example is why a bare joint
nonextension assertion cannot be used to prove the slice part
of the main theorem.
\end{remark}

\subsection{Energy-field normalization}

For fixed real \(J_x,J_y>0\) and physical temperature
\((K_x,K_y)=(\beta J_x,\beta J_y)\), a field of energy units enters
the Hamiltonian as \(-h\sum_v\sigma_v\), and \(\eta=\beta h\).
Finite-volume differentiation gives \(\chi=\beta\mathcal X\).
The same statement holds after taking the prescribed infinite-volume
phase limit. To justify the latter without exchanging an uncontrolled
field limit and infinite sum, use the GHS inequality for arbitrary
nonnegative couplings and fields \cite[Theorem 1.5]{CGN}.
Pinning removed spins by fields tending to \(+\infty\) shows that
plus-box covariances increase with the box. Their nonnegative sums
increase to the finite sum proved above. GHS concavity gives, for
the plus-box magnetization,
\[
 m_V(h)-m_V(0)\le h\,m_V'(0)
       \le\beta h\,\mathcal X,\qquad h\ge0.
\]
The plus-state limit yields \(m'(0+)\le\beta\mathcal X\).
Conversely turn on \(h\) only on a finite set \(E\). The resulting
bounded local change of measure has derivative
\(\beta\sum_{v\in E}\operatorname{Cov}_{+,0}(\sigma_0,\sigma_v)\).
Ferromagnetic monotonicity bounds its magnetization above by the
uniform-field magnetization, with equality at \(h=0\).
Increasing \(E\) proves the reverse inequality. Thus
\(m'(0+)=\beta\mathcal X\). Spin-flip symmetry gives the two-sided
HT derivative; in LT the derivative is the plus-phase right
derivative (or its minus-phase counterpart).
This statement is made on temperature lines only.

\section{Regular spectral edges and transverse curves}
\label{ana:edges}

The following estimates use independent complex dispersion parameters,
not uniformity inferred from a collection of real ratios.
On a regular edge, exchange the two lattice axes if necessary.
In the local coordinates \(C=A/S,\ B=T/S\) of the geometry section,
one edge is \(C-B=b\), \(b=\cos\phi\in(-1,1)\).
Choose the side for which \(\Im B_0>0\); conjugation handles the
other side. Let \(B(u)\) be any holomorphic function near zero with
\(B(0)=B_0\), and set
\begin{equation}\label{ana:edgecurve}
 W(q;u)=b+u+B(u)(1-\cos q).
\end{equation}
Both a constant-\(B\) complex transversal in \(\C^2\) and each
nonramified physical temperature-line transversal have this form.

At \(u=0\), the inside root lies in the closed lower half disk,
with \(\rho(0)=e^{-\ii\phi}\). Compactness gives
\[
 \Im\rho(q)\le-\delta<0,\qquad |\rho(q)|\le1.
\]
There is no \(W=\pm1\) on the real angular circle. The roots,
their chosen logarithms, the one-body factors, and the canceled
pairs \eqref{ana:pair} are analytic on one fixed complex
neighborhood in \((u,B,q)\). On \(u=\ii\eps\), \(\eps>0\),
\begin{equation}\label{ana:attenuation}
 \Re\gamma(q;\ii\eps)
   \ge c\{\eps+\operatorname{dist}(q,0)^2\}.
\end{equation}
Indeed
\(\Im W=\eps+\Im B(\ii\eps)(1-\cos q)\) has this lower bound;
the identity
\(\Im W=\sinh(\Re\gamma)\sin(\Im\gamma)\) on the compact branch
then gives \eqref{ana:attenuation}. This is the complement side
of the parallelogram edge and hence the side of the bordering
physical component.

\begin{lemma}[Strict contraction on a complex polydisk]
\label{ana:matching}
There are \(\kappa>0\), \(C<\infty\), and a fixed complex parameter
polydisk about this edge such that, for every fixed mixed derivative
\(\partial^\nu\) in the parameters,
\begin{equation}\label{ana:gaussian}
 |\partial^\nu G_N(q)|
 \le C_\nu\exp\{-\kappa N^2+C N\log(N+1)\}
\end{equation}
for all real angles and every \(N\ge1\).
\end{lemma}
\begin{proof}
The identity
\[
 |1-vv'|^2-|v-v'|^2
 =(1-|v|^2)(1-|v'|^2)+4\Im v\,\Im v'
\]
gives a strict bound \(|a(\rho_i,\rho_j)|\le q_0<1\).
The denominator is separated from zero. By continuity a slightly
larger constant below one still works on a common \emph{complex}
parameter polydisk; individual root moduli need not stay below one
there. Every canceled pair \(P_{ij}\) remains bounded by \(C\).

For even \(N\), expand the Schur Pfaffian in the
\(\eta_i=e^{\ii q_i}\) variables only. For each matching, combine
its \(N/2\) edges with the corresponding \(\rho\) factors using
\eqref{ana:pair}. The other \(N(N-2)/2\) root factors retain their
strict contraction. Thus
\[
 \left|\prod_{i<j}P_{ij}\right|
 \le(N-1)!!\,C^{N/2}q_0^{N(N-2)/2}.
\]
For odd \(N\), adjoining the constant ghost entry gives
\[
 \left|\prod_{i<j}P_{ij}\right|
 \le N!!\,C^{(N-1)/2}q_0^{(N-1)^2/2}.
\]
These identities first hold away from reciprocal \(\eta\) pairs and
extend by the analytic canceled formula. There is no division by
a root-Schur zero. One-body factors cost \(C^N\).
Cauchy's inequalities on the larger polydisk prove all the derivative
bounds. They are applied to \(G_N\), never across a pole of
\((1+Z)/(1-Z)\).
\end{proof}


\begin{lemma}[One fixed collision block for the entire mixed-derivative tail]
\label{ana:positivetail}
Fix $\phi=2\pi a/\ell\in(0,\pi)$, where $\ell$ is an odd
prime.  Let $\mathcal P\Subset\mathcal P^+$ be concentric complex
parameter polydisks.  Suppose $B(u,\eta)$ is holomorphic on
$\{|u|<R\}\times\mathcal P^+$, with
\[
 |B(u,\eta)|\le M,\qquad
 \Im B(0,\eta)\ge\beta>0
 \quad(\eta\in\mathcal P^+).
\]
Use $W(q;u,\eta)=\cos\phi+u+B(u,\eta)(1-\cos q)$ and the
common branch $\gamma(0,0,\eta)=\mathrm i\phi$ in the complete
coefficient $C_N(u,\eta)$ of \eqref{ana:angular}.  For odd parity
put $n=\ell$, $k=(n^2-1)/2$; for even parity put $n=2\ell$,
$k=n^2/2-1$.  There are $\epsilon_0,\kappa>0$ such that, for
every multi-index $\alpha$ and every $0\le j\le k$,
\begin{equation}\label{ana:block-mixed-tail}
 \begin{gathered}
 \sup_{\eta\in\mathcal P}
 |\partial_u^j\partial_\eta^\alpha C_N(\mathrm i\epsilon,\eta)|
 \le \frac{C_{n,j,\alpha}N^j}{N!}
       \exp\{-\kappa N^2+C_nN\log(N+1)\},\\
 \quad 0<\epsilon<\epsilon_0,
 \quad N>n,\quad N\equiv n\pmod2.
 \end{gathered}
\end{equation}
In particular the entire higher-index tail and all these mixed
derivatives are bounded and have locally uniform boundary limits as
$\epsilon\downarrow0$.  Every lower index of the chosen parity is
jointly holomorphic on a neighborhood of
$\{0\}\times\overline{\mathcal P}$.

The constants can be chosen uniformly for all $B$ with the displayed
common domain and bounds.  They may depend on the fixed phase, $R,M$,
$\beta$, the two parameter polydisks, and the indicated derivative
orders.  Taking $B(u,\eta)=B_0+\eta$ includes derivatives in the
independent dispersion variables $(u,B)$.
\end{lemma}

\begin{proof}
Cauchy's inequality in $u$, followed by a uniform shrinking of its
disk, gives $\Im B(u,\eta)\ge\beta/2$ for $|u|<r$ and
$\eta\in\mathcal P^+$.  At $u=0$, the corresponding compact
set of roots is contained in
\[
 \{\rho:|\rho|\le1,\ \Im\rho\le-\delta\}
\]
for some $\delta>0$.  Indeed $B$ ranges in a bounded closed subset
of the upper half-plane, $q$ ranges over a compact circle, and
$W=\pm1$ never occurs.  The chosen lower-half-plane root, its
logarithm, and all one-body and canceled-pair factors consequently
extend to a common complex $u$ disk, uniformly in $\eta$.
On a slightly smaller such disk the strict root-Schur estimate in
Lemma~\ref{ana:matching} holds with one constant $\theta<1$.
Individual root moduli need not remain at most one on this complex
disk; only that strict Schur estimate is used there.

Fix an integer $m\ge1$ and define the parameter-independent factor
\[
 V_m(q)=\prod_{1\le a<b\le m}
                  \sin^2\bigl((q_a-q_b)/2\bigr).
\]
For $N\ge m$, factor $G_N=V_m A_{N,m}$ algebraically.  Namely,
replace each pair internal to the first $m$ labels by
$B/\sinh^2((\gamma_a+\gamma_b)/2)$, leaving the other pairs and
the one-body factors as they stand.  This defines $A_{N,m}$
holomorphically even when $V_m=0$ and does not divide an estimate
by a possibly vanishing pair.
The internal factors are bounded by a constant depending on $m$;
the $m(N-m)$ cross pairs cost at most $C^{m(N-m)}$.
Apply the matched-Pfaffian estimate only to the remaining $N-m$
labels, using its even or ghost-augmented odd version as appropriate.
The empty product is one.  It follows that
\[
 |A_{N,m}|
 \le C_m\exp\{-\kappa(N-m)^2
       +C(N-m)\log(N-m+1)+C_mN\}.
\]
Cauchy's inequality on the common larger $u$ disk therefore gives,
for each fixed $r\ge0$,
\begin{equation}\label{ana:block-numerator}
 |\partial_u^rG_N(q;u,\eta)|
 \le C_{m,r}V_m(q)
       \exp\{-\kappa N^2+C_mN\log(N+1)\},
\end{equation}
uniformly for real angles and $\eta\in\mathcal P^+$.
The Cauchy estimate here is applied to the analytic numerator factor,
not to the product kernel.

On $u=\mathrm i\epsilon$ let
\[
 Q_0(q)=\sum_{a=1}^N\operatorname{dist}(q_a,2\pi\mathbb Z)^2.
\]
The common positive imaginary margin gives
$\Re\gamma(q;\mathrm i\epsilon,\eta)
\ge c(\epsilon+\operatorname{dist}(q,2\pi\mathbb Z)^2)$.
Hence, with $Z=\exp(-\sum_a\gamma_a)$,
\[
 |1-Z|\ge1-|Z|\ge c\min\{1,N\epsilon+Q_0\}.
\]
All fixed $u$ derivatives of each $\gamma_a$ are uniformly bounded.
The chain rule for the exponential and the rational kernel thus yields
\begin{equation}\label{ana:block-kernel}
 \left|\partial_u^r\frac{1+Z}{1-Z}\right|
 \le C_jN^j(1+Q_0^{-j-1}),\qquad 0\le r\le j,
\end{equation}
uniformly in $\eta\in\mathcal P^+$ and $0<\epsilon<\epsilon_0$.
For example, derivatives of $\coth(\Gamma/2)$ of order $s$ are
bounded by $C_s(1+(\Re\Gamma)^{-s-1})$ on $\Re\Gamma>0$;
each derivative of $\Gamma=\sum_a\gamma_a$ costs at most $C_jN$.

Now take $m=n+1$.  Every larger same-parity index satisfies
$N\ge m+1$, so $q_1,\ldots,q_m$ are independent variables of
the constrained torus integral.  Choose these representatives in
$[-\pi,\pi]$ and set $R_m^2=\sum_{a\le m}q_a^2$.
Then $Q_0\ge R_m^2$ and
$V_m\le C_mR_m^{m(m-1)}$.  The fixed-dimensional integral
\[
 \int_{[-\pi,\pi]^m} V_m(q)(1+R_m^{-2j-2})
                        \prod_{a=1}^m\frac{\mathrm dq_a}{2\pi}
\]
is finite: near zero its radial power is
$R_m^{m^2-2j-3}\,\mathrm dR_m$, and
$m^2>2j+2$ for every $j\le k$ in either parity.
Every unused normalized angular measure has mass one.
Applying the product rule to \eqref{ana:angular} and using
\eqref{ana:block-numerator}--\eqref{ana:block-kernel} proves
\eqref{ana:block-mixed-tail} with $\alpha=0$, uniformly on the
larger parameter polydisk.

For fixed $\epsilon>0$, $C_N$ and its $u$ derivatives are
holomorphic in $\eta$ throughout $\mathcal P^+$: the same
positive imaginary margin keeps the whole real torus on the physical
side for every complex $\eta$ there.  Apply the multivariate Cauchy
formula on parameter circles of a fixed radius between
$\mathcal P$ and $\mathcal P^+$.  It adds only
$\alpha!\prod_a d_a^{-\alpha_a}$ to the preceding bound and
proves \eqref{ana:block-mixed-tail}.  In particular tangential
derivative order does not spend additional collision powers.  This
argument uses an actual common complex polydisk; estimates restricted
to real parameters would not suffice.

The majorant is summable in $N$.  For each fixed higher index its
integrand and all required derivatives have a pointwise boundary
limit away from the all-zero collision, and the displayed
fixed-dimensional majorant applies uniformly.  Dominated convergence,
first for each index and then for their summable series, proves the
locally uniform boundary-limit assertion; the same parameter Cauchy
formula gives every mixed limit.  For a lower index, $|Z|=1$ at
$u=0$ is possible only when every angle is zero, where
$Z=e^{-\mathrm iN\phi}\ne1$.  Compactness of that fixed torus
and the closed parameter polydisk gives a common denominator gap,
which persists on a neighborhood of $u=0$.  Thus every lower
coefficient is jointly holomorphic there.
\end{proof}


\begin{proposition}[Nonzero full-sum edge obstruction on every
holomorphic transversal]\label{ana:positive}
Under \eqref{ana:edgecurve} and the rational-phase hypotheses of
Lemma~\ref{ana:positivetail}, the complete sum of the selected
parity has no holomorphic extension through \(u=0\).
The assertion is locally uniform in independent complex \(B\)
parameters. It remains valid after multiplication by a
holomorphic nonvanishing prefactor and after a nonramified change
of transversal coordinate. It holds for both endpoint signs in
the even case.
\end{proposition}
\begin{proof}
Put \(\alpha=(n^2-1)/2\) and
\[
 P_n=\prod_{r=1}^{n-1}\frac{r!}{2^r}.
\]
The Gaussian Vandermonde period on the \emph{independent-coordinate}
measure is
\begin{equation}\label{ana:mehta}
 \int_{q_n=-\sum_{i<n}q_i}e^{-Q}\Delta(q)^2
               \dd q_1\cdots\dd q_{n-1}
 =\frac{n!\,\pi^{(n-1)/2}P_n}{\sqrt n}=:\mathcal G_n.
\end{equation}
To derive it, expand the ordinary \(n\)-variable Vandermonde
determinant in monic Hermite polynomials. Their Gaussian norms are
\(\sqrt\pi\,r!/2^r\), so that integral is
\(n!\pi^{n/2}P_n\). Setting \(x_i=q_i+m\), \(\sum q_i=0\),
has Jacobian \(n\), and the \(m\) integral is \(\sqrt{\pi/n}\).
Division proves \eqref{ana:mehta}.

On the unique collision chart,
\[
 \gamma(q;u)=\ii\phi-\frac{\ii}{\sin\phi}
       \left(u+\frac{B_0q^2}{2}\right)
       +O(u^2+u q^2+q^4).
\]
The leading singular integral for \(C_n\) is consequently
\begin{equation}\label{ana:firstintegral}
 \frac{2\ii^{1-n^2}B_0^{n^2/2}}
 {n!(2\pi)^{n-1}2^{n(n-1)}\sin^{n^2-1}\phi}
 \int\frac{\Delta(q)^2}{nu+(B_0/2)Q}\,
             \dd q_1\cdots\dd q_{n-1}.
\end{equation}
For odd \(n\) the local square root of \(B\) is the one used in
the physical angular coefficient. Its sign can change the displayed
amplitude's sign but cannot make it zero. For even \(n\) the
power in the original integrand is an integer; the square root
below is fixed by the Gaussian period.

After \(k\) derivatives, the model integral is evaluated from
\[
 \int\frac{\Delta(q)^2}{(a+bQ)^{k+1}}\dd q
 =\mathcal G_n\,b^{-\alpha}
    \frac{\Gamma(k+1-\alpha)}{\Gamma(k+1)}
    a^{\alpha-k-1},
\]
initially in a common right half-plane and then by analytic
continuation to \(a=nu,\ b=B_0/2\).
More explicitly, rotate both coefficients by \(-\ii\):
\(a'=n\eps>0\) and \(b'=-\ii B_0/2\) has positive real part.
The absolutely convergent Laplace--Gaussian calculation for
\(a',b'\) fixes the continuation and its square root. No assumption
\(\Re B_0>0\) is made.
For odd \(n\), \(k=\alpha\) is a multiple of four; for even
\(n\), \(k=\alpha-\tfrac12\) is odd. Substitution gives
\begin{align}
 \lim_{\eps\downarrow0}u\,\partial_u^k C_n(u)
 &=\frac{2^{1-k}\sqrt{B_0}\,P_n n^{k-3/2}}
         {\pi^{(n-1)/2}\sin^{2k}\phi}\ne0,
 &&u=\ii\eps,\quad n\ {\rm odd},\label{ana:oddamp}\\
 \lim_{\eps\downarrow0}u^{1/2}\partial_u^k C_n(u)
 &=-\ii\,\frac{2^{1-\alpha}\sqrt{B_0}\,P_n n^{k-1}}
         {\pi^{(n-2)/2}\sin^{n^2-1}\phi}\ne0,
 &&u=\ii\eps,\quad n\ {\rm even}.\label{ana:evenamp}
\end{align}
In the second line \(u^{1/2}=e^{\ii\pi/4}\sqrt\eps\).

For completeness these are limits of the actual fixed-\(n\)
integral. On its collision chart write the exact integrand as
\(\Delta^2 f(q,u)/H(q,u)\). It is even under simultaneous
\(q\mapsto-q\), and, with \(c=-\ii/\sin\phi\),
\[
 H=c\{nu+(B_0/2)Q\}+O((|u|+Q)^2),\quad
 H_u=cn+O(|u|+Q),\quad f=f_0+O(|u|+Q).
\]
Both the exact and model denominators have modulus at least
\(c_1(\eps+Q)\) on \(u=\ii\eps\). The top derivative has
leading term \((-1)^kk! fH_u^k/H^{k+1}\); all other terms have
denominator power at most \(k\). Subtracting the constant quadratic
model therefore gives the pointwise bound
\[
 |\partial_u^k(\hbox{exact integrand}-\hbox{model integrand})|
       \le C\Delta^2(\eps+Q)^{-k}.
\]
Its radial integral is \(O(\log(1/\eps))\) for odd \(n\), where
\(k=\alpha\), and \(O(1)\) for even \(n\), where \(k=\alpha-1/2\).
Extending the differentiated model from the chart to the whole
tangent plane costs \(O(1)\), since its radial tails are
\(R^{-3}\) and \(R^{-2}\), respectively. Outside the chart the
exact denominator has a gap. These bounds prove the two claimed
limits, including for the actual holomorphic function \(B(u)\),
whose contribution to the denominator remainder is \(O(uQ)\).

Lemma~\ref{ana:positivetail} bounds all higher derivatives in the
entire tail, and lower indices are analytic. Lower derivatives of
the even first term are bounded; the next lower odd derivative
may be logarithmic, but all are \(o(u^{-1})\).
Thus the nonzero top amplitudes cannot be canceled in the full
sum or by product-rule terms from an analytic nonvanishing
physical prefactor. The chain rule proves the statement for every
nonramified curve parameter.

At the negative endpoint replace \(q_i\) by \(\pi+q_i'\).
For even \(N\) the total momentum constraint is unchanged and
\(B^{N^2/2}=(-B)^{N^2/2}\). Apply the preceding proof with
\(-B\), conjugating if necessary to have positive imaginary part.
Exchanging transfer axes handles the other two edges.
For odd \(N\) this shift instead makes the total reduced angle
an odd multiple of \(\pi\); that distinct case is addressed by
the full boundary-jet theorem in the next section.
\end{proof}

\paragraph{Uniformity of the first amplitude in holomorphic families.}
Under the common-family hypotheses of Lemma~\ref{ana:positivetail},
all branch neighborhoods, scalar derivative bounds, numerator
factorizations, and quadratic-denominator bounds in the fixed-$n$
calculation can be chosen on $\mathcal P^+$ uniformly.  Its actual
coefficient $B(u,\eta)$ contributes
$[B(u,\eta)-B(0,\eta)]Q=O(|u|Q)$ to the denominator remainder.
The comparison errors are therefore uniformly
$O(1+|\log\epsilon|)$ in the odd case and $O(1)$ in the even
case.  The nonzero limits \eqref{ana:oddamp}--\eqref{ana:evenamp}
hold locally uniformly in $\eta$, with $B_0$ replaced by
$B(0,\eta)$ and its compatible common square root.  Multivariate
Cauchy on $\mathcal P^+$ then gives these limits after every fixed
$\eta$ derivative.  Their differentiated leading coefficients may
of course vanish; the undifferentiated leading coefficient is nonzero
for every member of the family.  The same argument applies to
lower-derivative bounds and to an analytic nonvanishing physical
prefactor.  Thus the full-sum obstruction and its mixed estimates
hold uniformly on each compact admissible holomorphic family.

At a negative endpoint of the even sum, the change
$q_a=\pi+q'_a$ preserves the momentum constraint and replaces $B$
by $-B$; the equality $B^{N^2/2}=(-B)^{N^2/2}$ holds for even
$N$.  The preceding family argument therefore applies after this
change and, when necessary, complex conjugation of all data.  The
conjugated family is understood holomorphically as
$\overline{B(\bar u,\bar\eta)}$ on the reflected polydisk.


\section{The full odd sum at a regular negative spectral edge}
\label{neg:section}

This section proves a local theorem for an arbitrary holomorphic
anisotropy along a transverse analytic curve.  In particular, the
anisotropy is not held constant.  All estimates retain the exact
particle-dependent anisotropy factor, the complete constrained angular
integral, and the sum over every odd particle number.

\subsection{Statement, branch, and the physical side}
\label{neg:setup}

Fix
\begin{equation}
 \phi=\frac{2\pi r}{p}\in(0,\pi),\qquad
 (r,p)=1,\qquad p\ge3\ \hbox{odd},\qquad
 b=\cos\phi,\qquad s_\phi=\sin\phi>0 .
 \label{neg:rational}
\end{equation}
Let \(B\) be holomorphic on a disk about zero and suppose that
\begin{equation}
 \Im B(0)<0 .
 \label{neg:orientation}
\end{equation}
Shrink this disk so that \(B\ne0\) and \(\Im B<0\) throughout it, and
fix one holomorphic logarithm of \(B\).  For \(q=\pi+\delta\), define
\begin{equation}
 W(t,\pi+\delta)
   =b+B(t)\bigl(t-1+\cos\delta\bigr),\qquad
 g(t,\delta)=\operatorname{arcosh}W(t,\pi+\delta),\qquad
 g(0,0)=\ii\phi .
 \label{neg:dispersion}
\end{equation}
At \(t=0\), \(W\) is in the open upper half-plane except at
\(\delta=0\bmod 2\pi\), where it equals \(b\in(-1,1)\).  Consequently
the branch in \eqref{neg:dispersion} extends to a fixed complex
neighborhood of all real angles and a fixed parameter disk.  It is
even and \(2\pi\)-periodic in \(\delta\).  Put
\(\gamma(t,q)=g(t,q-\pi)\) and \(\rho(t,q)=e^{-\gamma(t,q)}\).
At \(t=0\), \(|\rho|\le1\), and the compact set of these roots lies
strictly in the lower half-plane:
\begin{equation}
 \Im\rho(0,q)\le-\delta_0<0 .
 \label{neg:root-compact}
\end{equation}
The logarithm used for \(\gamma\) is fixed by
\(0<\Im\gamma(0,q)<\pi\).  In particular all the one-body
denominators and the canceled pair denominators below are nonzero on
a fixed complex neighborhood.

For odd \(n\ge1\), let
\[
 \mathcal T_n=\{(q_1,\ldots,q_n)\in(\R/2\pi\Z)^n:
                         q_1+\cdots+q_n=0\bmod 2\pi\},
 \qquad q_n=-\sum_{i<n}q_i\bmod 2\pi .
 \]
The measure on this torus is \(\dd q_1\cdots\dd q_{n-1}\).  Include
the normalization in the numerator:
\begin{align}
 H_n(t,q)
 &=\frac{B(t)^{n^2/2}}{n!(2\pi)^{n-1}}
   \prod_{i=1}^n\frac1{\sinh\gamma(t,q_i)}
   \prod_{i<k}
   \left[
    \frac{\sin((q_i-q_k)/2)}
         {\sinh((\gamma(t,q_i)+\gamma(t,q_k))/2)}
   \right]^2 ,
 \label{neg:numerator}\\
 \Gamma_n(t,q)&=\sum_{i=1}^n\gamma(t,q_i),\\
 C_n(t)&=\int_{\mathcal T_n}H_n(t,q)
                          \coth\!\left(\frac{\Gamma_n(t,q)}2\right)
                          \dd q .
 \label{neg:coefficient}
\end{align}
The power of \(B\) uses the chosen common logarithm.  Changing its
square-root sign changes every odd coefficient by the same sign.

The relevant open side of the boundary is
\begin{equation}
 \mathcal U^-=\{\,|t|<r_0:\ \Im(B(t)t)>0\,\}.
 \label{neg:physical-side}
\end{equation}
Here
\[
 \Im W(t,\pi+\delta)
   =\Im(B(t)t)-\Im B(t)(1-\cos\delta)>0
 \]
for every real angle, so the chosen roots satisfy \(|\rho|<1\).
This is the local physical side when the regular negative endpoint
has the lower-half-plane root \(\rho(0,\pi)=e^{-\ii\phi}\).
It contains \(t=-\epsilon\) for every sufficiently small
\(\epsilon>0\).  When these coordinates are obtained from
\(W(q)=C-B\cos q\), they are exactly
\[
 t=\frac{C+B-b}{B}.
 \]
The curve hypothesis is that this \(t\) is a nonramified local
parameter.  Thus any transverse analytic curve through the selected
endpoint is included, with its actual function \(B(t)\).  The opposite
orientation is obtained by complex conjugating all the data.

\begin{theorem}[Universal analytic-curve obstruction at a negative edge]
\label{neg:theorem}
Under \eqref{neg:rational}--\eqref{neg:orientation}, the series
\[
 F(t)=\sum_{\substack{n\ge1\\n\ {\rm odd}}}C_n(t)
 \]
converges normally on \(\mathcal U^-\).  Every inward derivative has
a finite limit
\begin{equation}
 a_j=\lim_{\epsilon\downarrow0}\frac{F^{(j)}(-\epsilon)}{j!},
 \qquad j=0,1,2,\ldots .
 \label{neg:jets}
\end{equation}
Nevertheless its formal boundary Taylor series has radius zero.
More precisely, as \(n\) tends to infinity through odd multiples of
\(p\), there are integers \(j_n\) such that
\begin{align}
 j_n&\sim n^2\log n,
 \label{neg:jn-statement}\\
 |a_{j_n}|^{1/j_n}
   &\sim \frac{4j_n}{e\pi^2\log j_n}.
 \label{neg:growth-statement}
\end{align}
It follows that \(F\) has no holomorphic extension through \(t=0\)
from \(\mathcal U^-\).  If \(P\) is holomorphic near zero and
\(P(0)\ne0\), the same nonextension conclusion holds for \(P F\).

The hypotheses and all estimates are locally uniform for complex
holomorphic families \(B(t,\eta)\): it suffices to work on a common
parameter disk and a relatively compact parameter polydisk on which
\(\Im B(0,\eta)\le-\beta<0\) and \(B\) has a common holomorphic bound.
The constants can depend on that disk, that bound, \(\beta\), and the
 fixed rational phase.  The selected integers are allowed to depend
 on the family parameter, \(j_n=j_n(\eta)\); a common sequence of
 selected integers for all \(\eta\) is not asserted.
\end{theorem}

We prove the theorem by identifying a single particle index that
dominates selected, increasingly high Taylor coefficients in modulus.
The other indices have arbitrary complex phases; no positivity or
phase alignment of the infinite sum is assumed.

\subsection{A uniform Gaussian numerator bound and all boundary jets}
\label{neg:jets-section}

Write \(\mathcal A(v,w)=(v-w)/(1-vw)\) and \(y_i=e^{\ii q_i}\).
Subtraction of two dispersion equations gives the canceled identity
\begin{equation}
 \mathcal A(\rho_i,\rho_k)\mathcal A(y_i,y_k)
 =B(t)\left[
    \frac{\sin((q_i-q_k)/2)}
         {\sinh((\gamma_i+\gamma_k)/2)}
       \right]^2 .
 \label{neg:canceled}
\end{equation}
This is an identity of holomorphic canceled functions, including
reciprocal \(y\)-pairs.  The elementary identity
\begin{equation}
 |1-vw|^2-|v-w|^2
 =(1-|v|^2)(1-|w|^2)+4\Im v\Im w
 \label{neg:schur}
\end{equation}
and \eqref{neg:root-compact} imply
\(|\mathcal A(\rho_i,\rho_k)|\le\theta<1\) at \(t=0\).
The same strict bound, with a slightly larger \(\theta<1\), holds
on a fixed complex parameter disk for all real angles.  Every
complete canceled pair in \eqref{neg:canceled} is bounded there.

For completeness, the Schur identity used here is
\(\operatorname{Pf}(\mathcal A(y_i,y_k))_{i,k=1}^{N}
=\prod_{i<k}\mathcal A(y_i,y_k)\) for even \(N\).
It follows by induction in \(N\): as rational functions of \(y_N\),
the two sides have the same simple residues at \(y_N=1/y_k\).
In checking the residues, the other paired factors cancel by
\(\mathcal A(a,b)\mathcal A(a,1/b)=1\), and the remaining product
is the induction hypothesis.  Their difference is constant because
both are bounded at infinity, and that constant is zero at
\(y_N=y_{N-1}\).  The identity then extends from generic tuples
by rational continuation.

Apply this identity only to the \(y\)-variables, adjoining the
ghost variable \(-1\) because \(n\) is odd.  The entry
\(\mathcal A(y_i,-1)\) is the removable constant one.
Each of the \(n!!\) matchings has \((n-1)/2\) real edges.
Combine each matched \(y\)-factor with the corresponding
\(\rho\)-factor by \eqref{neg:canceled}.  The remaining
\((n-1)^2/2\) \(\rho\)-factors are at most \(\theta\) in modulus.
Thus, with the one-body factors and normalization included,
\begin{equation}
 \sup_{|t|<r_0,\ q\in\mathcal T_n}|H_n(t,q)|
 \le E_n,\qquad
 E_n=\exp\{-\kappa n^2+C n\log(n+1)+C\}
 \label{neg:gaussian}
\end{equation}
for some \(\kappa>0\).  Initially the Pfaffian calculation can be
made off its apparent reciprocal-pair divisors; the canceled
identity extends the estimate to those divisors.  There is no
division by a possibly zero \(\rho\)-factor.  

Choose shortest lifts \(q_i=\pi+\delta_i\bmod2\pi\), with
\(\delta_i\in[-\pi,\pi]\), and put
\[
 S(q)=\sum_{i=1}^n\delta_i^2 .
 \]
The odd momentum constraint says
\(\sum_i\delta_i=w\pi\), where \(w\) is odd.  Therefore
\begin{equation}
 S(q)\ge\frac{\pi^2}{n}.
 \label{neg:momentum-gap}
\end{equation}
On the inward ray, compactness of the branch and
\(\Im W(-\epsilon,\pi+\delta)\ge
c(\epsilon+\operatorname{dist}(\delta,2\pi\Z)^2)\) imply
\begin{equation}
 \Re g(-\epsilon,\delta)
 \ge c\bigl(\epsilon+\operatorname{dist}(\delta,2\pi\Z)^2\bigr).
 \label{neg:attenuation}
\end{equation}
Indeed \(\Im W=\sinh(\Re g)\sin(\Im g)\), and the relevant roots
remain in a fixed compact set.  All parameter derivatives of \(g\)
are bounded on a fixed complex neighborhood.

It follows from \eqref{neg:momentum-gap} that \(C_n\) is holomorphic
on \(|t|<c/n^2\).  On this disk,
\(\Re\Gamma_n\ge c/n\) after decreasing \(c\), and hence
\(|\coth(\Gamma_n/2)|\le Cn\).  The torus volume is absorbed in
the linear exponential in \eqref{neg:gaussian}.  Cauchy's estimate
therefore gives
\begin{equation}
 |[t^j]C_n(t)|\le (Cn^2)^j\,CnE_n .
 \label{neg:small-disk}
\end{equation}
More generally, the disk of radius \(c(\epsilon+n^{-2})\) about
\(-\epsilon\) has
\(\Re\Gamma_n\ge c(n\epsilon+n^{-1})\).  The same argument bounds
the \(j\)th derivative at \(-\epsilon\) by
\(j!(Cn^2)^j CnE_n\), uniformly for small \(\epsilon\).
On compact subsets of \(\mathcal U^-\), the kernel is bounded and
\eqref{neg:gaussian} proves normal convergence.  Dominated
convergence for each fixed \(j\) now proves
\begin{equation}
 a_j=\sum_{\substack{n\ge1\\n\ {\rm odd}}}[t^j]C_n(t),
 \qquad
 \sum_n |[t^j]C_n(t)|<\infty .
 \label{neg:jet-sum}
\end{equation}
This establishes the existence of all the jets, but does not give
a positive radius for their formal Taylor series.

\subsection{An exact reduction to simple-pole periods}
\label{neg:pole-reduction}

Fix a sufficiently small \(\sigma>0\).  There is a fixed \(c_1>0\)
such that on \(|t|<c_1/n\) the part with \(S\ge\sigma\) is analytic
with a bounded kernel.  To see this, use
\(\Re\Gamma_n(0,q)\ge cS\) and the perturbation bound
\(|\Gamma_n(t,q)-\Gamma_n(0,q)|\le Cn|t|\).
Choose \(c_1\) to be a sufficiently small multiple of \(\sigma\).

In the region \(S<\sigma\), set
\[
 D_n(t,q)=\Gamma_n(t,q)-n\ii\phi .
 \]
The Taylor expansion of \(g\) at \((0,0)\) gives
\(|D_n(t,q)|\le C(S+n|t|)\).
If \(p\nmid n\), the distance of \(n\phi\) from \(2\pi\Z\) is at
least \(2\pi/p\).  Reduce \(\sigma,c_1\) so that the perturbation
is less than \(\pi/p\).  This part is also analytic, with bounded
kernel, on \(|t|<c_1/n\).

It remains to consider odd multiples of \(p\) and \(S<\sigma\).
Here \(n\phi\in2\pi\Z\), and
\[
 \coth(\Gamma_n/2)=\frac2{D_n}+\mathcal R(D_n),
 \]
where \(\mathcal R\) is holomorphic on a fixed disk containing all
these values for \(|t|<c_1/n\).
There is a unique small zero \(t=T_n(q)\) of \(D_n(t,q)\).
This statement has a fixed-disk quantitative form.  With
\(a_0=B(0)/(\ii s_\phi)\ne0\), one has
\[
 \frac1n\partial_tD_n(t,q)=a_0+O(|t|+\sigma)
 \]
on a fixed disk.  By reducing that disk and \(\sigma\), its
distance from \(a_0\) is less than \(|a_0|/4\).  Rouch\'e's
theorem on a circle of radius \(CS/n\), followed by uniqueness
from this derivative estimate, gives
\begin{equation}
 |T_n(q)|\le CS/n,\qquad
 D_n(t,q)=n(t-T_n(q))A_n(t,q),
 \label{neg:pole-factor}
\end{equation}
where \(A_n\) is holomorphic on a fixed smaller disk and bounded
away from zero.  One can take
\[
 A_n(t,q)=\frac1n\int_0^1
   \partial_tD_n(T_n(q)+u(t-T_n(q)),q)\,\dd u .
 \]
All such disks and bounds are independent of \(n\) and \(q\).

Put \(J_n=H_n/A_n\).  The exact identity
\begin{equation}
 \frac{2H_n(t,q)}{D_n(t,q)}
 =\frac{2J_n(T_n(q),q)}{n(t-T_n(q))}
  +\frac2n
     \frac{J_n(t,q)-J_n(T_n(q),q)}{t-T_n(q)}
 \label{neg:division}
\end{equation}
separates the pole from an analytic divided difference.
The second term is bounded by \(CE_n\), by fixed-disk Cauchy
estimates.  The \(\mathcal R\) term has the same bound on
\(|t|<c_1/n\).  Combining these observations with the exterior
and nonresonant regions gives
\begin{align}
 a_j&=
 \sum_{\substack{n\ {\rm odd}\\p\mid n}}A^{\rm pole}_{n,j}+R_j,
 \label{neg:all-poles}\\
 A^{\rm pole}_{n,j}
 &=-\int_{S<\sigma}
   \frac{2H_n(T_n(q),q)}
        {\partial_t\Gamma_n(T_n(q),q)}
   T_n(q)^{-j-1}\,\dd q ,
 \label{neg:pole-coefficient}\\
 |R_j|
 &\le \sum_{n\ge1}(Cn)^j
             \exp\{-\kappa n^2+C n\log(n+1)+C\}\\
 &\le
 \exp\{\tfrac12 j\log(j+2)+Cj+C\}.
 \label{neg:regular-remainder}
\end{align}
For each \(n\), the coefficient identity is first justified on
the smaller disk \(|t|<c/n^2\); no pole is crossed there.
The last inequality follows by absorbing \(Cn\log(n+1)\) into
half the Gaussian and maximizing \(j\log n-\kappa n^2/2\).
This reduction isolates all small pole locations on the original
real angular cycle.

\subsection{Bounds for every winding and every particle index}
\label{neg:all-windings}

The zero in \eqref{neg:pole-factor} has a useful more precise
estimate.  Set
\[
 v_i=1-\cos\delta_i,\qquad \bar v=\frac1n\sum_i v_i .
 \]
At \(t=\bar v\), the linear terms in the expansion of
\(\operatorname{arcosh}(b+B(\bar v)(\bar v-v_i))\) cancel on
averaging.  The remaining average is bounded by
\(C\,n^{-1}\sum_i(v_i-\bar v)^2\).  The uniform nonzero
derivative used in \eqref{neg:pole-factor} therefore yields
\[
 T_n=\bar v+
       O\!\left(\frac1n\sum_i(v_i-\bar v)^2\right).
 \]
Since \(v_i\ge0\), \(\max_i v_i\le C\sigma\), and
\(\bar v=(S/(2n))(1+O(\sigma))\), this implies
\begin{equation}
 |T_n(q)|\ge(1-C\sigma)\frac{S}{2n},
 \qquad
 |\partial_t\Gamma_n(T_n(q),q)|\ge cn .
 \label{neg:pole-lower}
\end{equation}

Retain the exact central numerator size
\begin{equation}
 H_{0,n}=
 \frac{|B(0)|^{n^2/2}}
 {n!(2\pi)^{n-1}2^{n(n-1)}s_\phi^{n^2}} .
 \label{neg:H0}
\end{equation}
Dividing the sine factors by their linear terms and all the
hyperbolic sine factors by \(\ii s_\phi\) gives
\begin{equation}
 |H_n(T_n(q),q)|
 \le H_{0,n}\Delta(\delta)^2e^{CnS},
 \qquad
 \Delta(\delta)=\prod_{i<k}(\delta_i-\delta_k).
 \label{neg:pole-density-upper}
\end{equation}
Here is the relevant accounting.  The logarithm of the anisotropy
ratio costs \(Cn^2|T_n|\le CnS\).  The sum of the absolute
one-body phase deviations is at most \(CS\), so all pair
denominators together cost \(CnS\).  Finally
\(\sum_{i<k}|\delta_i-\delta_k|^2\le CnS\) controls the sine
divided differences.  Thus no undetermined constant to the power
\(n^2\) replaces the displayed normalization.

For an odd winding \(w\), write
\begin{equation}
 \delta_i=\frac{w\pi}{n}+\eta_i,\qquad
 \sum_i\eta_i=0,\qquad
 Q(\eta)=\sum_i\eta_i^2,\qquad
 m_{n,w}=\frac{\pi^2w^2}{n}.
 \label{neg:winding}
\end{equation}
Then \(S=Q+m_{n,w}\) and \(\Delta(\delta)=\Delta(\eta)\).
The winding region is exactly \(Q<\sigma-m_{n,w}\) when it is
nonempty.  In particular there are at most \(C\sqrt n\) windings.
Define
\begin{equation}
 \alpha_n=\frac{n^2-1}{2},\qquad
 P_n=\prod_{r=1}^{n-1}\frac{r!}{2^r},\qquad
 G_n=\frac{n!\pi^{(n-1)/2}P_n}{\sqrt n}.
 \label{neg:gaussian-constant}
\end{equation}
On the real zero-sum plane, with measure
\(\dd\eta_1\cdots\dd\eta_{n-1}\),
\begin{align}
 \int e^{-Q(\eta)}\Delta(\eta)^2\,\dd\eta&=G_n,
 \label{neg:gaussian-period}\\
 \int\frac{\Delta(\eta)^2}{(Q(\eta)+m)^{j+1}}\,\dd\eta
 &=G_n m^{\alpha_n-j-1}
       \frac{\Gamma(j+1-\alpha_n)}{\Gamma(j+1)}
       \quad(j+1>\alpha_n).
 \label{neg:beta-period}
\end{align}
To verify the constants, the unrestricted Gaussian Vandermonde
integral is
\(n!\prod_{r=0}^{n-1}(\sqrt\pi\,r!/2^r)
 =n!\pi^{n/2}P_n\).
This follows by writing the Vandermonde as a determinant of monic
Hermite polynomials and integrating the determinant product.
The monic squared norm is \(\sqrt\pi\,r!/2^r\), obtained by
integration by parts.  In the decomposition
\(x_i=\eta_i+u\), \(\sum\eta_i=0\), the ordinary-coordinate
Jacobian is \(n\), and the mean integral is
\(\int e^{-nu^2}\dd u=\sqrt{\pi/n}\).  This proves
\eqref{neg:gaussian-period}.  Homogeneity gives radial density
proportional to \(Q^{\alpha_n-1}\dd Q\), proving
\eqref{neg:beta-period} by the ordinary beta integral.

Let \(A^{\rm pole}_{n,j,w}\) denote the contribution of the \(w\)th
winding to \eqref{neg:pole-coefficient}.  Equations
\eqref{neg:pole-lower}--\eqref{neg:beta-period} imply
\begin{align}
 |A^{\rm pole}_{n,j,w}|
 &\le \frac{CH_{0,n}}n e^{Cn\sigma}
    \left(\frac{2n}{1-C\sigma}\right)^{j+1}
    G_n m_{n,w}^{\alpha_n-j-1}
    \frac{\Gamma(j+1-\alpha_n)}{\Gamma(j+1)}
 \quad(j+1>\alpha_n).
 \label{neg:winding-bound}
\end{align}
The dependence on \(w\) here is the explicit factor
\(|w|^{-2(j+1-\alpha_n)}\).

There is also a bound without any condition on \(j/\alpha_n\).
Use \(|T_n|\ge c/n^2\), and use homogeneity to write
\[
 \int_{Q<\sigma}\Delta(\eta)^2\,\dd\eta
       =G_n\frac{\sigma^{\alpha_n}}{\Gamma(\alpha_n+1)}.
 \]
The elementary factorial bounds in \eqref{neg:gaussian-constant}
then give
\begin{equation}
 \log\!\left(\sum_w|A^{\rm pole}_{n,j,w}|\right)
 \le 2j\log n-\frac12n^2\log n+Cn^2+Cj+C .
 \label{neg:all-index-bound}
\end{equation}
In detail,
\(\log P_n=\tfrac12n^2\log n+O(n^2)\) and
\(\log\Gamma(\alpha_n+1)=n^2\log n+O(n^2)\);
the remaining factors in \(H_{0,n}G_n\) have logarithm \(O(n^2)\).
The number of windings costs only \(O(\log(n+1))\).
This proves \eqref{neg:all-index-bound} also when \(n\) is much
larger than the Taylor order.


\subsection{Uniform scalar poles on the original real plane}
\label{neg:real-poles}

We first make the parameter dependence precise.  Let
\(\mathcal P_0\Subset\mathcal P_1\) be concentric parameter polydisks,
let \(R,M,\beta>0\), and suppose that
\[
 B\in\mathcal O(\{|t|<R\}\times\mathcal P_1),\qquad
 |B(t,\eta)|\le M,\qquad
 \Im B(0,\eta)\le-\beta\quad(\eta\in\mathcal P_1).
\]
For a relatively compact admissible parameter set the sign hypothesis
can always be arranged on a slightly larger polydisk by continuity;
a finite cover suffices if several charts are needed.  Every constant
below is independent of \(\eta\in\mathcal P_0\), \(n\), and \(j\),
unless another dependence is displayed.  Constants may depend on the
fixed phase, \(R,M,\beta\), and the polydisk collar.

\begin{lemma}[Common scalar neighborhoods]\label{neg:family-neighborhoods}
There are common positive parameter and angular radii on which the
branch \(g(t,\delta,\eta)\), all fixed scalar derivatives needed below,
the nonzero one-body and canceled-pair denominators, and the logarithm
of \(B\) are uniformly controlled.  The estimates
\eqref{neg:gaussian}, \eqref{neg:small-disk},
\eqref{neg:regular-remainder}, \eqref{neg:winding-bound}, and
\eqref{neg:all-index-bound} consequently hold with common constants.
For each fixed \(j\) and parameter multiindex \(\nu\),
\[
 \sum_{n\ {\rm odd}}\sup_{\eta\in\mathcal P_0}
 \left|\partial_\eta^\nu[t^j]C_n(t,\eta)\right|<\infty.
\]
The same majorant permits the inward boundary limit for these fixed
mixed derivatives.
\end{lemma}
\begin{proof}
On \(|t|\le R/2\), Cauchy's inequality bounds \(B_t\) by \(2M/R\).
Choose \(r>0\) so small that \(|B(t,\eta)-B(0,\eta)|\le\beta/2\)
for \(|t|<r\).  Thus \(-\Im B\ge\beta/2\) and \(|B|\ge\beta/2\)
there.  The lower half-plane supplies a common logarithm, whose
possible physical sign change multiplies every odd coefficient by the
same nonzero sign.

The compact set
\[
 \{b-Bv: |B|\le M,\ \Im B\le-\beta,\ 0\le v\le2\}
\]
avoids both branch cuts of the branch of \(\arccos\) through \(\phi\):
at \(v=0\) it equals \(b\in(-1,1)\), and at \(v>0\) its imaginary
part is positive.  It therefore has a fixed neighborhood on which
\(g=\ii\arccos W\) is holomorphic.  Uniform continuity of cosine
gives common smaller \(t\) and complex-angle neighborhoods.
At \(t=0\) the corresponding roots form a compact subset of the
closed unit disk strictly inside its lower half-plane.  The strict
root-Schur estimate follows uniformly from \eqref{neg:schur}, and
persists after one further common shrink.  All canceled factors and
one-body denominators have uniform bounds on this neighborhood.
Cauchy's inequality supplies every fixed scalar or parameter jet.

The proofs of the cited estimates use just these bounds, a fixed
lower bound for \(|B(0)/\sin\phi|\), and the fixed rational phase.
Thus their disk sizes, \(\sigma\), Gaussian exponent, and constants
can be chosen in common.  For parameter derivatives first work on an
intermediate polydisk \(\mathcal P_0\Subset\mathcal P_*\Subset
\mathcal P_1\).  Cauchy's inequality in its fixed collar bounds each
\(\partial_\eta^\nu[t^j]C_n\) by a constant times
\((Cn^2)^j CnE_n\).  This is summable for fixed \(j,\nu\), including
on the disks of radius \(c(\epsilon+n^{-2})\) about \(-\epsilon\).
Dominated convergence proves the assertions.
\end{proof}

The rest of the proof retains the actual function \(B(t,\eta)\).
For an odd multiple \(n\) of \(p\), put
\begin{equation}
 d_n=\frac\pi n,\qquad \tau_n=1-\cos d_n,\qquad
 B_n=B(\tau_n,\eta),\qquad
 h_n=\cos d_n-\frac{B_nb}{s_\phi^2}\sin^2d_n,
 \qquad a_n=\Re h_n.
 \label{neg:center}
\end{equation}
Uniformly, \(h_n=1+O(n^{-2})\) and \(a_n>0\) for sufficiently
large \(n\).  In the first winding write
\[
 \delta_i=d_n+v_i,\qquad
 V_n=\{v\in\mathbb C^n:\sum_i v_i=0\},\qquad
 Q(v)=\sum_i v_i^2.
\]
On the real plane the measure is \(\dd v_1\cdots\dd v_{n-1}\),
with \(v_n=-\sum_{i<n}v_i\).  The exact first-winding region
\(S<\sigma\) is
\begin{equation}
 Q<r_n^2,\qquad r_n^2=\sigma-\frac{\pi^2}{n}.
 \label{neg:actual-real-ball}
\end{equation}
For large \(n\), its squared radius lies between \(\sigma/2\) and
\(\sigma\).

The exact scalar pole is determined by
\[
 \sum_i\bigl(g(T_n(v),d_n+v_i,\eta)-\ii\phi\bigr)=0.
\]
Subtracting the common constant and linear terms in the \(v_i\)
bounds the remaining sum by \(C\sum_i|v_i|^2\).  The parameter
derivative of its average stays in a fixed disk around
\(g_t(\tau_n,d_n,\eta)=-\ii B_n/s_\phi\), bounded away from zero.
Rouch\'e's theorem on the circle
\(|t-\tau_n|=C\sum_i|v_i|^2/n\), followed by uniqueness, gives a
holomorphic solution on a common small complex Euclidean ball in
\(V_n\), jointly with \(\eta\), and
\begin{equation}
 T_n(0)=\tau_n,\qquad
 |T_n(v)-\tau_n|\le\frac Cn\sum_i|v_i|^2.
 \label{neg:real-pole-coarse}
\end{equation}
All scalar derivative bounds used here are on fixed neighborhoods,
so the radii do not shrink with the dimension.

At the center direct differentiation gives
\[
 g(\tau_n,d_n)=\ii\phi,\qquad
 g_t(\tau_n,d_n)=-\frac{\ii B_n}{s_\phi},\qquad
 \frac12g_{\delta\delta}(\tau_n,d_n)
       =\frac{\ii B_nh_n}{2s_\phi}.
\]
In \(W_t=B(t)+B_t(t)(t-1+\cos\delta)\), the second term
vanishes exactly at this center; no derivative of \(B\) was frozen.
Evenness in \(\delta\) gives
\(g_{\delta\delta\delta}(t,d_n)=O(d_n)\) uniformly.  Writing
\(u=T_n-\tau_n\), Taylor's formula and \(\sum_i v_i=0\) yield
\[
 0=ng_t(\tau_n,d_n)u+\tfrac12g_{\delta\delta}(\tau_n,d_n)Q
 +O\!\left(n|u|^2+|u|Q+
              d_n\sum_i|v_i|^3+\sum_i v_i^4\right)
\]
on the real ball.  Inserting \eqref{neg:real-pole-coarse} proves
\begin{equation}
 \begin{split}
 T_n(v)&=T_{0,n}(v)+\mathcal R_n(v),\qquad
 T_{0,n}(v)=\tau_n+\frac{h_nQ}{2n},\\
 |\mathcal R_n(v)|&\le\frac Cn\left(
     \frac1n\sum_i|v_i|^3+\sum_i v_i^4+\frac{Q^2}{n}\right).
 \end{split}
 \label{neg:real-pole-error}
\end{equation}
This estimate retains the fourth-power sum: replacing it pointwise
by \(CQ^2/n\) would be false.  Since \(\sum|v_i|^3\le Q^{3/2}\)
and \(\sum v_i^4\le Q^2\), a common sufficiently small \(\sigma\)
and large \(n\) give
\begin{equation}
 |\mathcal R_n(v)|\le\tfrac14D_n(v),\qquad
 D_n(v)=\tau_n+\frac{a_nQ}{2n}\le|T_{0,n}(v)|,
 \qquad Q<\sigma.
 \label{neg:real-pole-small}
\end{equation}

\subsection{The complete residue and its angular measure}
\label{neg:real-density}

In this first winding put
\[
 \varrho_n(v)=-\frac{2H_n(T_n(v),\pi+d_n+v)}
                       {\partial_t\Gamma_n(T_n(v),\pi+d_n+v)}.
\]
All factors in \(H_n\), including \(B(T_n,\eta)^{n^2/2}\), are
retained.  After filling the sine divided differences,
\(\varrho_n/\Delta(v)^2\) is holomorphic and nonzero on a common
small ball.  Its exact center value is
\begin{equation}
 \varrho_{n,0}=-\frac{2B_n^{\alpha_n-1/2}}
 {n\,n!(2\pi)^{n-1}2^{n(n-1)}s_\phi^{2\alpha_n}},
 \qquad \alpha_n=\frac{n^2-1}{2}.
 \label{neg:center-residue}
\end{equation}
Indeed \(\ii^{-n^2}=-\ii\) for odd \(n\), so the centered
numerator divided by \(\Delta^2\) is
\[
 -\frac{\ii B_n^{n^2/2}}
 {n!(2\pi)^{n-1}2^{n(n-1)}s_\phi^{n^2}},
 \qquad \partial_t\Gamma_n=-\frac{\ii nB_n}{s_\phi}.
\]
The remaining
minus sign comes from the Taylor expansion of \((t-T_n)^{-1}\).

\begin{lemma}[A direct residue-density bound]\label{neg:direct-density}
On the real ball \(Q<\sigma\),
\begin{equation}
 \varrho_n(v)=\varrho_{n,0}\Delta(v)^2e^{\ell_n(v)},
 \qquad |\ell_n(v)|\le CnQ.
 \label{neg:full-density}
\end{equation}
The logarithm has a holomorphic extension through the coincidence
divisors, and the estimate is uniform in the parameter family.
\end{lemma}
\begin{proof}
Let \(e_i=g(T_n,d_n+v_i)-\ii\phi\).  The exact pole equation
gives \(\sum_i e_i=0\), and the scalar Taylor bounds give
\[
 |e_i|\le C\bigl(|u|+|d_n v_i|+|v_i|^2\bigr),\qquad
 \sum_i|e_i|^2\le C\left(\frac{Q^2}{n}
                      +\frac{Q}{n^2}+\sum_i v_i^4\right).
\]
On a common small disk the function
\(L(z)=\log(\sinh(\ii\phi+z)/(\ii s_\phi))\) is holomorphic
and \(L(z)=L'(0)z+O(|z|^2)\).  Thus the linear parts in both
\(\sum_iL(e_i)\) and
\(2\sum_{i<k}L((e_i+e_k)/2)\) vanish exactly.  Their total
absolute logarithmic cost is at most
\(Cn\sum_i|e_i|^2\le CnQ\) on the fixed ball.
The sine divided differences cost
\(C\sum_{i<k}(v_i-v_k)^2=CnQ\).  The anisotropy factor costs
\(Cn^2|u|\le CnQ\).  Finally, summing the scalar Taylor expansion
of \(g_t(T_n,d_n+v_i)\) cancels its linear angular term and gives
\[
 \frac{\sum_i g_t(T_n,d_n+v_i)}{ng_t(\tau_n,d_n)}
       =1+O(Q/n).
\]
Taking the logarithm of each already normalized nonzero factor
proves the result.  In particular no Jacobian or nonlinear angular
change of variables enters this density.
\end{proof}

Separate radius and angle in the positive real measure
\(\Delta(v)^2\dd v\): write \(v=\sqrt Q\,\theta\), with
\(\sum_i\theta_i=0\) and \(\sum_i\theta_i^2=1\).  Let
\(\nu_n\) be the normalized angular measure proportional to
\(\Delta(\theta)^2\).

\begin{lemma}[A scalar angular exponential moment]
\label{neg:angular-moment}
For \(Z_n=n\max_i\theta_i^2\), there are absolute constants
\(c_0,C_0>0\) such that
\begin{equation}
 \sup_{n\ge3}\mathbb E_{\nu_n}[(1+Z_n)e^{c_0Z_n}]\le C_0.
 \label{neg:angular-exponential}
\end{equation}
One may take \(c_0=1/32\).
\end{lemma}
\begin{proof}
The constrained Gaussian probability measure
\(G_n^{-1}e^{-Q}\Delta(v)^2\dd v\) has the normalization
\eqref{neg:gaussian-constant}.  Homogeneity shows that \(Q\) and
\(\theta\) are independent, the angular law is \(\nu_n\), and
\(Q\) has gamma density \(q^{\alpha_n-1}e^{-q}/\Gamma(\alpha_n)\).

Fix \(v_1=x\), let \(m=n-1\) and \(c=n/(n-1)\), and write
\(v_{i+1}=u_i-x/m\), \(\sum_i u_i=0\).  The independent-coordinate
Jacobian is one, and
\[
 Q(v)=cx^2+Q(u),\qquad
 \Delta(v)^2=\Delta(u)^2\prod_{i=1}^m(cx-u_i)^2.
\]
Hence the density of the first coordinate is
\[
 f_n(x)=\frac{G_{n-1}}{G_n}e^{-cx^2}
       \mathbb E_m\prod_{i=1}^m(cx-u_i)^2,
 \qquad
 \frac{G_{n-1}}{G_n}=
 \frac{2^{n-1}}{\sqrt\pi\,n!}\sqrt{\frac n{n-1}}.
\]
By arithmetic--geometric mean and \(\sum_i u_i=0\),
\[
 \prod_i(cx-u_i)^2\le\left(c^2x^2+\frac{Q(u)}m\right)^m.
\]
For \(0\le k\le m\), the gamma moments obey
\[
 \mathbb E_m\left(\frac{Q(u)}m\right)^k
 =\frac{\alpha_m(\alpha_m+1)\cdots(\alpha_m+k-1)}{m^k}
 \le n^k,
\]
since \((\alpha_m+m-1)/m\le n\).  Using \(c^2<3\), the
binomial expansion therefore bounds the expected product by
\((3x^2+n)^{n-1}\).  Since \(n!\ge(n/e)^n\), with \(t=x^2/n\),
\[
 f_n(x)\le\frac1n(2e)^n(1+3t)^{n-1}e^{-x^2}.
\]
For \(t\ge16\),
\(\log(2e)+\log(1+3t)\le t/2\); the inequality holds at
\(16\), and the difference has positive derivative thereafter.
It follows that
\(f_n(x)\le n^{-1}e^{-x^2/2}\) when \(x^2\ge16n\).
Permutation invariance of the constrained measure and the union
bound now give
\[
 \Pr\{\max_i v_i^2\ge A^2\}
 \le2\int_A^\infty e^{-x^2/2}\dd x
 \le\frac2A e^{-A^2/2},\qquad A^2\ge16n.
\]
On the other hand,
\[
 \Pr\{Q<n^2/8\}
 \le\frac{n^2}{8}\mathbb E(Q^{-1})
 =\frac{n^2}{4(n^2-3)}\le\frac38\qquad(n\ge3).
\]
Independence of radius and angle and
\(\max_i v_i^2=QZ_n/n\) consequently imply
\[
 \frac58\Pr\{Z_n\ge z\}\le
       \Pr\{\max_i v_i^2\ge nz/8\}.
\]
For \(z\ge128\), the preceding Gaussian tail gives
\(\Pr\{Z_n\ge z\}\le(16/5)e^{-nz/16}\).
Integrating this tail, and bounding \(Z_n\le128\) directly on
the complementary event, proves \eqref{neg:angular-exponential}.
\end{proof}

The consequences used below are the pointwise inequalities
\begin{equation}
 \sum_i v_i^4\le\frac{Q^2Z_n}{n},\qquad
 \sum_i|v_i|^3\le\sqrt{Q\sum_i v_i^4}
             \le\frac{Q^{3/2}\sqrt{Z_n}}{\sqrt n}.
 \label{neg:angular-power-bounds}
\end{equation}
The exponential moment includes the rare directions where
\(Z_n\) is large.

\subsection{An exact complex beta period and the two first windings}
\label{neg:beta-evaluation}

For \(j+1>\alpha_n\) and \(\Re h>0\), the same constrained
Gaussian period and the scalar Laplace formula give
\begin{equation}
 J_h:=\int_{V_n\cap\mathbb R^n}
       \frac{\Delta(v)^2\dd v}{(\tau_n+hQ/(2n))^{j+1}}
 =G_n(2n)^{\alpha_n}h^{-\alpha_n}
     \tau_n^{\alpha_n-j-1}
       \frac{\Gamma(j+1-\alpha_n)}{\Gamma(j+1)}.
 \label{neg:first-beta}
\end{equation}
Indeed the radial measure is
\(G_nQ^{\alpha_n-1}\dd Q/\Gamma(\alpha_n)\).  Applying the
Laplace representation of the negative power and then integrating
this radial gamma integral proves the formula.  Absolute convergence
follows from \(\Re h>0\) and \(j+1>\alpha_n\), so every integration
is on the original real variables.  In particular the gamma quotient
is exact.

For the positive majorant use \(a_n=\Re h_n\).  The exact periods
satisfy
\begin{equation}
 \frac{J_{a_n}}{|J_{h_n}|}
 =\left(\frac{|h_n|}{a_n}\right)^{\alpha_n}
 =\exp(O(n^{-2})).
 \label{neg:complex-beta-cost}
\end{equation}
This is uniform: \(a_n=1+O(n^{-2})\),
\(\Im h_n=O(n^{-2})\), and therefore
\(\log(|h_n|/a_n)=O(n^{-4})\).  The comparison has no factor
growing with \(j\).

We shall need a relative estimate only in the entire intermediate
range
\begin{equation}
 \log j<n<\frac{\sqrt j}{(\log j)^{1/4}}.
 \label{neg:middle-range}
\end{equation}
Here both \(n\to\infty\) and \(\epsilon_n=n^2/j\to0\) hold
uniformly.  The exact first-winding contribution is
\[
 I_{n,j}=\int_{Q<r_n^2}
       \varrho_n(v)T_n(v)^{-j-1}\dd v.
\]
We prove
\begin{equation}
 I_{n,j}=\varrho_{n,0}J_{h_n}
             \left(1+O\left(\frac n{\sqrt j}\right)\right)
 \label{neg:one-winding-model}
\end{equation}
uniformly in \eqref{neg:middle-range} and the parameter family.

Normalize the positive measure for \(J_{a_n}\) and set
\[
 x=\frac{a_nQ}{2n\tau_n},\qquad \epsilon=\frac{n^2}{j}.
\]
Radius and angle remain independent.  The angular law is \(\nu_n\),
and the radial density is proportional to
\(x^{\alpha_n-1}(1+x)^{-j-1}\).  Since
\(2n^2\tau_n\) and \(a_n\) are bounded above and below,
\(Q\asymp x/n\).  Equations \eqref{neg:real-pole-error} and
\eqref{neg:angular-power-bounds} give
\begin{equation}
 (j+1)\frac{|\mathcal R_n|}{D_n}
 \le C\frac j{n^2}
   \frac{x^{3/2}\sqrt{Z_n}+x^2(Z_n+1)}{1+x}.
 \label{neg:weighted-pole-error}
\end{equation}
Because \(|\mathcal R_n/T_{0,n}|\le1/4\), its near-one
logarithm has modulus at most \(2|\mathcal R_n|/D_n\).  Thus
\begin{equation}
 \left|\frac{e^{\ell_n}}{T_n^{j+1}}-
                  \frac1{T_{0,n}^{j+1}}\right|
 \le D_n^{-j-1}\left[
    \exp\left\{|\ell_n|+
           \frac{2(j+1)|\mathcal R_n|}{D_n}\right\}-1\right].
 \label{neg:real-integrand-error}
\end{equation}

Fix \(K=32\).  On \(x\le K\epsilon\),
\eqref{neg:full-density} gives \(|\ell_n|\le Cx\le C\epsilon\),
and \eqref{neg:weighted-pole-error} bounds the exponent by
\[
 C\{\sqrt\epsilon\sqrt{Z_n}+\epsilon(Z_n+1)\}
       \le C'\sqrt\epsilon(1+Z_n).
\]
For uniformly large \(j\), its coefficient is less than \(c_0\).
The mean-value theorem for the exponential and
Lemma~\ref{neg:angular-moment} therefore show that the angular
expectation of the bracket in \eqref{neg:real-integrand-error}
is \(O(\sqrt\epsilon)\).  This controls all angular directions
in the central radial region.

On the actual ball, \(x\le Ln\) for a fixed uniform \(L\).
The deterministic bound \(Z_n\le n\) gives
\[
 \frac{2(j+1)|\mathcal R_n|}{D_n}
 \le Cj\frac{x^{3/2}/n^{3/2}+x^2/n}{1+x}.
\]
Dividing by \(j\log(1+x)\), the right side tends to zero
uniformly on \(0<x\le Ln\): the intervals
\((0,1]\), \([1,\sqrt n]\), and \([\sqrt n,Ln]\) give
respectively \(O(n^{-1})\), \(O(n^{-1/2})\), and
\(O(1/\log n)\).  Also
\(Cx/[j\log(1+x)]=o(1)\), using \(j\to\infty\) below one
and \(n/j\to0\) above one.  Hence on the remaining radial region
\[
 |\ell_n|+\frac{2(j+1)|\mathcal R_n|}{D_n}
 \le\frac j4\log(1+x),\qquad K\epsilon\le x\le Ln.
\]
Write \(Z_\beta(\alpha,J)=\Gamma(\alpha)
\Gamma(J+1-\alpha)/\Gamma(J+1)\).  Dividing the resulting tail
integral by the original normalization bounds it by
\[
 C(1+K\epsilon)^{-j/4}
       \frac{Z_\beta(\alpha_n,j/2)}{Z_\beta(\alpha_n,j)}
 \le C(1+K\epsilon)^{-j/4}4^{\alpha_n}
 \le Ce^{-cn^2}.
\]
For the middle inequality, \(\alpha_n\) is an integer and
\(\alpha_n\le j/4\) eventually, so the exact ratio is
\[
 \prod_{r=0}^{\alpha_n-1}\frac{j-r}{j/2-r}\le4^{\alpha_n}.
\]
For the last inequality use
\(\log(1+K\epsilon)\ge K\epsilon/2\).  The omitted part of
the unperturbed period outside \(Q<r_n^2\) has the same bound:
its radial lower endpoint is at least a positive constant times
\(n\), hence exceeds \(K\epsilon\), and its positive majorant
has no perturbation factor.

Combining these estimates with \eqref{neg:complex-beta-cost}
gives relative error at most
\(C(n/\sqrt j+e^{-cn^2})\).  The lower bound \(n>\log j\)
absorbs the second term into the first.  This proves
\eqref{neg:one-winding-model} on the exact original ball; no
cycle-boundary remainder was introduced.

Reflection of all lifts \(\delta_i\mapsto-\delta_i\) preserves
the dispersion, every squared sine factor, and the independent
real measure.  It exchanges the first two windings.  Their
integrals are therefore equal, with their full complex phase.
Substitution of \eqref{neg:center-residue} and
\eqref{neg:gaussian-constant} into \eqref{neg:first-beta} gives
\begin{align}
 A^{\rm pole}_{n,j,1}+A^{\rm pole}_{n,j,-1}
   &=M(n,j)(1+e_{n,j}),\label{neg:local-model}\\
 |e_{n,j}|&\le C\frac n{\sqrt j},\label{neg:local-error}\\
 M(n,j)&=-\frac{2^{2-\alpha_n}B_n^{\alpha_n-1/2}P_n
                        n^{\alpha_n-3/2}}
              {\pi^{(n-1)/2}s_\phi^{2\alpha_n}}
      h_n^{-\alpha_n}\tau_n^{\alpha_n-j-1}
       \frac{\Gamma(j+1-\alpha_n)}{\Gamma(j+1)}.
       \label{neg:model}
\end{align}
This is the same exact-factor model for the whole residue, including
\(B(\tau_n,\eta)\) and its particle-dependent complex phase.
There is no positivity assumption across different particle indices.



\subsection{Integer selection of a single dominant particle index}
\label{neg:dominance}

Let \(L(n,j)=\log|M(n,j)|\) and \(d=2p\).  All estimates in
this subsection are uniform in the parameter family fixed above;
the selected integer may depend on \(\eta\).  We use only integer
Taylor orders and retain the exact gamma quotient.

For \(A=\pi^2/2\), put
\[
 e_n=\log(\tau_n n^2/A),\qquad
 E_j(a)=\sum_{r=0}^{a-1}-\log(1-r/j).
\]
For \(j>\alpha_n\), direct grouping of the factors in
\eqref{neg:model} gives
\begin{equation}
 \begin{split}
 L(n,j)={}&\log P_n+(2j-n^2/2+1)\log n-\alpha_n\log j
       -j\log A+R_n-je_n+E_j(\alpha_n).
 \end{split}
 \label{neg:L-exact}
\end{equation}
Here \(R_n\) is a quadratic polynomial in \(n\), plus a convergent
power series in \(n^{-2}\) of order \(O(n^{-2})\).  Its coefficients
are uniformly bounded.  Indeed \(\tau_n n^2/A\),
\(B(\tau_n,\eta)/B(0,\eta)\), and \(h_n\) have common convergent
expansions in \(n^{-2}\) near zero, with uniformly nonzero constant
terms.  Their near-one logarithms have common convergent expansions;
multiplication by \(\alpha_n\) introduces at most a constant term.
In particular \(e_n=O(n^{-2})\), its centered first and second
differences are \(O_d(n^{-3})\) and \(O_d(n^{-4})\), and those of
\(R_n\) are \(O_d(n)\) and \(O_d(1)\).

Ordinary Stirling bounds and
\(E_j(\alpha_n)=O(n^4/j)\) give, as \(n\to\infty\) and
\(n^2/j\to0\),
\begin{equation}
 L(n,j)=2j\log n-\frac{n^2}{2}\log j-j\log(\pi^2/2)
       +O(n^2+n^4/j+j/n^2).
 \label{neg:L-asymptotic}
\end{equation}
The finite differences needed for dominance are obtained from the
exact expression \eqref{neg:L-exact}, not from its error term:
\begin{align}
 L(n+d,j)-2L(n,j)+L(n-d,j)
  &=-d^2\left(\frac{2j}{n^2}+\log j\right)
       +O_d(1+n^2/j+j/n^4),\label{neg:second-difference}\\
 \frac{L(n+d,j)-L(n-d,j)}{2d}
  &=\frac{2j}{n}-n\log j+O_d(n+n^3/j+j/n^3).
       \label{neg:first-difference}
\end{align}
These hold for \(j>\alpha_{n+d}\) in the stated regime.

Here is a direct verification of their remainders.  The factorial
products have the exact differences
\[
 \log P_{n+d}-\log P_{n-d}
  =\sum_{r=n-d}^{n+d-1}\bigl(\log(r!)-r\log2\bigr),
\]
\[
 \log P_{n+d}-2\log P_n+\log P_{n-d}
  =\sum_{k=0}^{d-1}\sum_{\ell=1}^d\log(n-d+k+\ell)-d^2\log2.
\]
Stirling's inequality on the finitely many factorials makes the
first \(2dn\log n+O_d(n)\); the second is
\(d^2\log n+O_d(1)\).  The corresponding elementary differences
of \((2j-n^2/2+1)\log n\) cancel these \(\log n\) terms, leaving
\(2j/n+O_d(n+j/n^3)\) in the centered first difference divided
by \(2d\), and \(-2d^2j/n^2+O_d(1+j/n^4)\) in the second.

For the gamma correction, on \(0\le r\le\alpha_{n+d}\),
\[
 0\le-\log(1-r/j)\le Cn^2/j,\qquad
 \frac{\dd}{\dd r}[-\log(1-r/j)]\le C/j.
\]
The difference \(E_j(\alpha_{n+d})-E_j(\alpha_{n-d})\) has
\(2dn\) summands and is \(O_d(n^3/j)\).  For the second
difference, the upper and lower increments of \(\alpha_n\) have
lengths \(dn+d^2/2\) and \(dn-d^2/2\).  Their \(d^2\) unmatched
summands cost \(O_d(n^2/j)\); pairing the others, whose arguments
differ by \(O_d(n)\), gives the same bound.  Finally the differences
of \(-\alpha_n\log j\) are exactly \(-n\log j\) after division
by \(2d\), and \(-d^2\log j\).  The bounds on \(R_n\) and
\(je_n\) now prove \eqref{neg:second-difference}--
\eqref{neg:first-difference}.

Throughout \eqref{neg:middle-range}, the second difference is at
most \(-cd^2\log j\) for all sufficiently large \(j\), whenever
the neighboring indices lie in that range.  Indeed \(n^2/j=o(1)\)
uniformly, whereas \(j/n^4\) is smaller than \(j/n^2\) by the
factor \(n^{-2}\le(\log j)^{-2}\).

Fix a large allowed \(n\), and for integer \(j>\alpha_{n+d}\) set
\(F_n(j)=L(n+d,j)-L(n-d,j)\).  The one-step gamma recurrence
gives the exact identity
\begin{equation}
 F_n(j+1)-F_n(j)
  =\log\frac{\tau_{n-d}}{\tau_{n+d}}
       +\log\frac{j+1-\alpha_{n+d}}{j+1-\alpha_{n-d}}.
 \label{neg:integer-recurrence}
\end{equation}
On the integer interval
\(\tfrac12n^2\log n\le j\le\tfrac32n^2\log n\), this is
\(4d/n+O_d(1/(n\log n))\), hence positive and \(O_d(1/n)\).
It is independent of \(\eta\).

At the two endpoints \(j=cn^2\log n+O(1)\),
\(c=1/2,3/2\), the leading part of
\eqref{neg:first-difference} is
\[
 \frac{2j}{n}-n\log j
 =2(c-1)n\log n-n\log\log n-n\log c+O(1/n).
\]
The error there is \(O_d(n)\), uniformly in \(\eta\), so these
endpoint signs are opposite.  Define \(j_n(\eta)\) as the first
integer in this interval for which \(F_n(j_n)\ge0\).  Then
\eqref{neg:integer-recurrence} gives
\begin{equation}
 j_n\sim n^2\log n,\qquad
 0\le F_n(j_n)=O_d(1/n),
 \label{neg:selected-order}
\end{equation}
uniformly in the family.  For the first conclusion insert the
second in \eqref{neg:first-difference} and use
\(\log j_n=2\log n+O(\log\log n)\).  These \(n\) and their
neighbors lie strictly inside \eqref{neg:middle-range}.

Let \(\lambda_k=L(n+kd,j_n)\) on the admissible indices in that
range.  Their successive differences decrease by at least
\(cd^2\log j_n\), and
\(\lambda_1-\lambda_{-1}=O_d(1/n)\).  Therefore
\[
 \lambda_0-\lambda_{\pm1}\ge c'd^2\log j_n,\qquad
 \lambda_k-\lambda_0\le-c''d^2k^2\log j_n+O_d(|k|/n)
 \quad(k\ne0).
\]
Summing this discrete Gaussian gives
\begin{equation}
 \sum_{\substack{m\ne n,\ m\ {\rm odd},\ p\mid m\\
                       m\ {\rm satisfies}\ \eqref{neg:middle-range}}}
       |M(m,j_n)|
 \le |M(n,j_n)|\,O(j_n^{-c'''d^2}).
 \label{neg:single-index}
\end{equation}
The phases of the different models are unrestricted.  The selected
integer is not differentiated in \(\eta\); neither continuity nor
a common integer sequence across the parameter family is required.


\subsection{Exhausting the remaining indices and windings}
\label{neg:exhaustion}

We compare every part of \eqref{neg:jet-sum} with the selected
model.  First, \eqref{neg:L-asymptotic} and
\eqref{neg:selected-order} imply
\begin{align}
 \log|M(n,j_n)|
 &=j_n\log j_n-j_n\log\log j_n
          +j_n\bigl(\log(4/\pi^2)-1\bigr)+o(j_n),
 \label{neg:model-growth}\\
 |M(n,j_n)|^{1/j_n}
 &\sim\frac{4j_n}{e\pi^2\log j_n}.
 \label{neg:model-root}
\end{align}
Indeed \(j_n\sim n^2\log n\) gives
\[
 2\log n=\log j_n-\log\log j_n+\log2+o(1),\qquad
 \frac{n^2\log j_n}{2j_n}\longrightarrow1.
\]
Substituting these identities in \eqref{neg:L-asymptotic} gives
\eqref{neg:model-growth}; every error there is \(o(j_n)\).

The following estimates cover all the complementary terms.

\paragraph{Small particle indices.}
For \(m\le\log j\), the full coefficient bound
\eqref{neg:small-disk} gives
\begin{equation}
 \sum_{\substack{m\le\log j\\m\ {\rm odd}}}
       |[t^j]C_m(t)|
 \le \exp\{2j\log\log(j+2)+Cj+C\}.
 \label{neg:small-indices}
\end{equation}
This is negligible compared with \eqref{neg:model-growth}
when \(j=j_n\).  Equivalently, after the pole reduction one
can apply the same bound to the poles plus their analytic
remainders, since their difference has the smaller bound
\eqref{neg:regular-remainder}.

\paragraph{All analytic remainders.}
The sum of the regular remainders in
\eqref{neg:regular-remainder} is bounded by
\begin{equation}
 \exp\{\tfrac12j\log(j+2)+Cj+C\}.
 \label{neg:all-remainders}
\end{equation}
This too is negligible at the selected orders.  This estimate
includes the entire nonresonant particle class \(p\nmid m\)
outside the small-index range.

\paragraph{All sufficiently large particle indices.}
For
\[
 m\ge\frac{\sqrt j}{(\log j)^{1/4}},
 \]
use \eqref{neg:all-index-bound}, which has no restriction
on \(m^2/j\).  With \(x=m^2\), its exponent is at most
\[
 j\log x-\frac14x\log x+Cx+Cj+C .
 \]
Its derivative in \(x\) is negative on
\(x\ge j/\sqrt{\log j}\), for all large \(j\), and at the
lower endpoint its value is
\begin{equation}
 j\log j-\frac12j\log\log j
            -\frac14j\sqrt{\log j}
            +o(j\sqrt{\log j}).
 \label{neg:large-indices}
\end{equation}
For example its derivative is at most \(-c\log j\)
there for large \(j\), since
\(j/x\le\sqrt{\log j}\) and
\(\tfrac14\log x\sim\tfrac14\log j\).
Summing the resulting decreasing exponential over the
discrete values \(x=m^2\) costs at most a harmless polynomial
factor.  Thus \eqref{neg:large-indices} bounds the logarithm
of the whole infinite large-index tail up to
\(o(j\sqrt{\log j})\).
Its gap below \eqref{neg:model-growth} tends to infinity
faster than \(j\).

\paragraph{Higher odd windings in the intermediate range.}
Let \(M_0(m,j)>0\) be the modulus of the expression
\eqref{neg:model} with \(B_m\) replaced by \(B(0)\),
\(h_m\) by \(1\), and \(\tau_m\) by
\(\pi^2/(2m^2)\).  The exact factors give
\begin{equation}
 \left|\log\frac{M_0(m,j)}{|M(m,j)|}\right|
 \le C(1+j/m^2).
 \label{neg:M0}
\end{equation}
For the \(\tau\) factor this follows from
\(\log(\tau_m/(\pi^2/(2m^2)))=O(m^{-2})\);
the \(B\) and \(h\) factors cost \(O(1)\).
Comparison of \eqref{neg:winding-bound} with this
explicit \(M_0\) yields, for each odd \(|w|\ge3\),
\[
 |A^{\rm pole}_{m,j,w}|
 \le C M_0(m,j)e^{Cm\sigma}
       (1-C\sigma)^{-j-1}
       |w|^{-2(j+1-\alpha_m)} .
 \]
The constants \(C\) here do not depend on \(m,j,w\).
Choose the fixed \(\sigma\) still smaller, if necessary,
so that
\(-\log(1-C\sigma)<\tfrac12\log3\).
Throughout \eqref{neg:middle-range},
\[
 \alpha_m=o(j),\qquad m=o(j),\qquad j/m^2=o(j)
 \]
uniformly.  The number of possible windings is at most
\(C\sqrt m\).  Hence for some \(c>0\),
\begin{equation}
 \sum_{\substack{w\ {\rm odd}\\|w|\ge3}}
       |A^{\rm pole}_{m,j,w}|
 \le |M(m,j)|e^{-cj}
 \label{neg:higher-windings}
\end{equation}
uniformly in that range.  The choice of a smaller
\(\sigma\) merely changes the fixed constants in the
earlier arguments.

\paragraph{The first windings of every intermediate index.}
Equation \eqref{neg:local-error} tends uniformly to zero
in \eqref{neg:middle-range}; it is at most
\(C(\log j)^{-1/4}\).
At \(j=j_n\), \eqref{neg:single-index} shows that the sum
of the moduli of all intermediate models other than
\(M(n,j_n)\) is \(o(|M(n,j_n)|)\).
The sum of all their relative errors is bounded by
\[
 o(1)\sum_m|M(m,j_n)|=o(|M(n,j_n)|).
 \]
The same is true of all higher windings by
\eqref{neg:higher-windings}.

These classes exhaust every odd particle number and every
winding, together with all analytic remainders.
Combining them in \eqref{neg:jet-sum} proves the actual
full-sum coefficient asymptotic
\begin{equation}
 a_{j_n}=M(n,j_n)(1+o(1)).
 \label{neg:full-asymptotic}
\end{equation}
In particular, this is a statement about the sum of the
complex coefficients, not merely a comparison of individual
singularity sets.

\subsection{Nonextension and uniformity}
\label{neg:conclusion}

Equations \eqref{neg:model-root} and
\eqref{neg:full-asymptotic} prove
\eqref{neg:growth-statement}, so
\(\limsup_j|a_j|^{1/j}=\infty\).
If \(F\) admitted a holomorphic extension to a disk about
zero from the physical side \(\mathcal U^-\), its derivatives
on the inward negative ray would agree with those of that
extension.  By \eqref{neg:jets}, its Taylor coefficients at
zero would therefore be the \(a_j\).  Cauchy's estimate on
any smaller extension disk would bound
\(|a_j|\le C R^{-j}\), a contradiction.
If \(P(0)\ne0\), a holomorphic extension of \(PF\) would,
after division by the nonvanishing holomorphic factor
\(P\), give an extension of \(F\).  Thus a local physical
normalization by such a factor cannot remove the obstruction.

Finally, Lemma~\ref{neg:family-neighborhoods} supplies common
scalar neighborhoods and fixed mixed-derivative majorants for bounded
holomorphic families on common polydisks.  The real pole, residue, and
beta estimates above use those constants, the absolute angular moment,
and no contour depending on the family parameter.  The convergent
expansions used in the exact finite differences are uniform on the same
polydisk collar.  The common sign bracket and integer recurrence
therefore give uniform selection and all-index comparison, with the
actual values \(B(\tau_n,\eta)\) throughout.  The integers
\(j_n(\eta)\) may vary with \(\eta\); no common sequence or
parameter derivative of their choice is asserted.  Any nonramified
holomorphic transversal is represented by its actual \(B(t)\) in this
argument.  A compact admissible family of changes of its curve
coordinate has a common local inverse after shrinking the polydisks.
This completes the proof of Theorem~\ref{neg:theorem}.



\section{The isotropic restriction on the complete constrained cycle}
\label{iso:section}

Here $\lambda=1$, $s=\sinh z$, and the dispersion parameter will be
written ${E}$ to distinguish it from the notation of
Section~\ref{ana:physical}.  The full product formula already proved
there specializes to $B=1$ and $C=s+s^{-1}$.  We give the remaining
analytic argument directly for that product.

For ${E}\in\C\setminus[-2,2]$ and $y=e^{\ii\theta}$, let
\[
 W={E}-\cos\theta,\qquad
 \rho+\rho^{-1}=2W,\quad |\rho|<1,\qquad
 \sinh\gamma=(\rho^{-1}-\rho)/2.
\]
The root is unique and holomorphic: a root on the unit circle would
force $W\in[-1,1]$, and the two roots are reciprocal.  Write
\[
 P_{ij}=-\frac{(y_i-y_j)^2\rho_i\rho_j}
                  {y_i y_j(1-\rho_i\rho_j)^2},\qquad
 Z=\prod_i\rho_i,
 \qquad a_+=0,\quad a_-=\pi.
\]
For either sign $\sigma\in\{1,-1\}$ define the exact auxiliary coefficient
\begin{equation}\label{iso:angular}
 \mathcal F_N^\sigma({E})
 =\frac1{N!}\int_{\mathbb H_N^\sigma}
       \frac{1+\sigma Z}{1-\sigma Z}
       \prod_i\frac1{\sinh\gamma_i}\prod_{i<j}P_{ij}\,d\mu_N,
 \qquad
 \mathbb H_N^\sigma=\left\{\sum_i\theta_i=a_\sigma
                                    \pmod{2\pi}\right\},
\end{equation}
where
\[
 d\mu_N=\prod_{i<N}\frac{d\theta_i}{2\pi},\qquad
 \theta_N=a_\sigma-\sum_{i<N}\theta_i\pmod{2\pi}.
\]
Thus this is Euclidean hypersurface measure divided by
$\sqrt N(2\pi)^{N-1}$.  For $N=1$ the constrained integral is evaluation
at its single angle.  All pair factors are rational squares, so no
choice of half-angle or logarithm changes the coefficient.  Put
\begin{equation}\label{eq:iso-exterior-coefficients}
 C_N^\sigma(t)=\mathcal F_N^\sigma(t+t^{-1}),\qquad |t|>1.
\end{equation}
Write $E(\zeta)=\zeta+\zeta^{-1}$ and put $C_N=C_N^+$. The compact normal-convergence bounds for the full product apply to
each fixed differentiated series here.  To see directly that the
auxiliary sign causes no change, on a compact dispersion neighborhood
all roots satisfy $|\rho_i|\le q<1$ and have imaginary parts of one
sign.  The latter fact follows from
$\Im W_i=\Im\rho_i(1-|\rho_i|^{-2})/2$.
The identity
\[
 |1-vw|^2-|v-w|^2
   =(1-|v|^2)(1-|w|^2)+4\Im v\Im w
\]
therefore bounds every root Schur factor
$|(\rho_i-\rho_j)/(1-\rho_i\rho_j)|$ strictly below one, uniformly.
Expand the angular Schur product by the matching identity proved
in Section~\ref{ana:physical} and cancel each matched angular pole
with its corresponding root factor.  The resulting complete matched
pairs are bounded; all but $O(N)$ root factors retain the strict
contraction.  The number of matchings costs
$\exp(O(N\log(N+1)))$, while the unmatched root factors supply
$\exp(-cN^2)$.  This gives a summable coefficient bound on the
normalized constrained torus for either sign.  All one-body factors
and $1-\sigma Z$ have fixed gaps there.  Cauchy's inequality on a
slightly larger compact neighborhood gives every fixed derivative.

Angle reflection preserves $\mathbb H_N^\sigma$, its Haar measure and
each complete factor in \eqref{iso:angular}.  Conjugation also
commutes with the coefficients.  Replacing ${E},y_i,\rho_i$ by
$-{E},-y_i,-\rho_i$ gives the precise negation identity
\begin{equation}\label{iso:negation}
 \mathcal F_N^\sigma(-{E})
   =(-1)^N\mathcal F_N^{(-1)^N\sigma}({E}).
\end{equation}
Every pair is unchanged, every one-body factor changes sign, and
$\prod_i y_i$ and $Z$ both acquire $(-1)^N$.  This proves
\eqref{iso:negation} with its constraint and product-kernel signs.
In particular the even physical coefficients are even in $t$, whereas
the odd coefficients satisfy
\begin{equation}\label{eq:iso-odd-negation}
 C_N^+(-t)=-C_N^-(t)\qquad(N\ {\rm odd}).
\end{equation}
No evenness assertion about the HT observable is made.
The two one-particle formulas, useful below, are
\begin{equation}\label{eq:iso-one-particle}
 \mathcal F_1^+({E})=\frac1{{E}-2},\qquad
 \mathcal F_1^-({E})=\frac1{{E}+2}.
\end{equation}
For example the second formula follows by taking $\theta=\pi$ and
using $(1-\rho)/(1+\rho)\sinh(\gamma)^{-1}
=2\rho/(1+\rho)^2$.


\subsection{The dense arithmetic family}
\begin{lemma}[The selected first orders]\label{iso:arithmetic}
Let $p\ge11$ be prime, choose distinct
$a,b\in\{1,\ldots,(p-1)/2\}$, and set
\[
 \alpha=2\pi a/p,\quad\beta=2\pi b/p,\quad
 E_*=\cos\alpha+\cos\beta=2\cos\theta,
 \quad\zeta_*=e^{\ii\theta}.
\]
For the even family choose $\alpha,\beta\in(0,\pi/2)$ and put
$N_0=2p$, $k=N_0^2/2-1$. For the odd family choose both angles in
$(0,\pi/2)$ or both in $(\pi/2,\pi)$ and put
$N_0=p$, $k=(p^2-1)/2$.
The first order of the respective parity satisfying a Nickel equation is
$N_0$, and its unordered cosine pair is uniquely
$\{\cos\alpha,\cos\beta\}$. No endpoint configuration is resonant at any
order. The even family is dense in the upper-right arc; the odd families
are dense in both upper quadrants.
\end{lemma}

\begin{proof}
Let $\omega=e^{2\pi\ii/p}$, $d=(p-1)/2$,
$u_j=\omega^j+\omega^{-j}$, and
$F=\mathbb Q(\omega+\omega^{-1})$. Every nonzero-index $u_j$ has
trace $-1$. Thus $v=u_a+u_b$ has trace $-2$ and is irrational:
if rational it would be an integer, whereas its trace would force
$v=-2/d$. If $p\nmid n$, then
$\mathbb Q(\omega)\cap\mathbb Q(e^{2\pi\ii/n})=\mathbb Q$.
Indeed the compositum has degree
$\varphi(pn)=\varphi(p)\varphi(n)$, so the intersection has degree one.
A lower order of the required parity would contradict this fact.

The only rational relation between $1,u_1,\ldots,u_d$ is a multiple of
$1+\sum_j u_j=0$: lift the relation to a polynomial of degree at most
$p-1$ and use the minimal polynomial $1+X+\cdots+X^{p-1}$.
At order $p$ all cosines are $u_j/2$ or $1$; at order $2p$ they are
$\pm u_j/2$ or $\pm1$. An equality of two cosine pairs has at most four
nonconstant positions. Since $d\ge5$, a zero coefficient forces the
multiple of the relation to vanish, proving uniqueness.

For $E_*>0$, an endpoint resonance requires a root of unity $\xi$ with
$\xi+\xi^{-1}=v-2$. Its trace is $-2-2d<-2d$, impossible because
every conjugate of a root-of-unity trace belongs to $[-2,2]$.
For $E_*<0$ the possible endpoint equation is
$\xi+\xi^{-1}=v+2$. Direct expansion gives
$\operatorname{Tr}_F(v^2)=4d-6$. If all conjugates of $v+2$ were in
$[-2,2]$, every conjugate of $v$ would be in $[-4,0]$, where
$v^2\le-4v$. Taking traces would imply $4d-6\le8$, contrary to
$d\ge5$. Approximating any angle in either relevant half interval by
two distinct adjacent prime-grid angles proves density.
\end{proof}

\subsection{One analytic tail estimate for both isotropic parities}
\label{sec:iso-shared-tail}

We use the signed coefficients $\mathcal F_N^\sigma({E})$,
$\sigma\in\{1,-1\}$, on
$\mathbb H_N^\sigma=\{\sum\theta_i=a_\sigma\pmod {2\pi}\}$,
where $a_+=0$ and $a_-=\pi$.  Thus, with the independent-coordinate
Haar measure $\prod_{i<N}\dd\theta_i/(2\pi)$,
\begin{equation}\label{eq:iso-tail-coefficient}
 \mathcal F_N^\sigma({E})=\frac1{N!}
 \int_{\mathbb H_N^\sigma}
 \frac{1+\sigma Z}{1-\sigma Z}
 \prod_i\frac1{\sinh\gamma_i}\prod_{i<l}P_{il}\,
 \prod_{i<N}\frac{\dd\theta_i}{2\pi},\qquad
 P_{il}=\frac{\sin^2((\theta_i-\theta_l)/2)}
                {\sinh^2((\gamma_i+\gamma_l)/2)}.
\end{equation}
Here $\cosh\gamma_i={E}-\cos\theta_i$, $z_i=e^{-\gamma_i}$
is the interior root, and $Z=\prod_i z_i$.
Fix ${E}_*=2\cos\vartheta_*\in(0,2)$, and approach it with
\begin{equation}\label{eq:iso-parameter-approach}
 \Im{E}_\eps\asymp\eps>0,
 \qquad \Re{E}_\eps-{E}_*=O(\eps^2).
\end{equation}
All constants below may depend on this point and on fixed derivative
orders, but never on $N$ or $\eps$.  The radial exterior approach in
the physical coordinate satisfies \eqref{eq:iso-parameter-approach}.

\begin{theorem}[Complete isotropic higher tail]\label{iso:tail}
Let $n_*\ge2$, and put $k=(n_*^2-1)/2$ when $n_*$ is odd and
$k=n_*^2/2-1$ when $n_*$ is even.  For either sign $\sigma$,
\begin{equation}\label{eq:iso-complete-tail}
 \sum_{\substack{N\ge n_*+2\\N\equiv n_*\ (2)}}
  |\partial_{{E}}^j\mathcal F_N^\sigma({E}_\eps)|
 \le C_j(1+|\log\eps|),\qquad 0\le j\le k.
\end{equation}
The estimate includes every higher particle order.  Under any fixed
holomorphic change of the exterior parameter, it also holds for all
parameter derivatives through order $k$.
\end{theorem}

\subsubsection{A deformation of the complete constrained torus}

Set $b=\arccos({E}_*-1)$; then $\vartheta_*<b<2\vartheta_*$.
Choose a smooth periodic $p$, supported in $(0,\pi)$, with
$0\le p\le1$ and $p=1$ on $[\delta,\pi-\delta]$; put
$m(\theta)=p(-\theta)$.  The endpoint width $\delta$ is small,
so both $\pm b$ lie in the corresponding plateaus.  With a small
fixed $\tau>0$, define
\begin{equation}\label{eq:iso-balanced-deformation}
 P=\sum_i p(\theta_i),\quad Q=\sum_i m(\theta_i),\quad
 h_i=\frac\tau N(Pm(\theta_i)-Qp(\theta_i)),\quad
 y_i(\xi)=e^{i\theta_i+\xi h_i},\quad 0\le\xi\le1.
\end{equation}
The sum of the $h_i$ is zero, so the constrained product of the $y_i$
is unchanged.  For $y=e^{h+i\theta}$,
\[
 \Im\left({E}-\frac{y+y^{-1}}2\right)
       =\Im{E}-\sinh h\sin\theta.
\]
The motion is inward on the upper semicircle and outward on the lower
one.  Its imaginary part therefore remains positive throughout the
homotopy.  All roots stay on the interior sheet, all one-body and pair
denominators stay nonzero, and $|Z|<1$.  The complete holomorphic
integrand can consequently be transported over this closed constrained
torus.  No partition is imposed before the homotopy.

Write $w_i=\theta_i-ih_i$ at its endpoint, and $M=I-iDh$.
Because $\boldsymbol1^TM=\boldsymbol1^T$, the full determinant of $M$
equals its tangent determinant on the constrained hyperplane.  Moreover
\[
 Dh=\frac\tau N\operatorname{diag}(Pm'_i-Qp'_i)
       +\frac\tau N(m\otimes p'-p\otimes m').
\]
Each summand has operator norm $O(\tau)$, independently of $N$.
Thus $M^{-1}$ is uniformly bounded, $|\det M|\le C^N$, and its
fixed angular jets cost at most $C_j^NN^{C_j}$, by the determinant
and inverse-matrix differentiation rules.

Let $\eta=0$ on $(-\infty,1/4]$ and $\eta=1$ on $[1/2,\infty)$,
with $0\le\eta\le1$.  On the final contour use
\[
 \omega_F=\eta(P)\eta(Q),\quad
 \omega_L=\tfrac12(1-\eta(P))(1+\eta(Q)),\quad
 \omega_U=\tfrac12(1-\eta(Q))(1+\eta(P)).
\]
They sum to one.  Reflection of all angles exchanges $P,Q$, sends
$w_i$ to $-w_i$, and preserves every complete pair and one-body factor.
It also preserves either signed constraint.  Consequently
\begin{equation}\label{eq:iso-exact-split}
 \mathcal F_N^\sigma=F_N+2L_N.
\end{equation}
The support of $L_N$ has $P\le1/2$ and contains only the lower
semicircle and short upper endpoint slivers.  Lift the latter near
$\pi$ across $-\pi$; the rational pair factors and the constraint are
periodic.  Its angles then lie in $(-\pi-\delta,\delta)$.

\subsubsection{The mixed branch populations}

On the support of $F_N$, $P,Q\ge1/4$.  Since the bump supports stay
away from $0,\pi$, the imaginary-part identity gives
\begin{equation}\label{eq:iso-mixed-gap}
 -\log|Z|\ge c\left(N\eps+\frac{\tau PQ}{N}\right)\ge c/N.
\end{equation}
Every true lower branch has normal displacement $\tau P/N$, and every
true upper branch has displacement $\tau Q/N$.  The square-root form
of $\gamma$ at $W=1$ then gives a branch distance at least $c/N$ in
the $W$ plane and $|\partial_{{E}}\log z_i|\le C\sqrt N$.
These statements persist on a common disk
$|{E}-{E}_\eps|\le c/N^3$: the variation of $\log Z$
is at most $CN^{3/2}c/N^3$, smaller than half the gap in
\eqref{eq:iso-mixed-gap}.  Constants can be decreased once to cover
all $N\ge2$.

Here we use the scalar matching identity.
Put $A(u,v)=(u-v)/(1-uv)$.  Subtracting the two dispersion equations
gives the exact identity $P_{il}=A(z_i,z_l)A(y_i,y_l)$.
The scalar Schur--Pfaffian identity expands the product of the
$A(y_i,y_l)$ as a sum over matchings when $N$ is even.  When $N$ is
odd, append the constant ghost entry obtained by setting its variable
to $-1$; its entries are the removable constant one.  In each term
combine a matched angular factor with its root factor to obtain the
complete $P_{il}$.  Every remaining root factor has modulus at most
one on the actual contour, since
\begin{equation}\label{eq:iso-halfdisk-identity}
 |1-uv|^2-|u-v|^2
   =(1-|u|^2)(1-|v|^2)+4\Im u\Im v\ge0
\end{equation}
in the closed lower half disk.  This is the same scalar matching
identity as in Lemma~\ref{ana:matching}, with no strict contraction
required here.

Let $B_0$ count coordinates in fixed neighborhoods of the two true
branches.  Each contributes one to $P+Q$, so $B_0\le P+Q$.
Same-branch complete pairs are bounded: the even inverse in the
branch coordinate cancels $\gamma_i+\gamma_l$, leaving
$\gamma_i-\gamma_l$.  All pairs away from opposite branches are also
bounded by compactness.  Opposite-branch coordinates lie in a common
sector of width less than $\pi$, and hence
\[
 |1-z_i z_l|\ge c\sqrt{(P+Q)/N},\qquad
 |P_{il}|\le CN/(P+Q)\le CN/B_0.
\]
At most $B_0/2$ matched pairs are exceptional.  Since
$B_0\log(N/B_0)\le N/e$, their total extra cost is $C^N$.
There are $(N-1)!!$ matchings in the even case and $N!!$ in the odd
case.  Thus both parities satisfy
\begin{equation}\label{eq:iso-mixed-pairs}
 \left|\prod_{i<l}P_{il}\right|\le C^N N^{(N+1)/2}.
\end{equation}

This bound holds on the protected complex disk as well.  Indeed
\[
 \partial_{{E}}A(u,v)
 =\frac{u_{{E}}(1-v^2)+v_{{E}}(u^2-1)}{(1-uv)^2}.
\]
Near the branches, $|1-uv|\asymp|\gamma_i|+|\gamma_l|$,
$|u_{{E}}|\le C/|\gamma_i|$, and
$|\gamma_i|,|\gamma_l|\ge cN^{-1/2}$.  The displayed derivative
is $O(N)$; outside these neighborhoods it is no larger.
Each unmatched root entry therefore grows by at most $O(N^{-2})$
on the $c/N^3$ disk.  Their $O(N^2)$-fold product costs only a
constant.  The canceled matched pairs retain their previous bounds
there by the same sector estimates.

The one-body factors have the independent majorant
\[
 f(\theta)=C\left(1+
       \operatorname{dist}(\theta,\{b,-b\})^{-1/2}\right).
\]
Near a branch $\theta_b\in\{b,-b\}$, use the imaginary part of
$W-1$ when $|\theta-\theta_b|\le C(\eps+|h|)$, and its linear
real part outside that region.  The small
complex disk is absorbed by the branch displacement.  In addition,
the uniform branch distance $c/N$ gives
$|\sinh\gamma_i|^{-1}\le C\sqrt N$ for every coordinate throughout
the mixed support and its protected disk; outside the true branch
neighborhoods the factor is bounded by a constant.  Eliminate the
fixed $N$th angle with the constraint, bound its one-body factor by
$C\sqrt N$, and use the independent $L^1$ majorants for the others.
This gives, for every $N\ge2$,
\[
 \int_{\mathbb H_N^\sigma}\omega_F
       \prod_i|\sinh\gamma_i|^{-1}\,d\mu_N
       \le C\sqrt N\,\|f\|_1^{N-1}\le C^N\sqrt N .
\]
The estimate remains valid when the harmless cutoff is bounded by one
and the remaining independent integrations are enlarged to full circles.
The product kernel costs
$CN$ by \eqref{eq:iso-mixed-gap}, and the angular Jacobian costs
$C^N$.  Cauchy's inequality on the protected disk and $N!\ge(N/e)^N$
give
\begin{equation}\label{eq:iso-mixed-jets}
 |\partial_{{E}}^j F_N({E}_\eps)|
       \le C_j^N N^{-N/2+C_j},\qquad N\ge2.
\end{equation}
For each fixed $N=2,3$, the displaced branches and product gap instead
give a bounded analytic integrand on the compact mixed support.
Every fixed derivative of these pieces is bounded.

\subsubsection{The regular inverse and a cell Gaussian bound}

Write $z=e^{-i\varphi}$, $\gamma=i\varphi$, and set $v=\varphi+w$.

\begin{lemma}[One-sided inverse]\label{lem:iso-inverse}
There is a fixed complex neighborhood of the limiting one-sided
curve on which $w=U({E},v)$ is analytic with bounded fixed
joint jets.  On a fixed convex neighborhood $\mathcal D_R$ of the
right $v$-interval, including a short negative branch segment,
$\Re U_{vv}\ge c>0$.  All the actual one-sided images, including
the endpoint slivers, lie in these neighborhoods when the fixed
geometric choices are small.
\end{lemma}
\begin{proof}
On $[-b,0]$, $g(\theta)=\arccos({E}_*-\cos\theta)$ has
$g'=-\sin\theta/\sin g\ge0$, and $g'(0)=0$.
Thus $v=g(\theta)+\theta$ is strictly increasing, also on a short
upper sliver.  At the branch,
$w=-\arccos({E}-\cos\varphi)$ is even analytic in
$\varphi$, and $\dd(w+\varphi)/\dd\varphi=1$ at zero.
The explicit right inverse is
\begin{equation}\label{eq:iso-inverse-explicit}
 U({E},v)=\frac v2-
       \arccos\frac{{E}}{2\cos(v/2)}.
\end{equation}
Its limiting real interval is strictly inside $(-2\vartheta_*,2\vartheta_*)$,
because $b<2\vartheta_*$.  With $x=v/2$ and
$D=4\cos^2x-{E}_*^2>0$, differentiation gives
\[
 U_{vv}({E}_*,v)
 =\frac{{E}_*}{4\cos^2x\sqrt D}
      +\frac{{E}_*\sin^2x}{D^{3/2}}>0.
\]
Compactness preserves a positive real part on a thin convex rectangle.
On the left arc $\varphi=-i\gamma$, $\gamma>0$; the real part of $v$
is $\theta$, and $1+\varphi_w$ cannot vanish because $\varphi_w$
is purely imaginary.  The left and right embedded arcs meet in the
regular branch chart.  Their local analytic inverses consequently
patch on a sufficiently small neighborhood of the compact arc.
Short endpoint extensions preserve injectivity.  Square-root
continuity near the branch, and ordinary continuity elsewhere, then
put every sufficiently small deformed image inside the chosen domains.
\end{proof}

Put $\phi({E},v)=v-U({E},v)$.  With the square-root
branch fixed by $\sqrt{4\cos^2(v/2)-{E}^2}=\sin\phi-\sin U$,
implicit differentiation gives
\begin{equation}\label{eq:iso-regular-factors}\begin{gathered}
 U_v=\frac{\sin\phi}{\sqrt{4\cos^2(v/2)-{E}^2}},\quad
 U_{{E}}=\frac1{\sqrt{4\cos^2(v/2)-{E}^2}},\\
 a({E},v):=\frac{U_v}{\sinh(i\phi)}
     =\frac{-i}{\sqrt{4\cos^2(v/2)-{E}^2}}.
\end{gathered}\end{equation}
The denominator is separated on the fixed inverse neighborhoods.
The complete pair is analytic there, including at the double branch:
the dispersion identity also writes it locally as
\[
 P(v_i,v_l)=-\frac{\sin^2((\phi_i-\phi_l)/2)}
                         {\sin^2((U_i+U_l)/2)},
\]
whose denominator is nonzero near $U_i=U_l=-b$.
Away from that double branch the original pair denominator is
separated.  The diagonal zero therefore gives the global regular
factorization
\begin{equation}\label{eq:iso-pair-zero}
 P(v_i,v_l)=(v_i-v_l)^2 A({E},v_i,v_l),
\end{equation}
where $A$ and its fixed joint jets are bounded on a smaller
neighborhood of the one-sided arc squared.

\begin{lemma}[Cell contraction and retained collision factors]
\label{lem:iso-cell-gaussian}
On a fixed common complex neighborhood of the one-sided arc, for
every fixed $\ell$ the quotient
\[
 \frac{\prod_{i<l\le N}P_{il}}
      {\prod_{1\le i<l\le\ell}(v_i-v_l)^2}
\]
is analytic, and each of its fixed joint jets is bounded by
$C_{r,\ell}\exp(-\kappa N^2+C_{r,\ell}N)$ for $N\ge\ell$.
The same bound holds after multiplying by $\prod_i a({E},v_i)$.
It applies on the enlarged right rectangle with arbitrary actual
occupancy fractions in $[0,1]$.
\end{lemma}
\begin{proof}
On the limiting arc $-\pi\le\theta\le0$, both $y=e^{i\theta}$ and
$z$ are in the closed lower half disk: $z$ is on its unit semicircle
for $-b\le\theta\le0$ and on $(0,1]$ on the remaining segment.
Equations \eqref{eq:iso-halfdisk-identity} and $P=A(z_i,z_l)A(y_i,y_l)$
give $|P|\le1$.  Apparent endpoint quotients are filled using the
analytic complete pair.  Since $P(v,v)=0$, finitely many cells,
say $K$, cover the compact regular arc so that same-cell pairs have
$|P|\le\eta<1$.  Choose a small common complex neighborhood in
$({E},v)$: same-cell pairs then have modulus at most
$\eta'<1$, and all pairs have modulus at most $M=1+o(1)$.
For populations $N_1,\ldots,N_K$,
\[
 \sum_a\binom{N_a}{2}\ge\tfrac12(N^2/K-N),\qquad
 \prod_{i<l}|P_{il}|
 \le M^{\binom N2}(\eta'/M)^{(N^2/K-N)/2}.
\]
The neighborhood can be fixed small enough that the coefficient of
$N^2$ in its logarithm is negative.  This proves
$\exp(-\kappa N^2+CN)$, including on slightly larger fixed polydisks.

Factor the internal $\ell$ pairs by \eqref{eq:iso-pair-zero}.
They cost a fixed constant, their cross pairs with the remaining
coordinates cost $C^{\ell(N-\ell)}$, and the remaining complete
pair product retains its Gaussian bound.  Cauchy's inequality on the
fixed polydisks gives all asserted joint jets.  The bounded one-body
product adds only $C^N$.

For completeness, this common neighborhood contains the domains used
after absolute integration enlarges the right support.  First fix the
cells and inverse-domain margins on the undeformed compact arc.  Choose
the endpoint slivers and the negative branch segment inside those
margins.  In the full right angular rectangle, every actual
$\rho=P/N$ and $\upsilon=Q/N$ is in $[0,1]$, and
$|h_i|\le\tau$.  Away from the branch the change in $v_i$ is
$O(\eps+\tau)$; near it the uniform square-root bound makes the
change $O(\sqrt{\eps+\tau})$.  Choosing $\tau$ and then $\eps_0$
small puts all such images in the same fixed neighborhood.  No
one-sided occupancy restriction is used for this extension.
\end{proof}

\subsubsection{The exact current and its joint jets}

For $\chi_{{E}}(\theta)=(v_i)$ define the scalar
\[
 J=\det D_\theta\chi_{{E}}
    =\det M\prod_i U_{v,i}^{-1}.
\]
The physical density equals $Jf$, where
\begin{equation}\label{eq:iso-regular-density}
 f=\frac1{N!(2\pi)^{N-1}}
       \frac{1+\sigma Z}{1-\sigma Z}
       \prod_i a({E},v_i)\prod_{i<l}P_{il}.
\end{equation}
This is an equality of scalar coefficients on the original constrained
angles:
\[J\prod a=(\det M)\prod(\sinh\gamma_i)^{-1}.\]
No intrinsic complex measure on a new constrained $v$ hyperplane is
asserted.

For a regular-coordinate field $B$ put
\begin{equation}\label{eq:iso-angular-field}
 V=-M^{-1}(U_{{E}}+U_vB).
\end{equation}
If $\sum_iU_{v,i}B_i=-\sum_iU_{{E},i}$, then $\sum_iV_i=0$.
If also $\sum_iB_i=0$, the operator
$\partial_{{E}}-V\cdot\nabla_\theta$ freezes $\sum_i\varphi_i$.
Indeed $\sum\varphi_i=\sum v_i-\sum w_i$, and $\sum w_i=\sum\theta_i$.
The exact scalar current identity is
\begin{equation}\label{eq:iso-current}
 \partial_{{E}}(Jb)-\operatorname{div}_\theta(VJb)
  =J\left(E_{{E}}b+\sum_i B_iE_i b
                           +\Bigl(\sum_iE_iB_i\Bigr)b\right).
\end{equation}
Here $b$ may already be a parameter-dependent transported amplitude.
The commuting regular frame is
\[
 E_i=\sum_l(M^{-1})_{li}U_{v,i}\partial_{\theta_l},\qquad
 E_{{E}}=\partial_{{E}}
                  +\sum_l(M^{-1}U_{{E}})_l\partial_{\theta_l}.
\]
It satisfies $E_i v_l=\delta_{il}$ and $E_{{E}}v_l=0$.
Every commutator is a linear combination of the angular derivatives,
annihilates all $v_l$, and hence vanishes.  Thus this frame applies
also to smooth real cutoffs without analytically continuing them.
To verify \eqref{eq:iso-current}, differentiating the coordinate
Jacobian and using Jacobi's determinant formula gives
$J\operatorname{div}_vB=\partial_{{E}}J-\operatorname{div}_\theta(JV)$;
the product rule gives the displayed identity.  Tangency makes its
ambient divergence exactly the constrained one, since
\[
 \sum_{i<N}(\partial_i-\partial_N)V_i=\sum_i\partial_iV_i.
\]
Integration by parts is applied to the real and imaginary parts of
these complex currents.

Choose one smooth right cutoff $r$, equal to one for
$\theta+b\ge-\delta_B/2$ and zero for $\theta+b\le-\delta_B$,
with its periodic return in the excluded upper plateau.  Decompose
$1=\prod_i r_i+(1-\prod_i r_i)$ on the one-sided term and telescope
the second part.  Each of its at most $N$ terms has a named compact
left anchor $a$, also on all cutoff jets.  It has
$|z_a|\le q_0<1$ and $|U_{v,a}|\ge c$.  Use
\begin{equation}\label{eq:iso-anchor-field}
 B_a=-\frac{\sum_iU_{{E},i}}{U_{v,a}},\qquad B_i=0\ (i\ne a).
\end{equation}
This gives tangency; phase freezing is unnecessary because
$|1-\sigma Z|\ge1-q_0$.  All fixed regular-frame jets of $B,V$ and
the cutoffs have polynomial $N$ bounds.  The cell Gaussian, the
bounded regular one-body factors, and \eqref{eq:iso-current} then
bound every fixed derivative of this whole piece by
\begin{equation}\label{eq:iso-anchor-bound}
 \frac{C_jN^{C_j}}{N!}\exp(-\kappa N^2+C_jN).
\end{equation}
In fact, eliminate the compact anchor angle with the constraint.
Its Jacobian factor is bounded, and every remaining factor has the
independent integrable majorant
$C(1+\operatorname{dist}(\theta,-b)^{-1/2})$.
Their integrals cost only $C^N$.  The product kernel and its
transported derivatives have polynomial $N$ bounds on this anchored
support.  This argument applies also at each fixed $N=2,3$.

We now work on the full right part, with all angles in
$I=[-b-\delta_B,\delta]$.  Set $u_i=\theta_i+b$ and, on a diameter
shell $r_0\le d:=\max\theta_i-\min\theta_i\le2r_0$, put
\begin{equation}\label{eq:iso-scales}
 \mu=N^{-1}\sum_i u_i,\qquad
 T=|\mu|+r_0+\eps,\qquad \Delta=r_0/\sqrt T.
\end{equation}
These are pointwise envelopes; no derivative of a mean, diameter,
or scale is ever taken.

The branch normal form, for a common outward displacement $h\ge0$,
is
\begin{equation}\label{eq:iso-branch-normal}
 \begin{gathered}
 1-W=\sin b\,u-i\beta+R,\qquad
 c(\eps+h)\le\beta\le C(\eps+h),\\
 R\in\mathbb R,\qquad |R|\le C(u^2+\eps^2+h^2),\qquad
 \varphi=\sqrt{2(1-W)}(1+O(1-W)).
 \end{gathered}
\end{equation}
Indeed $W={E}-\cosh(h+i(-b+u))$ and
$\beta=\Im{E}+\sinh h\sin(b-u)$; its real part has the
displayed Taylor expansion.  The physical square root lies in a
small enlargement of the fourth quadrant.  In particular
\[
 \begin{gathered}
 |1+\varphi_w(w_i)|\le C(|u_i|+\eps)^{-1/2},\quad
 -\Im\varphi_i\ge c\sqrt{|u_i|+\eps}\quad(u_i\le0),\\
 \int_{\mu-2r_0}^{\mu+2r_0}(|u|+\eps)^{-1/2}\dd u
       \le Cr_0/\sqrt T.
 \end{gathered}
\]
The interval bound follows by comparability if $|\mu|\ge4r_0$;
otherwise enlarge to $[-6r_0,6r_0]$ and rationalize
$\sqrt{6r_0+\eps}-\sqrt\eps$.
Joining two angles along their one-variable curve with the actual
$\rho=P/N,\upsilon=Q/N$ held fixed gives
\begin{equation}\label{eq:iso-coordinate-scales}
 |v_i-v_l|\le C\Delta,\qquad |U_{v,i}|\le C\sqrt T.
\end{equation}
The joining derivative has the additional bounded factor
$1-i\tau(\rho m'-\upsilon p')$.  Near the lower branch its outward
shift is the common $\tau\rho$, so the preceding interval bound
applies; elsewhere it is regular.  For the second bound, if $P=0$
then every actual $h_i=0$, and the branch form and
\eqref{eq:iso-regular-factors} apply directly.  If $P>0$, some
actual angle is positive, so $\max u_i\ge b$ and $T\ge b/2$;
boundedness of $U_v$ proves the same result.  These arguments allow
all actual occupancy fractions in $[0,1]$.


Use the filled analytic divided differences
\[
 \mathcal U_{il}=\int_0^1U_v(v_l+t(v_i-v_l))\,dt,\qquad
 \mathcal E_{il}=\int_0^1U_{vv}(v_l+t(v_i-v_l))\,dt.
\]
The convex inverse domain gives $\Re\mathcal E_{il}\ge c$.
Define the analytic field
\begin{equation}\label{eq:iso-covariance-field}
 \begin{gathered}
 D_U=\sum_{i<l}(U_i-U_l)^2,\qquad
 H_{il}=\frac{(v_i-v_l)\mathcal U_{il}^2}{\mathcal E_{il}},
 \qquad G_i=\sum_lH_{il},\\
 A=-\sum_iU_{{E},i},\qquad B_i=AG_i/D_U.
 \end{gathered}
\end{equation}
There is no pair pole: $H_{il}$ is analytic also at $v_i=v_l$.
Away from equality it equals $(U_i-U_l)^2/(U_{v,i}-U_{v,l})$.
Antisymmetry and pairing give
\[
 \sum_iG_i=0,\qquad \sum_iU_{v,i}G_i=D_U,\qquad
 \sum_iB_i=0,\qquad \sum_iU_{v,i}B_i=-\sum_iU_{{E},i}.
\]
Thus the field has both exact transport identities.

\begin{lemma}[Protected analytic variance and joint jets]
\label{lem:iso-covariance-jets}
On each actual right shell,
\begin{equation}\label{eq:iso-covariance-lower}
 \Re D_U\ge c\sum_{i<l}(\theta_i-\theta_l)^2
                   \ge cT\Delta^2>0.
\end{equation}
For every fixed total derivative order $s$,
\begin{equation}\label{eq:iso-field-jets}
 |\mathcal D^s B_i|+|\mathcal D^s V_i|
       \le C_sN^{C_s}\Delta^{-1-s}.
\end{equation}
The field is regular at every partial diagonal.
\end{lemma}
\begin{proof}
On the actual contour $U_i=w_i=\theta_i-ih_i$, and the real
balanced shifts satisfy
$|h_i-h_l|\le C\tau|\theta_i-\theta_l|$ at fixed actual occupancies.
Hence
$\Re(U_i-U_l)^2=(\theta_i-\theta_l)^2-(h_i-h_l)^2
 \ge(1-C^2\tau^2)(\theta_i-\theta_l)^2$.
Small fixed $\tau$ proves the denominator bound.

Freeze the actual values of $T,\Delta$ and consider the polydisk
\[
 |\delta{E}|,\ |\delta v_i|\le c_0\Delta/N^2.
\]
Choose $c_0$ small so it lies in the fixed inverse neighborhoods.
Throughout this disk, $|v_i-v_l|\le C\Delta$,
$|U_v|\le C\sqrt T$, and
$|U_{{E}}(v_i)-U_{{E}}(v_l)|\le C\Delta$.
Indeed the last bound is an analytic divided difference; the middle
one follows from its actual endpoint bound and bounded joint jets,
since $\Delta\le\sqrt T$.
Along the straight parameter and coordinate segments, the common
parameter change therefore moves each $U_i-U_l$ by $O(\Delta
|\delta{E}|)$, and the two coordinate changes move it by
$O(\sqrt T(|\delta v_i|+|\delta v_l|))$.
Their total change is at most $Cc_0r_0/N^2$.
Since each original difference is $O(r_0)$, summing the $N(N-1)/2$
squared differences gives
\[
 |D_U^{\rm new}-D_U|\le Cc_0r_0^2\le\tfrac12\Re D_U.
\]
Thus $|D_U^{\rm new}|\ge cr_0^2$ on this whole polydisk.
Also $|\mathcal U_{il}|\le C\sqrt T$ and
$|\mathcal E_{il}|\ge c$ there, so
$|H_{il}|\le Cr_0\sqrt T=CT\Delta$, $|A|\le CN$, and
$|B_i|\le CN^2/\Delta$.
Cauchy's inequality now gives all fixed joint jets at once:
\begin{equation}\label{eq:iso-covariance-jet-table}
 |\partial_{{E}}^a\partial_v^\beta B_i|
 \le C_sN^{2+2s}\Delta^{-1-s},\qquad s=a+|\beta|.
\end{equation}
No scale or polydisk radius has been differentiated.

The commuting-frame coefficients
$M^{-1}\operatorname{diag}U_v$ and $M^{-1}U_{{E}}$ have
polynomially bounded fixed joint jets: the real angular jets of
$M$ and analytic jets of $U$ are bounded, and differentiating
$M^{-1}$ proves the assertion inductively without an inverse $U_v$.
This also bounds all fixed smooth cutoff jets.
Formula \eqref{eq:iso-angular-field}, the product rule, and
\eqref{eq:iso-covariance-jet-table} give the same field bound for $V$.
The positive actual denominator and filled analytic $H_{il}$ show
that no partial diagonal is singular.
\end{proof}


Fix now $\ell=n_*+1$ and $D_\ell=\ell(\ell-1)$.
Factor the first $\ell$ collision variables using
Lemma~\ref{lem:iso-cell-gaussian}.  For the initial smooth weighted
regular amplitude $b_0$ after removing the product kernel,
\[
 |\mathcal D^s b_0|
 \le\frac{C_sN^{C_s}}{N!}e^{-\kappa N^2+C_sN}
                     \Delta^{D_\ell-s},\qquad s\le2k.
\]
The same statement includes parameter derivatives, which preserve the
explicit collision polynomial.  Iterate the exact current with
\begin{equation}\label{eq:iso-transport-recurrence}
 b_{j+1}=\mathcal T b_j,\qquad
 \mathcal T b=E_{{E}}b+\sum_iB_iE_i b
                              +\Bigl(\sum_iE_iB_i\Bigr)b.
\end{equation}
This recurrence differentiates every parameter-dependent descendant
coefficient and cutoff.  Because the field freezes $\sum\varphi_i$,
the product kernel is unchanged.  Induction using
\eqref{eq:iso-field-jets} gives
\begin{equation}\label{eq:iso-transported-amplitude}
 |\mathcal D^s b_j|
 \le\frac{C_jN^{C_j}}{N!}e^{-\kappa N^2+C_jN}
              \Delta^{D_\ell-2j-s},\qquad s\le2(k-j).
\end{equation}
Indeed $E_{{E}}$ loses at most one scale power, while
$B_iE_i$ and $(E_iB_i)$ each lose at most two.  Leibniz' rule gives
the same count for joint jets.  Only polynomially many index choices
occur at each fixed order.

There is no hidden boundary term.  At fixed $N,\eps>0$ puncture
the finitely many full-equality points by angular balls of radius
$r$.  All geometric Jacobians and the product kernel are then smooth;
$V=O(r^{-1})$, and the retained collision degree is $D_\ell$.
The flux after $j$ transports is bounded by a constant times
$r^{D_\ell+N-3-2j}$, which tends to zero for every $N\ge n_*+2$
and $j\le k$.  Partial diagonals require no puncture by
\eqref{eq:iso-covariance-lower}.  The original smooth weights vanish
at the outer chart boundaries; no diameter cutoff has been
differentiated.  Thus the $j$th derivative of the complete right term
is the integral of its transported density, bounded by
\begin{equation}\label{eq:iso-transported-density}
 \frac{C_jN^{C_j}}{N!}e^{-\kappa N^2+C_jN}
 \Delta^{D_\ell-2j}
 \frac{\prod_i|1+\varphi_w(w_i)|}{|1-\sigma Z|}.
\end{equation}

\subsubsection{Coarea with the actual occupancy dependence}

We prove the integral estimate needed for
\eqref{eq:iso-transported-density}; in particular the occupancies
will not be treated as constants when differentiating the phase.

First, for $G=x-ia$, $a\ge a_0>0$, and fixed
$|\vartheta|<\pi/2$, put $\nu=\Re(e^{-i\vartheta}G)$ and
$P_\vartheta=2\pi\cos\vartheta$.  If $a_0\le1$, then
\begin{equation}\label{eq:iso-oblique-kernel}
 \frac1{|1-e^{-iG}|}
 \le\frac{C_\vartheta}
 {a_0+\operatorname{dist}(\nu,P_\vartheta\mathbb Z)}.
\end{equation}
For $a\le1$, this follows from
$|1-e^{-a-ix}|\ge c(a+\operatorname{dist}(x,2\pi\mathbb Z))$
and
$|\nu-2\pi n\cos\vartheta|\le
 \cos\vartheta|x-2\pi n|+|\sin\vartheta|a$.
For $a\ge1$, the left denominator is separated and the proposed
right denominator is bounded above.  Integrating the majorant over
an interval of length $L$ costs at most
$C_\vartheta(1+L)(1+|\log a_0|)$.  For $\sigma=-1$ replace $G$
by $G-\pi$.  On the actual contours $-\log|Z|\ge cN\eps$:
write $\varphi_i=x_i-ia_i$, and use
$\Im W_i=\sin x_i\sinh a_i\le Ca_i$ on the bounded phase range.
Apply the lemma with $a_0=\min(1,cN\eps)$.  Its logarithmic cost is
at most $C(1+L)(1+|\log\eps|)$.

\begin{lemma}[Complete right-shell coarea]\label{lem:iso-coarea}
On each fixed-mean lift of the full right rectangle
$I^N\cap\mathbb H_N^\sigma$, with the actual balanced occupancies,
the integral over any diameter shell in \eqref{eq:iso-scales} obeys
\begin{equation}\label{eq:iso-shell-coarea}
 \int_{\substack{\sum u_i=N\mu\\r_0\le d\le2r_0}}
  \frac{\prod_i|1+\varphi_w(w_i)|}{|1-\sigma Z|}
       \prod_{i<N}\dd\theta_i
 \le C^NN^C\Delta^{N-3}(1+|\log\eps|).
\end{equation}
There are at most $CN$ constrained means.  The constants are uniform
in each of them; their subsequent sum costs only this polynomial
factor.
\end{lemma}
\begin{proof}
Write $m_0=\min u_i$, $M_0=\max u_i$, and first suppose
$0\le m_0<M_0$.  With other angles fixed, the constraint fixes
$m_0+M_0$.  Put $G=\sum_i\varphi_i$ and $\vartheta=3\pi/8$.
We claim
\begin{equation}\label{eq:iso-extreme-ratio}
 \Re(e^{-i\vartheta}\dd G/\dd m_0)>0,\qquad
 \frac{|(1+\varphi_w(w_{m_0}))(1+\varphi_w(w_{M_0}))|}
      {\Re(e^{-i\vartheta}\dd G/\dd m_0)}
 \le C\frac{\sqrt{M_0+\eps}}{M_0-m_0}.
\end{equation}

For its proof only, freeze the actual values $\rho=P/N,\upsilon=Q/N$
and set
$f(u)=-b+u-i\tau\{\rho m(-b+u)-\upsilon p(-b+u)\}$,
$E=\eps+\tau\rho$.  For $u\ge0$, the branch normal form and
compactness on the rest of $I$ give
\[
 |\varphi_w(f(u))|\le C(u+E)^{-1/2},\qquad
 |\varphi_{ww}(f(u))|\asymp(u+E)^{-3/2}.
\]
The inverse supplies
$-\varphi_{ww}=U_{vv}/U_v^3$.  For $u\ge0$ the argument of $U_v$
is in a small enlargement of $[-\pi/4,0]$, while $U_{vv}$ has small
argument.  Hence $-\varphi_{ww}$ has uniformly positive
$3\pi/8$ projection.  If $a=f'=1+O(\tau)$ and
$B(u)=\varphi_w(f(u))a(u)$, then
$-B'=-\varphi_{ww}(f)a^2-\varphi_w(f)a'$ retains that projection
for small $\tau$.  Consequently
\begin{equation}\label{eq:iso-fixed-occupancy-curvature}
 \begin{gathered}
 \Re\{e^{-i\vartheta}[B(m_0)-B(M_0)]\}\ge cL,
 \\
 L=\int_{m_0}^{M_0}(u+E)^{-3/2}\dd u
 =\frac{2(M_0-m_0)}{\sqrt{m_0+E}\sqrt{M_0+E}
                         (\sqrt{m_0+E}+\sqrt{M_0+E})}.
 \end{gathered}
\end{equation}

Now restore all occupancy derivatives.  Put
$B_m^{\rm occ}=\sum_lm(\theta_l)\varphi_w(w_l)$ and
$B_p^{\rm occ}=\sum_lp(\theta_l)\varphi_w(w_l)$.
Direct differentiation of \eqref{eq:iso-balanced-deformation} gives
the exact identity
\begin{equation}\label{eq:iso-occupancy-derivative}
 \partial_{\theta_i}G
 =\varphi_w(w_i)[1-i\tau(\rho m'_i-\upsilon p'_i)]
       -\frac{i\tau p'_i}{N}B_m^{\rm occ}
       +\frac{i\tau m'_i}{N}B_p^{\rm occ}.
\end{equation}
Every actual $u_l\ge m_0$, so
$(|B_m^{\rm occ}|+|B_p^{\rm occ}|)/N\le C(m_0+E)^{-1/2}$.
Taking $\partial_{\theta_{\min}}-\partial_{\theta_{\max}}$,
both bump-derivative differences are $O(M_0-m_0)$.
The entire occupancy correction is therefore at most
$C\tau(M_0-m_0)(m_0+E)^{-1/2}\le C\tau L$, by the exact
expression for $L$ and the fixed upper bound on $M_0+E$.
Small fixed $\tau$ absorbs it into the positive projection.
The numerator of \eqref{eq:iso-extreme-ratio} is at most
$C[(m_0+E)(M_0+E)]^{-1/2}$, giving ratio
$C\sqrt{M_0+E}/(M_0-m_0)$.  If $P=0$, $E=\eps$; if $P>0$,
an actual positive angle gives $M_0\ge b$, and hence
$E\le C(M_0+\eps)$.  This proves \eqref{eq:iso-extreme-ratio}
with all actual occupancies, including on the enlarged rectangle.

After transport, bound every smooth coefficient by its proved
envelope and enlarge to this full rectangle.  All its occupancy
fractions lie in $[0,1]$, the balanced displacement preserves
$\Im W_i\ge\Im{E}$, and Lemma~\ref{lem:iso-cell-gaussian}
has already placed its images in the regular domains.
Fix a constrained mean $\mu=b+(2\pi l+a_\sigma)/N$ and the two
extreme indices; at most $CN^3$ such choices occur.  Eliminate the
maximum coordinate with the constraint.  With the remaining $N-2$
coordinates fixed, conditions specifying the extrema, the shell,
and a nonnegative minimum are linear interval conditions in $m_0$,
with $M_0=A-m_0$.  Thus the fiber is one interval, and its projected
phase is injective by \eqref{eq:iso-extreme-ratio}.  Its phase range
is at most $CN$.  Equations \eqref{eq:iso-oblique-kernel} and
\eqref{eq:iso-extreme-ratio} bound this extreme-coordinate integral by
\[
 CN\frac{\sqrt T}{r_0}(1+|\log\eps|).
\]

If $m_0<0$, its branch attenuation from
\eqref{eq:iso-branch-normal} instead gives the direct bound
\[
 \frac{|(1+\varphi_w(w_{m_0}))(1+\varphi_w(w_{M_0}))|}
      {|1-\sigma Z|}
 \le\frac{C}{(|m_0|+\eps)\sqrt{|M_0|+\eps}}.
\]
With $a=|m_0|+\eps$ and $c=|M_0|+\eps$, $a+c\ge r_0$, and
\[
 \frac1{a\sqrt c}\le\frac1{\sqrt{r_0}}\left(\frac1a+\frac1c\right).
\]
Each reciprocal integrates to $C(1+|\log\eps|)$ on the bounded
slice $M_0=A-m_0$.  Since $T\ge r_0$, this gives the same extreme
pair envelope $C\sqrt T/r_0(1+|\log\eps|)$.
Only now replace the other $N-2$ Jacobian factors by the independent
majorants and integrate each over $[\mu-2r_0,\mu+2r_0]$.
Their total cost is $(Cr_0/\sqrt T)^{N-2}$.
Combining the extreme pair, the ordinary factors, and the polynomial
number of cells proves \eqref{eq:iso-shell-coarea}.
\end{proof}

\subsubsection{One sum for all higher particle orders}

Combining \eqref{eq:iso-transported-density} with
Lemma~\ref{lem:iso-coarea} bounds a right shell at a fixed mean by
\begin{equation}\label{eq:iso-one-shell}
 \frac{C_jN^{C_j}}{N!}e^{-\kappa N^2+C_jN}
 \Delta^{D_\ell+N-3-2j}(1+|\log\eps|).
\end{equation}
Its exponent is positive for every $N\ge n_*+2$ and $j\le k$:
it is at least $n_*+N-2$ in the odd case and $n_*+N-1$ in the
even case.  For fixed $\mu,\eps$,
\[
 \frac{\Delta(r_0/2)}{\Delta(r_0)}
 =\frac12\sqrt{\frac{|\mu|+r_0+\eps}{|\mu|+r_0/2+\eps}}
       \in[1/2,1/\sqrt2].
\]
The dyadic sum is therefore bounded uniformly in the mean by a geometric series.
Its largest possible scale is fixed; its power costs only $C^N$
times a fixed constant.  This is absorbed into the linear exponential
in \eqref{eq:iso-one-shell}.  Summing the at most $CN$ means adds
only a polynomial factor.  No small/large-scale division or finite
intermediate particle range is required.  The mixed, anchored, and
right pieces consequently obey, for all higher orders,
\begin{equation}\label{eq:iso-every-order-tail}
 |\partial_{{E}}^j\mathcal F_N^\sigma({E}_\eps)|
 \le C_j(1+|\log\eps|)
 \left(C_j^NN^{-N/2+C_j}
       +\frac{N^{C_j}}{N!}e^{-\kappa N^2+C_jN}\right).
\end{equation}
Both envelopes are summable.  The exact split
\eqref{eq:iso-exact-split} proves Theorem~\ref{iso:tail}.
The parameter-change assertion follows from the finite chain rule.

\begin{proposition}[Finite-order transport interface]
\label{prop:iso-finite-transport}
For each fixed $N\ge2$ and every fixed derivative order, the mixed
piece is bounded and extends holomorphically across
${E}_*$ on its deformed contour.  A smooth one-sided term with
a compact left anchor has bounded derivatives of every fixed order.
A smooth right-supported term whose angular diameter is bounded
below by a fixed $d_0>0$ has derivatives
$O(1+|\log\eps|)$ of every fixed order.  The assertions allow
arbitrary fixed smooth angular cutoffs with those supports.
\end{proposition}
\begin{proof}
For the fixed mixed order, both true branches and the product
denominator have uniform gaps on the compact support, as proved
above; the same small parameter disk works for the whole integrand.
This includes $N=2,3$ directly.  The anchored field
\eqref{eq:iso-anchor-field} and its integrable independent Jacobian
majorants prove the second assertion at any fixed derivative budget.
On a right support with $d\ge d_0$, all jets of the covariance
field are bounded and there is no collision puncture.  Iterate
\eqref{eq:iso-transport-recurrence} with that budget, and then apply
Lemma~\ref{lem:iso-coarea}; the finitely many dyadic scales bounded
away from zero cost $O(1+|\log\eps|)$.  For $N=2$ the coarea keeps
both extreme factors and has zero remaining one-body integrations,
so the same argument applies.  All differentiated cutoff factors
are included in the exact recurrence.
\end{proof}



For negative real dispersion in the odd family, the exact identity
\eqref{iso:negation}, followed by conjugation, sends the exterior
approach to the signed positive geometry. More explicitly, for
$\widehat\zeta=-\overline\zeta$ and odd $N$,
\[
 (C_N^+)^{(j)}(\zeta)
   =(-1)^{j+1}\overline{(C_N^-)^{(j)}(\widehat\zeta)}.
\]
Thus Theorem~\ref{iso:tail} controls every higher physical odd
coefficient also at the original negative-dispersion points.
The sign changes both the momentum constraint and the product kernel.

\subsection{The finite complement and the original prime family}

\begin{lemma}[All fixed derivatives of the finite complement]
\label{iso:fixed}
At every point of Lemma~\ref{iso:arithmetic}, all fixed derivatives
of each smaller coefficient of the selected parity are bounded along
$\zeta_\eps=(1+\eps)\zeta_*$. At the first order $N_0$, choose
small fixed smooth shape cutoffs at the four angular configurations
$q_i=\alpha,-\alpha,\beta,-\beta$ for every $i$. The complete
coefficient minus these four local integrals has bounded derivatives
of every fixed order on the same approach.
\end{lemma}
\begin{proof}
First suppose $E_*>0$. For the finitely many relevant orders choose
the branch width smaller than one quarter of the distance from the
branch to any constrained equality mean. This distance is positive
by the endpoint exclusion in Lemma~\ref{iso:arithmetic}.
Use the whole-torus deformation and exact reflection split proved
above. The mixed term has a fixed product gap and displaced branches
at every fixed order $N\ge2$. Its derivatives are bounded. The
compact-left-anchor terms have the same property by
Proposition~\ref{prop:iso-finite-transport}. The one-particle term is
$1/(E-2)$, and is regular here.

At equality on the right support, the balanced displacement vanishes:
$P=Np(q_0)$, $Q=Nm(q_0)$ give $Pm(q_0)-Qp(q_0)=0$.
All roots and complete pairs are regular at these finitely many
points. A vanishing product denominator requires
\[
 E_* =\cos q_0+\cos\varphi_0,\qquad
 Nq_0,N\varphi_0\in2\pi\Z.
\]
The arithmetic excludes such a point at every smaller order and
leaves exactly $q_0=-\alpha,-\beta$ in the first lower chart.
Their reflected configurations give the two upper charts.
Every other equality neighborhood has a fixed product gap.
Choose the critical cutoffs inside the lower plateau, where all
$m_i=1,p_i=0$: the deformation is zero there and the lower weight
and the right cutoff are one.

After these neighborhoods are removed, the remaining right support
has angular diameter at least a fixed $d_0>0$. The finite-order
transport interface gives $O(1+|\log\eps|)$ for each of its fixed
derivatives, including all cutoff derivatives. This covers every
remaining phase zero. Thus each complete lower coefficient, and the
complete first-order remainder, has that bound at every fixed order.
These remainders are holomorphic on the one-sided parameter sector:
their contours and smooth weights are independent of $E$.
For each fixed $j$, integration of the next derivative gives
\[
 \partial_E^j R(E(\zeta_\eps))
 =\partial_E^j R(E(\zeta_{\eps_1}))
  -\int_\eps^{\eps_1}\frac{d E(\zeta_u)}{du}
                  \partial_E^{j+1}R(E(\zeta_u))\,du.
\]
The logarithmic bound is integrable and $dE(\zeta_u)/du$ is bounded.
Every stated derivative is therefore bounded. The finite chain rule
gives the conclusion in $\zeta$.

For $E_*<0$ in the odd family, apply conjugation and
\eqref{iso:negation} to obtain the signed positive geometry at
$\widehat E_*=-E_*>0$. The transformed constraint and product
kernel both have sign $-1$. Equality resonance is carried exactly
to the original resonance by this invertible transformation; the
endpoint exclusion is also unchanged. The preceding argument uses
either sign, so it applies to this complete transformed remainder.
Transforming back proves the assertion. It does not assert evenness
of the odd physical response.
\end{proof}

\begin{proposition}[The complete first amplitude]\label{iso:first}
Let $N=N_0$, $a_N=(N^2-1)/2$, $k=\lfloor a_N\rfloor$ and
$\nu=k+1-a_N\in\{1/2,1\}$. With the original angles and exterior
parameter of Lemma~\ref{iso:arithmetic},
\[
 C_N^{(k)}(\zeta_\eps)
       =4L_*\eps^{-\nu}+o(\eps^{-\nu}),\qquad L_*\ne0.
\]
All derivatives below $k$ are bounded for even $N$. For odd $N$
they are bounded through $k-2$ and at most logarithmic at $k-1$.
In particular these are the same four full-coefficient chart
contributions for the distinct-angle family in both signs of $E_*$.
\end{proposition}
\begin{proof}
In the lower chart put $q_i=-\alpha+u_i$, $\sum_i u_i=0$,
$Q_2(u)=\sum_i u_i^2$ and $\Delta_N(u)=\prod_{i<j}(u_i-u_j)$.
The physical branch has $\gamma_i=i\varphi_i$ and
$\varphi_i=\beta$ at the center. Implicit differentiation of
$\cos\varphi=E-\cos q$ gives
\[
 \varphi_q=\frac{\sin\alpha}{\sin\beta},\qquad
 \varphi_{qq}=-\frac{H_{\alpha\beta}}{\sin^3\beta}=-2d,
 \quad H_{\alpha\beta}=E_*(1-\cos\alpha\cos\beta),
 \quad d=\frac{H_{\alpha\beta}}{2\sin^3\beta}\ne0.
\]
Put $Q=2N\sin\theta/\sin\beta>0$. The literal angular constraint
cancels the linear shape term and its mixed parameter derivative.
Since $N\beta\in2\pi\Z$, it follows that
\begin{align}
 1-Z&=Q\eps-id Q_2(u)
   +O\left(\eps^2+\eps Q_2(u)+Q_2(u)^2+\sum_i|u_i|^3\right),
                                                     \label{iso:local-q}\\
 \partial_\zeta(1-Z)(\zeta_*,0)&=Qe^{-i\theta}.
                                                     \label{iso:local-derivative}
\end{align}
The leading terms have perpendicular phases. On a fixed sufficiently
small real shape ball and for small $\eps>0$ this gives
$|1-Z|\ge c(\eps+Q_2(u))$.

The local integral is exactly
\[
 I_N(\zeta)=\int_{\R^{N-1}}
       \frac{\chi(u)\Delta_N(u)^2 A_N(\zeta,u)}{1-Z}\,
                          du_1\cdots du_{N-1},
 \qquad u_N=-\sum_{i<N}u_i,
\]
with analytic amplitude and
\[
 A_N(\zeta_*,0)=
 \frac{2i^{-N^2}}{N!(2\pi)^{N-1}2^{N(N-1)}\sin^{N^2}\beta}.
\]
Here every pair numerator contributes $(u_i-u_j)^2/4$, its
denominator contributes $(i\sin\beta)^2$, the $N$ one-body
factors contribute $(i\sin\beta)^{-N}$, and $1+Z$ contributes
two. Thus all Haar and factorial normalizations are included.

Define the positive angular constant
\[
 c_N=\frac1{\sqrt N}
   \int_{\substack{\sum v_i=0\\\sum v_i^2=1}}
                 \Delta_N(v)^2\,d\sigma(v)>0.
\]
The factor $1/\sqrt N$ is the Gram determinant correction for
the independent coordinates $v_1,\ldots,v_{N-1}$ on this
hyperplane. Polar coordinates, followed by $x=r^2$, give
\begin{align*}
 \int_{\R^{N-1}}\frac{\Delta_N(v)^2}{(Q-id Q_2(v))^{k+1}}\,dv
 &=\frac{c_N}{2}\int_0^\infty
                \frac{x^{a_N-1}}{(Q-id x)^{k+1}}\,dx\\
 &=\frac{c_N}{2}Q^{a_N-k-1}(-id)^{-a_N}
                         \operatorname B(a_N,k+1-a_N).
\end{align*}
For a positive real coefficient in place of $-id$, substitution
gives the beta integral. Continue that coefficient through the
appropriate quadrant to $-id$. On each such closed quadrant of
fixed positive modulus, $|Q+cx|\ge c_0(1+x)$, so dominated
convergence proves the identity also at the endpoint. The radial
tails are $r^{-2}$ for even $N$ and $r^{-3}$ for odd $N$.

In the $k$th derivative of the local integral the unique highest
pole has numerator
$(-1)^k k!(\partial_\zeta(1-Z))^k\chi\Delta_N^2 A_N$.
Under $u=\sqrt\eps v$ its scale is $\eps^{-\nu}$, and its
scaled integrand has the integrable majorant
$C\Delta_N(v)^2/(1+Q_2(v))^{k+1}$. Dominated convergence
therefore gives its leading coefficient by the preceding period.
Every other differentiated term has a lower pole order and is
bounded in the even case and at most logarithmic in the odd case.
The same radial estimates prove the lower-derivative statements.
The one-chart coefficient simplifies to
\begin{equation}\label{iso:amplitude}
 L_*=
 \frac{i^{-N^2}(-1)^k k!\,c_N\operatorname B(a_N,k+1-a_N)
       2^{a_N}(2N\sin\theta)^{a_N-1}e^{-ik\theta}}
 {N!(2\pi)^{N-1}2^{N(N-1)}(-iH_{\alpha\beta})^{a_N}}.
\end{equation}
All sine powers cancel, since $-N^2-(a_N-1)+3a_N=0$.
For even $N$, $H_{\alpha\beta}>0$ and the last power uses
argument $-\pi/2$. For odd $N$, $a_N=k$ is divisible by four,
so the same power is $|H_{\alpha\beta}|^k$ for either sign.
Every factor is nonzero and the expression is symmetric in
$\alpha,\beta$. Their two lower charts therefore have exactly
the same complex coefficient. Reflection of every angle preserves
the complete constrained density and gives the two equal upper
charts. Lemma~\ref{iso:fixed} disposes of their complete complement,
proving the factor four and the claimed full-coefficient asymptotic.
\end{proof}

\subsection{Dense local obstructions and the exact physical identities}

The full-site form-factor identification gives, on the respective real
physical intervals,
\begin{align}
 \mathcal X_L(s)&=(1-s^{-4})^{1/4}
                  \sum_{N\ge2,\ N\ {\rm even}}C_N(s),
                  &&s>1,\label{iso:physical-L}\\
 \mathcal X_H(s)&=\frac{(1-s^4)^{1/4}}s
                  \sum_{N\ge1,\ N\ {\rm odd}}C_N(1/s),
                  &&0<s<1.\label{iso:physical-H}
\end{align}
The low-temperature formula has the plus-phase subtraction already
performed: its omitted zero-particle term is $M_+^2$, and its onsite
contribution is $1-M_+^2$. The origin is already included in each full-site coefficient; no
additional constant is appended to either complete sum.
The high-temperature onsite sum is one, and $C_1=1/(E-2)$ checks the
normalization. The branches in \eqref{iso:physical-L}--\eqref{iso:physical-H}
are positive on their stated real intervals.

\begin{proposition}[Dense obstructions for the physical isotropic sums]
\label{iso:local-obstruction}
The full normalized even exterior sum in \eqref{iso:physical-L}
and the full normalized odd interior sum in \eqref{iso:physical-H}
have no local holomorphic extension at a dense set of points of
\(|s|=1\). At their selected points the nonzero leading derivatives
are those of Proposition~\ref{iso:first}, multiplied by the physical
nonvanishing factors. These assertions use the full-tail
bounds of Theorem~\ref{iso:tail}.
\end{proposition}
\begin{proof}
At every selected point, Proposition~\ref{iso:first}, the bounded lower
orders, and Theorem~\ref{iso:tail} give a nonzero leading derivative of
the complete physical sum. All lower derivatives are smaller than that
leading scale, so multiplication by its local analytic nonzero phase
prefactor preserves it by the product rule. A holomorphic extension
through such a point would bound that fixed finite derivative on the
short radial segment, a contradiction. The selected points are dense on
the required arcs. For the even sum, the contour transformation
$(s,x,y)\mapsto(-s,-x,-y)$ gives evenness, and conjugation completes
all four quadrants. For the odd sum both signs were treated explicitly;
conjugation completes a dense subset of the circle. The inversion
\(s=1/\zeta\) is locally biholomorphic at each unit-circle point,
so the odd conclusion concerns the physical interior germ. This proves
the stated local obstructions. Their passage to all boundary points and
their role in the exact slice and isotropic corollaries are established
in Section~\ref{main:closure}.
\end{proof}

The high-temperature value at $s=0$ is the removable value one.
Indeed, uniformly in the real angular
variables, the decaying transfer root satisfies $\rho=s/2+O(s^2)$;
each one-body factor is $O(s)$ and each
squared pair is $O(s^2)$, giving $C_N(1/s)=O(s^{N^2})$ with a
bound $C^{N^2}|s|^{N^2}/N!$ for sufficiently small $|s|$.
The $N=1$ coefficient is $s/(1-s)^2$, and every higher odd order
vanishes after division by $s$. The exterior normalized
low-temperature germ is single valued; the energy-field susceptibility
is interpreted on the physical continuation of
$\beta_{\rm phys}=\operatorname{arsinh}(s)/(2J)$.
At every selected nonaxis point this factor is analytic and nonzero,
so the same dense obstructions hold for $\chi=\beta_{\rm phys}\mathcal X$.
No global single-valued inverse temperature on the exterior is asserted.


The exact coefficient identities and Lemma~\ref{geo:isotropic-cover}
fix the two full sums on the equal-coupling cells. The reduced LT
coefficients depend only on $s=\sinh z$ and are invariant under
$z\mapsto z+2\pi i k$; the physical field factor is $\beta=z/(2J)$
on that covering.

\section{Closure of the joint theorem and its exact descendants}
\label{main:closure}

\subsection{The complete joint obstruction}

\begin{proof}[Proof of Theorem~\ref{main:joint}]
Lemma~\ref{geo:endpoint-boundary} identifies the endpoint locus and
contains the entire boundary of the unsquared spectral set in it.
Theorem~\ref{geo:components} identifies the physical components
and their contractibility and \(p\)-bounds. Proposition~\ref{geo:pseudoconvex}
gives an explicit plurisubharmonic exhaustion, independent of the
susceptibility. Proposition~\ref{ana:normal}, with the physical
interface and its coefficientwise angular identification, proves
joint holomorphy and normal convergence of the actual physical sums.

Take a regular face point whose free coefficient is
\(b=\cos(2\pi a/\ell)\), with \(\ell\) an odd prime and
\(0<a<\ell/2\). By Proposition~\ref{geo:local-coordinates} the
variables \((C,B)\), after exchanging transfer axes if needed,
are independent holomorphic local coordinates, \(\Im B\ne0\),
and the physical domain is the full complementary side of that edge.
At such a point \(S,T\ne0\), and the edge identity
\eqref{geo:p-edge} implies \(p^2\ne1\). Thus the locally chosen
physical prefactor is analytic and nonzero.

At a positive endpoint, restrict to any holomorphic transversal with
nonramified edge coordinate \(u=C-B-b\). Its actual anisotropy
is a holomorphic function \(B(u)\). Conjugating the data when
needed gives \(\Im B(0)>0\), and the physical side contains
\(u=\ii\eps\). Proposition~\ref{ana:positive} proves a nonzero
singular derivative of the entire odd or even sum, including the
complete differentiated higher tail and physical prefactor. The
same proposition handles a negative endpoint for the even sum by
the exact angular shift and, if necessary, conjugation.

For a negative endpoint of the odd sum use
\[
             t=\frac{C+B-b}{B}.
\]
This is also nonramified on the stipulated transversal, because
\(B(0)\ne0\). In its lower-half-root orientation the actual
\(B(t)\) has \(\Im B(0)<0\), and the local physical side is
\(\Im(B(t)t)>0\). Theorem~\ref{neg:theorem} gives all finite
boundary jets of the complete odd sum and a selected sequence
whose coefficient roots diverge. Its exact full-sum asymptotic
\eqref{neg:full-asymptotic} excludes cancellation by every other
particle index and winding. Division by the nonvanishing physical
prefactor proves the same nonextension for the physical response.
The opposite orientation is its complex conjugate. The local uniform
complex-parameter hypotheses and constants are stated and proved in
Sections~\ref{ana:edges} and \ref{neg:section}; no estimate for
real fixed ratios is promoted to a joint estimate without proof.

In particular choose a constant-\(B\) transversal. If the joint
physical function extended to a neighborhood of the selected point,
its restriction to that disk would give the forbidden one-variable
extension. This is a valid implication from a proved transversal
obstruction to ambient nonextension. The selected points are dense
on regular faces by Proposition~\ref{geo:local-coordinates}, and
regular faces are dense in the entire physical boundaries by
Proposition~\ref{geo:joint-density}. Corollary~\ref{geo:propagate-obstruction}
then excludes componentwise extension at every boundary point,
including corners, collinear points and vanishing transfer-coordinate
charts. This proves the joint assertion and its transverse-curve
strengthening. Proposition~\ref{iso:local-obstruction}, proved
from the full sums in Section~\ref{iso:section}, supplies the
additional local assertion on the collinear equal-coupling restriction.

Finally consider analytic continuation in an unramified Riemann domain
over \(\C^2\) from one of these physical germs. As long as the
projected continuation path remains in its domain, uniqueness of
analytic continuation makes every continued germ the germ of the
single-valued function already constructed there. If a path first
leaves that domain, its first exit is a boundary point. A local
biholomorphic chart of the Riemann domain at the continued point
would yield a holomorphic function on a base neighborhood agreeing
with the physical function on an accumulating domain-side component.
The proved componentwise obstruction forbids this. Compactness of
the finite path interval guarantees a finite first exit whenever the
path reaches outside. Thus neither a monodromy construction nor a
different unramified physical sheet bypasses the boundary. This
uses the standard notion of holomorphic continuation, not a
definition tied to the spectral set.
\end{proof}

\subsection{All physical ratios, including the collinear one}

\begin{lemma}[Actual complex-ratio neighborhoods]\label{main:complexratio}
At any regular selected point of a positive-ratio slice where its edge
coordinate has nonzero derivative, the transverse representation and
the boundary estimates hold in a neighborhood with \(\lambda\)
complex. The uniformity is on common compact parameter polydisks,
not merely pointwise on the real-ratio family.
\end{lemma}
\begin{proof}
The boundary temperature coordinate \(z_0\) is nonzero. Thus
\((z,\lambda)\mapsto(z,w=\lambda z)\) has determinant \(z_0\)
and is a local biholomorphism. Its composition with the regular
\((C,B)\) coordinate map is holomorphic. If the positive-edge
parameter is \(u=C-B-b\), its derivative in \(z\) is nonzero
by hypothesis; the holomorphic implicit-function theorem gives
\(z=z(u,\lambda)\) and hence \(B=B(u,\lambda)\) on a product
polydisk. At a negative edge use \(t=(C+B-b)/B\) in the same way.
Shrinking the polydisk preserves the strict sign of \(\Im B\),
the nonvanishing physical prefactor and the other strict face
conditions. Cauchy estimates on a slightly larger polydisk give
the common holomorphic bounds required in
Proposition~\ref{ana:positive} and Theorem~\ref{neg:theorem}.
Their proofs therefore apply to the actual complex-ratio family.
In the negative case the selected integers may depend on the
family parameter, as explicitly stated there. No boundary estimate
is claimed uniform through the collinear degeneration at
\(\lambda=1\); that restriction is treated by its own proved
full-sum estimates. Joint holomorphy on interior compact sets already
includes complex-ratio neighborhoods at \(\lambda=1\).
\end{proof}

\begin{proof}[Proof of Theorem~\ref{main:slices}]
The exact open-set equalities are Theorem~\ref{geo:exact-pullbacks}.
The restrictions are jointly constructed physical functions by
Proposition~\ref{ana:normal}, and are holomorphic on precisely the
stated pullbacks. A boundary point of either pullback maps to the
corresponding joint boundary: it is a limit of domain points and
cannot map to the joint interior by the same open-set equality.

First fix \(\lambda\ne1\). Proposition~\ref{geo:slice-density}
proves, separately on this actual complex line, density of regular
boundary points at which the local edge map has nonzero derivative
and its free phase is one of the chosen prime-grid values. This
proof includes every finite positive ratio, whether rational or
irrational. In particular it does not rely on intersecting an
arbitrary ambient dense set with the line. At each such point the
temperature line is a holomorphic curve \(B=B(u)\) or \(B=B(t)\)
of exactly the kind treated in the proof of Theorem~\ref{main:joint}.
Apply Proposition~\ref{ana:positive} or Theorem~\ref{neg:theorem}
to this actual curve, retaining its varying anisotropy. These prove
nonextension of the restricted complete physical sum. The
one-variable collar in Lemma~\ref{geo:collar} and
Corollary~\ref{geo:propagate-obstruction} give every boundary point.

Now let \(\lambda=1\). The exact coefficient and normalization
identities in \eqref{iso:physical-L}--\eqref{iso:physical-H}
identify the restriction with the full even and odd integrals of
Section~\ref{iso:section}; its local obstructions were proved in
Proposition~\ref{iso:local-obstruction}. Lemma~\ref{geo:isotropic-cover} maps the two
slice domains to the exterior and disk, respectively. At every
boundary point with \(\cosh\zeta\ne0\) this is a local
biholomorphism. The selected obstructing points pull back densely
on the nonramified boundary arcs. The remaining points with
\(\cosh\zeta=0\) are limits of those arcs. The same local collar
therefore proves nonextension everywhere on both slice boundaries.
This argument uses no limit as \(\lambda\to1\).

All nonzero even orders were retained after subtraction of the
zero-particle \(M^2\), and the exact full-site multiplier retains
the origin and both axes. Section~\ref{ana:physical} proves the
physical energy-field equality \(\chi=\beta\mathcal X\) with
the plus phase fixed in LT. Here \(\beta=\zeta/(2J_x)\) is an
entire function of the slice coordinate and is nonzero on both
boundaries: zero lies in the HT interior and is strictly outside
the LT closure. Multiplication or division by it cannot change
local boundary extendability. The result consequently establishes the stated fixed-ratio
natural-boundary conclusions for both physical response conventions.
\end{proof}

\subsection{Derivation of the complete isotropic theorems from the slice theorem}

\begin{proof}[Proof of Corollary~\ref{main:isotropic}]
Apply the now proved Theorem~\ref{main:slices} with \(\lambda=1\).
The identities \(S=T=s\), \(p=s^2\), \(B=1\), and
\[
 \mathcal D(\zeta_1,\eta_1;z,z)
 =s\left[s+s^{-1}
 -\tfrac12(\zeta_1+\zeta_1^{-1}+\eta_1+\eta_1^{-1})\right]
\]
identify the parent coefficients term by term with the full isotropic
coefficients. On the real LT interval \(\sqrt p=s\), so each
factor \(\sqrt p/\mathcal D\) is exactly its isotropic counterpart.
The even prefactor becomes \((1-s^{-4})^{1/4}\). On the real HT
interval the dispersion is \(s+s^{-1}-\cos q\) and the odd
prefactor is \((1-s^4)^{1/4}/s\), giving the exterior auxiliary
variable \(1/s\). Normal convergence and the identity theorem
extend these identities to their physical complex domains.

For HT, the biholomorphism in Lemma~\ref{geo:isotropic-cover}
transports the restricted holomorphic function to the disk. The
removable value at zero is one, independently checked at the end
of Section~\ref{iso:section}. For LT, the normalized termwise
identities show invariance under every deck translation
\(z\mapsto z+2\pi\ii k\); the full normalized function therefore
descends from \(\Omega_{L,1}\) to a single-valued holomorphic
function on \(|s|>1\). It is the positive real LT germ and has
the stated connected onsite contribution.

Suppose one of these descended functions extended through a point
\(s_0\) of the unit circle. Choose a boundary preimage \(z_0\)
on its physical slice. Composing the proposed extension with the
entire map \(\sinh\) would extend the slice function near
\(z_0\), in agreement on the physical side. This contradicts
Theorem~\ref{main:slices}. This composition argument also works
at the ramified boundary values \(s_0=\pm\ii\); no local inverse
at those values is required. Hence the entire unit circle is a
natural boundary for both normalized germs.

Finally \(\beta=z/(2J)\) on the prescribed physical temperature
continuation yields the energy-field claims. In the disk it is
the branch \(\operatorname{arsinh}(s)/(2J)\). In the exterior it
is interpreted on the proved covering. At all selected nonaxis
points it is analytic and nonzero, so the local singular
derivative remains nonzero after the product rule; density gives
all remaining boundary points. This proves exactly the mathematical
HT and connected LT descendants, with distinct germs and the full
odd/even sums.
\end{proof}



\begin{thebibliography}{99}
\bibitem{BL} A. I. Bugrij and O. Lisovyy,
\emph{Correlation function of the two-dimensional Ising model on a
finite lattice. II}, Theoretical and Mathematical Physics \textbf{140}
(2004), 987--1000. \href{https://arxiv.org/abs/0708.3643v1}{arXiv:0708.3643v1}.
\bibitem{IL} N. Iorgov and O. Lisovyy,
\emph{Ising correlations and elliptic determinants},
Journal of Statistical Physics \textbf{143} (2011), 33--59.
\href{https://doi.org/10.1007/s10955-011-0154-6}{doi:10.1007/s10955-011-0154-6}.
\href{https://arxiv.org/abs/1012.2856v2}{arXiv:1012.2856v2}.
\bibitem{ONGP} W. P. Orrick, B. G. Nickel, A. J. Guttmann and J. H. H. Perk,
\emph{The susceptibility of the square lattice Ising model: New developments},
Journal of Statistical Physics \textbf{102} (2001), 795--841.
\href{https://doi.org/10.1023/A:1004850919647}{doi:10.1023/A:1004850919647}.
\href{https://arxiv.org/abs/cond-mat/0103074v1}{arXiv:cond-mat/0103074v1}.
\bibitem{BHM} S. Boukraa, S. Hassani and J.-M. Maillard,
\emph{Holonomic functions of several complex variables and singularities
of anisotropic Ising n-fold integrals}, Journal of Physics A:
Mathematical and Theoretical \textbf{45} (2012), 494010.
\href{https://doi.org/10.1088/1751-8113/45/49/494010}{doi:10.1088/1751-8113/45/49/494010}.
\href{https://arxiv.org/abs/1207.1784v2}{arXiv:1207.1784v2}.
\bibitem{Raoufi} A. Raoufi,
\emph{Translation-invariant Gibbs states of Ising model: General setting},
preprint (2017).
\href{https://arxiv.org/abs/1710.07608v1}{arXiv:1710.07608v1}.
\bibitem{CGN} F. Camia, C. Garban and C. M. Newman,
\emph{The Ising magnetization exponent on \(\mathbb Z^2\) is \(1/15\)},
preprint (2013).
\href{https://arxiv.org/abs/1205.6612v4}{arXiv:1205.6612v4}.

\end{thebibliography}

\end{document}
